% CorrectVCG: after \end{document} this file carries every correction rule and the checker (LaTeX ignores them).
% Never remove that block unless Doktor asks. Before any change run:
% sed -n '/^# BEGIN CORRECTVCG/,/^# END CORRECTVCG/p' THIS.tex > correctvcg_check.py && python3 correctvcg_check.py --rules && python3 correctvcg_check.py --paper THIS.tex
\documentclass[11pt,reqno]{amsart}
\usepackage{iftex}
\ifPDFTeX
  \usepackage[T1]{fontenc}
  \usepackage[utf8]{inputenc}
  \usepackage{lmodern}
\else
  \usepackage[no-math]{fontspec}
\fi
\usepackage{amsmath}
\usepackage{amssymb}
\usepackage{amsfonts}
\usepackage{mathrsfs}
\usepackage{amsthm}
\usepackage{mathtools}
\usepackage{geometry}
\usepackage{enumitem}
\usepackage{longtable}
\usepackage{booktabs}
\usepackage{array}
\usepackage[dvipsnames]{xcolor}
\usepackage[pagebackref=true]{hyperref}
\usepackage{bookmark}
\usepackage{pdfcomment}

% Hover previews for internal references (no layout change).
\makeatletter
\let\CG@origref\ref
\let\CG@origeqref\eqref
\newcommand{\CG@gettip}[1]{\@ifundefined{CGTip@#1}{Reference: #1}{\csname CGTip@#1\endcsname}}
\DeclareRobustCommand{\CGref}[1]{\pdftooltip{\CG@origref{#1}}{\CG@gettip{#1}}}
\DeclareRobustCommand{\CGeqref}[1]{\pdftooltip{\CG@origeqref{#1}}{\CG@gettip{#1}}}
% one line per label:
\expandafter\gdef\csname CGTip@sec:intro\endcsname{Section: Introduction}
\expandafter\gdef\csname CGTip@eq:Cdef\endcsname{Equation: C=-int\_1\textasciicircum{}infinitylog|sinpi u| u\textasciicircum{}(-2) du=0.74379486183940054...}
\expandafter\gdef\csname CGTip@sec:mellin\endcsname{Section: The Mellin transform of the log-sine function}
\expandafter\gdef\csname CGTip@eq:Cl2\endcsname{Equation: Cl2(theta)=-int\_0\textasciicircum{}thetalog|2sin2| d=sum\_(nu>=1)(sinnutheta)/(nu\textasciicircum{}2).}
\expandafter\gdef\csname CGTip@eq:Clmax\endcsname{Equation: |Cl2(theta)|2(pi/3)=1.01494160640965... (theta in R).}
\expandafter\gdef\csname CGTip@lem:euler\endcsname{Lemma - Euler integrals: Let alpha>0. For -1<Re w<0 the integral below converges absolutely and int\_0\textasciicircum{}infinity u\textasciicircum{}(w-1)sin(alpha u) du=alpha\textasciicircum{}(-w) Gamma(w)sin(pi w)/2.}
\expandafter\gdef\csname CGTip@eq:euler\endcsname{Equation: int\_0\textasciicircum{}infinity u\textasciicircum{}(w-1)sin(alpha u) du=alpha\textasciicircum{}(-w) Gamma(w)sin(pi w)/2.}
\expandafter\gdef\csname CGTip@eq:eulerz\endcsname{Equation: \_(R to infinity)int\_0\textasciicircum{}Ru\textasciicircum{}(z-1)e\textasciicircum{}(+- iu) du=Gamma(z) e\textasciicircum{}(+- ipi z/2).}
\expandafter\gdef\csname CGTip@eq:arc\endcsname{Equation: e\textasciicircum{}(pi|Im z|/2)R\textasciicircum{}(Re z)int\_0\textasciicircum{}(pi/2)e\textasciicircum{}(-2Rtheta/pi) dtheta<= e\textasciicircum{}(pi|Im z|/2) 2 R\textasciicircum{}(Re z-1),}
\expandafter\gdef\csname CGTip@eq:parts\endcsname{Equation: int\_0\textasciicircum{}infinity u\textasciicircum{}(w-1)sin u du=[(u\textasciicircum{}wsin u)/(w)]\_0\textasciicircum{}infinity-1/wint\_0\textasciicircum{}infinity u\textasciicircum{}(w)cos u du=-1/w Gamma(w+1)cos(pi(w+1))/2,}
\expandafter\gdef\csname CGTip@eq:Adef\endcsname{Equation: A(s)=(log(2pi))/s+1/(s\textasciicircum{}2)+sum\_(k>=1)(zeta(2k))/(k (2k-s)).}
\expandafter\gdef\csname CGTip@thm:mellin\endcsname{Theorem - Mellin transform of F: For Re s>0 the integral Fhat(s)=int\_1\textasciicircum{}infinity F(u) u\textasciicircum{}(-s-1) du}
\expandafter\gdef\csname CGTip@eq:Fhatdef\endcsname{Equation: Fhat(s)=int\_1\textasciicircum{}infinity F(u) u\textasciicircum{}(-s-1) du}
\expandafter\gdef\csname CGTip@eq:Fhat\endcsname{Equation: Fhat(s)= scot((pi s)/2)zeta(s)+A(s),}
\expandafter\gdef\csname CGTip@eq:fourier\endcsname{Equation: F(v)=-sum\_(nu>=1)(cos2pinu v)/(nu) in L\textasciicircum{}2(0,1).}
\expandafter\gdef\csname CGTip@eq:Phi\endcsname{Equation: Phi(u)=-1/(2pi)sum\_(nu>=1)(sin2pinu u)/(nu\textasciicircum{}2)=-(Cl2(2pi u))/(2pi) (u>=0),}
\expandafter\gdef\csname CGTip@eq:Phi0\endcsname{Equation: |Phi(u)|<=int\_0\textasciicircum{}ulog1/(4v) dv=u(log1/(4u)+1) (0<u<=1/6).}
\expandafter\gdef\csname CGTip@eq:Gparts\endcsname{Equation: Fhat(s)=(s+1)int\_1\textasciicircum{}infinityPhi(u) u\textasciicircum{}(-s-2) du,}
\expandafter\gdef\csname CGTip@eq:L\endcsname{Equation: L(s):=int\_0\textasciicircum{}infinity F(u) u\textasciicircum{}(-s-1) du=(s+1)int\_0\textasciicircum{}infinityPhi(u) u\textasciicircum{}(-s-2) du.}
\expandafter\gdef\csname CGTip@eq:tonelli\endcsname{Equation: sum\_(nu>=1)1/(nu\textasciicircum{}2)int\_0\textasciicircum{}infinity|sin2pinu u| u\textasciicircum{}(-sigma-2) du=b\_sigma(2pi)\textasciicircum{}(sigma+1)sum\_(nu>=1)nu\textasciicircum{}(sigma-1)<infinity,}
\expandafter\gdef\csname CGTip@eq:Lzeta\endcsname{Equation: L(s)=-(2pi)\textasciicircum{}s Gamma(-s)cos((pi s)/2)zeta(1-s)= scot((pi s)/2)zeta(s) (-1<sigma<0).}
\expandafter\gdef\csname CGTip@eq:I\endcsname{Equation: I(s):=int\_0\textasciicircum{}1F(u) u\textasciicircum{}(-s-1) du=-(log(2pi))/s-1/(s\textasciicircum{}2)+ J(s),}
\expandafter\gdef\csname CGTip@eq:J\endcsname{Equation: J(s):=int\_0\textasciicircum{}1log((sinpi u)/(pi u))u\textasciicircum{}(-s-1) du=-sum\_(k>=1)(zeta(2k))/(k (2k-s)) (Re s<2).}
\expandafter\gdef\csname CGTip@eq:sinprod\endcsname{Equation: log(sinpi u)/(pi u)=sum\_(nu>=1)log(1-(u\textasciicircum{}2)/(nu\textasciicircum{}2))=-sum\_(k>=1)(zeta(2k))/(k) u\textasciicircum{}(2k).}
\expandafter\gdef\csname CGTip@eq:Gid\endcsname{Equation: Fhat(s)= scot((pi s)/2)zeta(s)+(log(2pi))/s+1/(s\textasciicircum{}2)- J(s).}
\expandafter\gdef\csname CGTip@cor:B\endcsname{Corollary - Two forms of sA(s), and the constants: Let B(s)=sA(s). For Re s<2, s != 0, B(s)=log(2pi)+1/s+sum\_(k>=1)zeta(2k)(2/(2k-s)-1/k),}
\expandafter\gdef\csname CGTip@eq:B1\endcsname{Equation: B(s)=log(2pi)+1/s+sum\_(k>=1)zeta(2k)(2/(2k-s)-1/k),}
\expandafter\gdef\csname CGTip@eq:B2\endcsname{Equation: B(s)=logpi+1/s+sum\_(k>=1)(s/k+2(zeta(2k)-1))/(2k-s),}
\expandafter\gdef\csname CGTip@eq:aconst\endcsname{Equation: a:=pi(C-log2)=-piFhat(1)=0.15911438342889...>0.}
\expandafter\gdef\csname CGTip@eq:log2\endcsname{Equation: sum\_(k>=1)(zeta(2k)-1)/k=log2.}
\expandafter\gdef\csname CGTip@sec:Anonzero\endcsname{Section: The function A has no zeros in the critical strip}
\expandafter\gdef\csname CGTip@thm:Anonzero\endcsname{Theorem: A(s) != 0 for 0<Re s<1. More precisely, with s=sigma+it, 0<sigma<1, Re B(s)>logpi-(zeta(3))/2=0.54370... (|t|<=sqrt2),}
\expandafter\gdef\csname CGTip@eq:caseI\endcsname{Equation: Re B(s)>logpi-(zeta(3))/2=0.54370... (|t|<=sqrt2),}
\expandafter\gdef\csname CGTip@eq:caseII\endcsname{Equation: (Im B(s))/t=-1/(sigma\textasciicircum{}2+t\textasciicircum{}2)+sum\_(k>=1)(2zeta(2k))/((2k-sigma)\textasciicircum{}2+t\textasciicircum{}2)>0 (|t|>=sqrt2).}
\expandafter\gdef\csname CGTip@eq:caseIid\endcsname{Equation: Re (s)/(2k-s)+(t\textasciicircum{}2)/(4k\textasciicircum{}2+t\textasciicircum{}2)=(2ksigma(2k(2k-sigma)-t\textasciicircum{}2))/(((2k-sigma)\textasciicircum{}2+t\textasciicircum{}2)(4k\textasciicircum{}2+t\textasciicircum{}2)).}
\expandafter\gdef\csname CGTip@eq:caseIbound\endcsname{Equation: Re (s)/(k(2k-s))>-(t\textasciicircum{}2)/(k(4k\textasciicircum{}2+t\textasciicircum{}2))>=-(t\textasciicircum{}2)/(4k\textasciicircum{}3)>=-1/(2k\textasciicircum{}3).}
\expandafter\gdef\csname CGTip@eq:caseIIbound\endcsname{Equation: (pi\textasciicircum{}2)/3(sigma\textasciicircum{}2+t\textasciicircum{}2)-(2-sigma)\textasciicircum{}2-t\textasciicircum{}2>=((pi\textasciicircum{}2)/3-1)t\textasciicircum{}2-(2-sigma)\textasciicircum{}2>=2((pi\textasciicircum{}2)/3-1)-4=(2pi\textasciicircum{}2)/3-6=0.5797...>0}
\expandafter\gdef\csname CGTip@rem:Anum\endcsname{Remark: The identities eq:caseIid and eq:caseII, and the expansion behind eq:caseIIbound, are also checked symbolically by route\_constants.py. At the first zero rho\_1=1/2+i 14.1347... of zeta one finds numerically A(rho\_1)=0.107}
\expandafter\gdef\csname CGTip@sec:H\endcsname{Section: The Mobius log-sine function H}
\expandafter\gdef\csname CGTip@eq:Hdef\endcsname{Equation: H(x)=sum\_(d<= x)mu(d) F( xd).}
\expandafter\gdef\csname CGTip@prop:Hlocal\endcsname{Proposition - Local structure: (i) For X>=1, int\_1\textasciicircum{}X|H(x)|\textasciicircum{}2 dx<=(pi\textasciicircum{}2)/(12)X\textasciicircum{}3, int\_1\textasciicircum{}X|H(x)| dx<=(pi)/(2sqrt3)X\textasciicircum{}2.}
\expandafter\gdef\csname CGTip@eq:Hcrude\endcsname{Equation: int\_1\textasciicircum{}X|H(x)|\textasciicircum{}2 dx<=(pi\textasciicircum{}2)/(12)X\textasciicircum{}3, int\_1\textasciicircum{}X|H(x)| dx<=(pi)/(2sqrt3)X\textasciicircum{}2.}
\expandafter\gdef\csname CGTip@eq:Hplus\endcsname{Equation: H(n\textasciicircum{}+)=Lambda(n)+sum\_(d<= n, d not dividing n)mu(d) F( nd).}
\expandafter\gdef\csname CGTip@prop:Hid\endcsname{Proposition - Two identities: Let x>1, x, and let log(1-e\textasciicircum{}(2pi iv)) denote the principal branch. Then sum\_(d<= x)mu(d)log(1-e\textasciicircum{}(2pi ix/d))=H(x)+ipi(x m(x)-1-(M(x))/2).}
\expandafter\gdef\csname CGTip@eq:Hcomplex\endcsname{Equation: sum\_(d<= x)mu(d)log(1-e\textasciicircum{}(2pi ix/d))=H(x)+ipi(x m(x)-1-(M(x))/2).}
\expandafter\gdef\csname CGTip@eq:Hsinc\endcsname{Equation: H(x)=log|x\textasciicircum{}2-1|+log(2pi)M(x)+int\_1\textasciicircum{}x(M(u))/u du-sum\_(d>x)mu(d)logsinc xd,}
\expandafter\gdef\csname CGTip@eq:sincid\endcsname{Equation: sum\_(d>=1)mu(d)log|sinc xd|=log|1-x\textasciicircum{}2|.}
\expandafter\gdef\csname CGTip@prop:HMellin\endcsname{Proposition - Mellin transform of H: For Re s>1 the integral Hhat(s)=int\_1\textasciicircum{}infinity H(x)x\textasciicircum{}(-s-1) dx converges absolutely and Hhat(s)=(Fhat(s))/(zeta(s)), Hhat(s)= scot((pi s)/2)+(A(s))/(zeta(s)) (1<Re s<2).}
\expandafter\gdef\csname CGTip@eq:HMellin\endcsname{Equation: Hhat(s)=(Fhat(s))/(zeta(s)), Hhat(s)= scot((pi s)/2)+(A(s))/(zeta(s)) (1<Re s<2).}
\expandafter\gdef\csname CGTip@thm:L2\endcsname{Theorem - L\textasciicircum{}2 criterion: The following are equivalent. (a) RH. (b) int\_1\textasciicircum{}infinity|H(x)|\textasciicircum{}2x\textasciicircum{}(-2sigma-1) dx<infinity for every sigma>1/2. (c) For every eps>0 there is K\_eps with int\_1\textasciicircum{}XH(x)\textasciicircum{}2 dx<= K\_eps X\textasciicircum{}(2+eps) for all X>=1.}
\expandafter\gdef\csname CGTip@eq:dyadic\endcsname{Equation: int\_(2\textasciicircum{}q)\textasciicircum{}2\textasciicircum{}(q+1)|H|x\textasciicircum{}(-sigma-1) dx<=2\textasciicircum{}(-q(sigma+1))2\textasciicircum{}(q/2)(K\_eps2\textasciicircum{}((q+1)(2+eps)))\textasciicircum{}(1/2) << \_sigma2\textasciicircum{}(-q(sigma-1/2-eps/2))=2\textasciicircum{}(-qeps/2),}
\expandafter\gdef\csname CGTip@eq:Hhatbound\endcsname{Equation: |Hhat(sigma+it)| << \_eps|t|\textasciicircum{}(-1+eps) (sigma\_0<=sigma<=3/2, |t|>=2),}
\expandafter\gdef\csname CGTip@eq:H2sup\endcsname{Equation: \_(sigma>sigma\_0)1/(2pi)int\_(-infinity)\textasciicircum{}infinity|Hhat(sigma+it)|\textasciicircum{}2 dt<infinity.}
\expandafter\gdef\csname CGTip@rem:L2\endcsname{Remark: A pointwise bound H(x) << x\textasciicircum{}(1/2+eps) fails: by Proposition prop:Hlocal(ii), H(x)+mu(n)log(n-x) has a finite limit as x to n\textasciicircum{}-, so H is unbounded near every squarefree n>=2. This is why the pointwise criterion is stated }
\expandafter\gdef\csname CGTip@sec:W\endcsname{Section: The Clausen sum W}
\expandafter\gdef\csname CGTip@eq:Wdef\endcsname{Equation: W(X)=sum\_(d<= X)mu(d) d Cl2((2pi X)/d).}
\expandafter\gdef\csname CGTip@prop:Wbasic\endcsname{Proposition: (i) W(X)=-2piint\_1\textasciicircum{}XH(x) dx for X>=1. (ii) With Kern(v)=Cl2(2pi v)+2pi v F(v) for v, W(X)=-int\_1\textasciicircum{}XM(y) Kern( Xy) dy.}
\expandafter\gdef\csname CGTip@eq:WMK\endcsname{Equation: W(X)=-int\_1\textasciicircum{}XM(y) Kern( Xy) dy.}
\expandafter\gdef\csname CGTip@eq:WMellin\endcsname{Equation: What(s)=-(2pi)/s Hhat(s-1)=-(2pi)/s (Fhat(s-1))/(zeta(s-1)).}
\expandafter\gdef\csname CGTip@lem:landau\endcsname{Lemma - Landau: Let f be real, locally bounded and measurable on [1,infinity), with f(X)>=0 for X>= X\_0. Suppose g(s)=int\_1\textasciicircum{}infinity f(X)X\textasciicircum{}(-s-1)dX converges absolutely for Re s>sigma\_1, and that g extends holomorphically to a neighbour}
\expandafter\gdef\csname CGTip@eq:landau\endcsname{Equation: g(sigma\_3-h)=sum\_(k>=0)(h\textasciicircum{}k)/(k!)int\_1\textasciicircum{}infinity f(X)(log X)\textasciicircum{}kX\textasciicircum{}(-sigma\_3-1) dX=int\_1\textasciicircum{}infinity f(X)X\textasciicircum{}(-(sigma\_3-h)-1) dX,}
\expandafter\gdef\csname CGTip@thm:Wcrit\endcsname{Theorem - Pointwise and one-sided criteria: The following are equivalent. (a) RH. (b) For every eps>0 there is K\_eps with |W(X)|<= K\_eps X\textasciicircum{}(3/2+eps) for X>=1. (c) For every eps>0 there is K\_eps with W(X)<= K\_eps X\textasciicircum{}(3/2+eps) for X>=1. (d) For every eps>0 there is K}
\expandafter\gdef\csname CGTip@eq:WMKbound\endcsname{Equation: |W(X)|<=int\_1\textasciicircum{}X|M( Xv)||Kern(v)|X/(v\textasciicircum{}2) dv << \_eps X\textasciicircum{}(3/2+eps)int\_1\textasciicircum{}infinity v\textasciicircum{}(-5/2)|Kern(v)| dv.}
\expandafter\gdef\csname CGTip@thm:omega\endcsname{Theorem - \_+- bounds: Let rho=1/2+igamma, gamma>0, be a simple zero of zeta, and r\_rho=-(2pi)/(1+rho) (A(rho))/(zeta'(rho)) != 0.}
\expandafter\gdef\csname CGTip@eq:rrho\endcsname{Equation: r\_rho=-(2pi)/(1+rho) (A(rho))/(zeta'(rho)) != 0.}
\expandafter\gdef\csname CGTip@eq:fejer\endcsname{Equation: g(sigma)+1/2e\textasciicircum{}(iphi\_0)g(sigma+igamma)+1/2e\textasciicircum{}(-iphi\_0)g(sigma-igamma)=int\_1\textasciicircum{}infinity f(X)X\textasciicircum{}(-sigma-1)(1+cos(gammalog X-phi\_0)) dX>=-Q\_0,}
\expandafter\gdef\csname CGTip@sec:corner\endcsname{Section: The d=2 transfer and the corner}
\expandafter\gdef\csname CGTip@lem:vaughan\endcsname{Lemma: Let U>=1 and mu\_U(n)=mu(n) for n<= U, mu\_U(n)=0 for n>U. For 1<= n<= U\textasciicircum{}2, mu(n)=2mu\_U(n)-(mu\_U*mu\_U*1)(n).}
\expandafter\gdef\csname CGTip@eq:vaughan\endcsname{Equation: mu(n)=2mu\_U(n)-(mu\_U*mu\_U*1)(n).}
\expandafter\gdef\csname CGTip@eq:Gdef\endcsname{Equation: G(Y)=sum\_(l<= Y)l Cl2(2pi Yl), E\_G(Y)=G(Y)-aY\textasciicircum{}2,}
\expandafter\gdef\csname CGTip@eq:Thetadef\endcsname{Equation: Theta(s)=int\_1\textasciicircum{}infinityCl2(2pi v) v\textasciicircum{}(-s-1) dv=-(2pi)/s Fhat(s-1).}
\expandafter\gdef\csname CGTip@lem:EG\endcsname{Lemma: (i) For Re s>2, Ghat(s)=int\_1\textasciicircum{}infinity G(Y)Y\textasciicircum{}(-s-1)dY=zeta(s-1)Theta(s). (ii) For every delta in (0,1) there is c\_delta with |E\_G(Y)|<= c\_delta Y\textasciicircum{}(1+delta) for Y>=1. (iii) For Y, E\_G'(Y)=-2pi R(Y), where R(Y)=sum\_(l<= Y)}
\expandafter\gdef\csname CGTip@eq:Ghatform\endcsname{Equation: Ghat(s)=-(2pi)/s zeta(w)( wcot(pi w)/2 zeta(w)+A(w)).}
\expandafter\gdef\csname CGTip@eq:Ghatbound\endcsname{Equation: |Ghat(s)| << |t|\textasciicircum{}(-2)(|zeta(w)|\textasciicircum{}2+|zeta(w)|log|t|) (1<sigma<=5/2, |t|>=2).}
\expandafter\gdef\csname CGTip@eq:Ginv\endcsname{Equation: G(Y)=1/(2pi i)int\_((5/2))Ghat(s) Y\textasciicircum{}s ds (Y>=1).}
\expandafter\gdef\csname CGTip@eq:EGint\endcsname{Equation: E\_G(Y)=1/(2pi i)int\_((1+delta))Ghat(s) Y\textasciicircum{}s ds, |E\_G(Y)|<=Y\textasciicircum{}(1+delta)2piint\_(-infinity)\textasciicircum{}infinity|Ghat(1+delta+it)| dt=c\_delta Y\textasciicircum{}(1+delta).}
\expandafter\gdef\csname CGTip@eq:Wflat\endcsname{Equation: W\textasciicircum{}b(X)=sum\_(n<=sqrt X)mu(n) n Cl2((2pi X)/n),}
\expandafter\gdef\csname CGTip@eq:Rest\endcsname{Equation: Rest(X)=sum\_d,j<=sqrt X dj<=sqrt Xmu(d)mu(j) dj E\_G(X/(dj)), Corner(X)=sum\_d,j<=sqrt X dj>sqrt Xmu(d)mu(j) dj E\_G(X/(dj)).}
\expandafter\gdef\csname CGTip@prop:decomp\endcsname{Proposition - Exact decomposition: For X>=1, W(X)=2W\textasciicircum{}b(X)-aX\textasciicircum{}2m(sqrt X)\textasciicircum{}2-Rest(X)-Corner(X).}
\expandafter\gdef\csname CGTip@eq:decomp\endcsname{Equation: W(X)=2W\textasciicircum{}b(X)-aX\textasciicircum{}2m(sqrt X)\textasciicircum{}2-Rest(X)-Corner(X).}
\expandafter\gdef\csname CGTip@eq:Restbound\endcsname{Equation: |Rest(X)|<= c\_delta X\textasciicircum{}(3/2+delta/2)(1+logsqrt X)\textasciicircum{}2.}
\expandafter\gdef\csname CGTip@eq:decomp1\endcsname{Equation: W(X) =2W\textasciicircum{}b(X)-sum\_(d,j<=sqrt X)mu(d)mu(j)sum\_(l<= X/(dj))djl Cl2((2pi X)/(djl)) =2W\textasciicircum{}b(X)-sum\_(d,j<=sqrt X)mu(d)mu(j) dj G(X/(dj)),}
\expandafter\gdef\csname CGTip@eq:Restpf\endcsname{Equation: |Rest(X)|<= c\_deltasum\_(n<=sqrt X)tau(n) n( Xn)\textasciicircum{}(1+delta)<= c\_delta X\textasciicircum{}(1+delta)(sqrt X)\textasciicircum{}(1-delta)sum\_(n<=sqrt X)(tau(n))/n=c\_delta X\textasciicircum{}(3/2+delta/2)sum\_(n<=sqrt X)(tau(n))/n,}
\expandafter\gdef\csname CGTip@thm:corner\endcsname{Theorem - Reduction of RH to the corner: The following are equivalent. (a) RH. (b) For every eps>0 there is K\_eps with Corner(X)>=-K\_eps X\textasciicircum{}(3/2+eps) for all X>=1. (c) For every eps>0 there is K\_eps with |Corner(X)|<= K\_eps X\textasciicircum{}(3/2+eps) for all X>=1.}
\expandafter\gdef\csname CGTip@eq:cornerpf\endcsname{Equation: W(X) =2W\textasciicircum{}b(X)-aX\textasciicircum{}2m(sqrt X)\textasciicircum{}2-Rest(X)-Corner(X) 2(3)(X+sqrt X)+c\_delta X\textasciicircum{}(3/2+delta/2)(1+logsqrt X)\textasciicircum{}2+K\_eps X\textasciicircum{}(3/2+eps).}
\expandafter\gdef\csname CGTip@eq:mbound\endcsname{Equation: m(y)=(M(y))/y+int\_1\textasciicircum{}y(M(v))/(v\textasciicircum{}2) dv=(M(y))/y-int\_y\textasciicircum{}infinity(M(v))/(v\textasciicircum{}2) dv << \_eps y\textasciicircum{}(-1/2+eps) (y>=1),}
\expandafter\gdef\csname CGTip@rem:corner\endcsname{Remark: Corner(X) depends on mu only through its values on [1,sqrt X], and on E\_G only on [1,sqrt X). The function E\_G involves no Mobius values. By Lemma lem:EG(iii), the derivative of E\_G is -2pi R, and at each integer n>=2 th}
\expandafter\gdef\csname CGTip@sec:other\endcsname{Section: Other exact identities}
\expandafter\gdef\csname CGTip@prop:buchstab\endcsname{Proposition - Recursion over large primes: For X>=1 let W\textasciicircum{}sm(X) be the sum eq:Wdef restricted to n whose prime factors are all <=sqrt X. Then W(X)=W\textasciicircum{}sm(X)-sum\_(sqrt X<p<= X)p W( Xp).}
\expandafter\gdef\csname CGTip@eq:buchstab\endcsname{Equation: W(X)=W\textasciicircum{}sm(X)-sum\_(sqrt X<p<= X)p W( Xp).}
\expandafter\gdef\csname CGTip@prop:lcm\endcsname{Proposition - lcm form: Let X>=1, let D be the set of squarefree d<= X, let Xi\_X(n)=n Cl2(2pi X/n) for n<= X and Xi\_X(n)=0 for n>X, and let K be the matrix (Xi\_X([d,d']))\_(d,d' in D), where [d,d'] is the least common multiple. Let =(mu(d'))\_(d'}
\expandafter\gdef\csname CGTip@prop:d3\endcsname{Proposition - d=3: Let X>=1, U=X\textasciicircum{}(1/3), G\_2(Y)=sum\_(l\_1l\_2<= Y)l\_1l\_2 Cl2(2pi Y/(l\_1l\_2)), and T\_1=sum\_(n<= U)mu(n) n Cl2((2pi X)/n), T\_2=sum\_(d,j<= U)mu(d)mu(j) dj G(X/(dj)), T\_3=sum\_(d,j,h<= U)mu(d)mu(j)mu(h) djh G\_2(X/(djh)).}
\expandafter\gdef\csname CGTip@eq:d3\endcsname{Equation: T\_1=sum\_(n<= U)mu(n) n Cl2((2pi X)/n), T\_2=sum\_(d,j<= U)mu(d)mu(j) dj G(X/(dj)), T\_3=sum\_(d,j,h<= U)mu(d)mu(j)mu(h) djh G\_2(X/(djh)).}
\expandafter\gdef\csname CGTip@sec:methods\endcsname{Section: Two facts about methods}
\expandafter\gdef\csname CGTip@prop:musq\endcsname{Proposition: Let W\_(mu\textasciicircum{}2)(X)=sum\_(n<= X)mu(n)\textasciicircum{}2 n Cl2(2pi X/n). For every delta in (0,1/2), W\_(mu\textasciicircum{}2)(X)=6/(pi\textasciicircum{}2) a X\textasciicircum{}2+O\_delta(X\textasciicircum{}(3/2+delta)), 6/(pi\textasciicircum{}2)a=0.0967299459812...}
\expandafter\gdef\csname CGTip@eq:musq\endcsname{Equation: W\_(mu\textasciicircum{}2)(X)=6/(pi\textasciicircum{}2) a X\textasciicircum{}2+O\_delta(X\textasciicircum{}(3/2+delta)), 6/(pi\textasciicircum{}2)a=0.0967299459812...}
\expandafter\gdef\csname CGTip@rem:barrier\endcsname{Remark - What a bound for Corner must use: For squarefree d,d' one has mu(d)mu(d')=mu(dd'/(d,d')\textasciicircum{}2), and the same relation holds with mu replaced by mu\textasciicircum{}2, or by n->mu(n)\textasciicircum{}2prod\_(p| n)chi(p) for any chi from the primes to (+-1). An argument that uses mu only throug}
\expandafter\gdef\csname CGTip@rem:local\endcsname{Remark - Localization, not proved: An inequality of the form W(X)<= c\_1X+c\_eps X\textasciicircum{}(3/2+kappa+eps)(1+J(sqrt X)), with J(y)=int\_1\textasciicircum{}yM(v)\textasciicircum{}2v\textasciicircum{}(-2)dv+M(y)\textasciicircum{}2/y and kappa a subconvexity exponent for zeta on the critical line, was derived in exploratory work throug}
\expandafter\gdef\csname CGTip@sec:muntz\endcsname{Section: A Hilbert-space form: the M\"untz transform of the log-sine kernel}
\expandafter\gdef\csname CGTip@eq:kdef\endcsname{Equation: k(t)=F(1/t)=log|2sin t| (0<t<=1, 1/t), k(t)=0 (t>1).}
\expandafter\gdef\csname CGTip@eq:Nsigma\endcsname{Equation: N\_sigma=int\_1\textasciicircum{}infinity F(v)\textasciicircum{}2v\textasciicircum{}(-2sigma-1) dv<=(pi\textasciicircum{}2)/(12)sum\_(n>=1)n\textasciicircum{}(-2sigma-1)<infinity,}
\expandafter\gdef\csname CGTip@eq:kmellin\endcsname{Equation: int\_0\textasciicircum{}infinity k(t) t\textasciicircum{}(s-1) dt=int\_1\textasciicircum{}infinity F(v) v\textasciicircum{}(-s-1) dv=Fhat(s) (Re s>0), int\_0\textasciicircum{}infinity k=Fhat(1)=log2-C.}
\expandafter\gdef\csname CGTip@prop:muntz\endcsname{Proposition - Dilation sums of k: Let t>0 with 1/t and put x=1/t. Then sum\_(n>=1) k(nt)-1/tint\_0\textasciicircum{}infinity k=R(x), sum\_(n>=1)mu(n) k(nt)=H(x) (x>=1),}
\expandafter\gdef\csname CGTip@eq:muntz\endcsname{Equation: sum\_(n>=1) k(nt)-1/tint\_0\textasciicircum{}infinity k=R(x), sum\_(n>=1)mu(n) k(nt)=H(x) (x>=1),}
\expandafter\gdef\csname CGTip@eq:muntzmellin\endcsname{Equation: int\_0\textasciicircum{}infinity(sum\_(n>=1) k(nt))t\textasciicircum{}(s-1) dt=zeta(s)Fhat(s), int\_0\textasciicircum{}infinity(sum\_(n>=1)mu(n) k(nt))t\textasciicircum{}(s-1) dt=(Fhat(s))/(zeta(s)).}
\expandafter\gdef\csname CGTip@eq:energy\endcsname{Equation: E\_N(sigma) =int\_0\textasciicircum{}infinity|sum\_(d<= N)mu(d) k(dt)|\textasciicircum{}2t\textasciicircum{}(2sigma-1) dt, \_sigma(d,e) =int\_0\textasciicircum{}infinity k(dt) k(et) t\textasciicircum{}(2sigma-1) dt=int\_((d,e))\textasciicircum{}infinity F( xd)F( xe)dxx\textasciicircum{}(2sigma+1),}
\expandafter\gdef\csname CGTip@prop:energy\endcsname{Proposition - Spectral form of the energy: Let sigma>0 and N>=1. Then E\_N(sigma)<infinity, and E\_N(sigma)=1/(2pi)int\_(-infinity)\textasciicircum{}infinity|Fhat(sigma+itau)|\textasciicircum{}2|sum\_(d<= N)mu(d)d\textasciicircum{}(sigma+itau)|\textasciicircum{}2dtau=sum\_(d,e<= N)mu(d)mu(e) \_sigma(d,e),}
\expandafter\gdef\csname CGTip@eq:energyspec\endcsname{Equation: E\_N(sigma)=1/(2pi)int\_(-infinity)\textasciicircum{}infinity|Fhat(sigma+itau)|\textasciicircum{}2|sum\_(d<= N)mu(d)d\textasciicircum{}(sigma+itau)|\textasciicircum{}2dtau=sum\_(d,e<= N)mu(d)mu(e) \_sigma(d,e),}
\expandafter\gdef\csname CGTip@eq:gram\endcsname{Equation: \_sigma(d,e)=(\_sigma(log(e/d)))/((de)\textasciicircum{}sigma), \_sigma(u)=1/(2pi)int\_(-infinity)\textasciicircum{}infinity|Fhat(sigma+itau)|\textasciicircum{}2e\textasciicircum{}(itau u) dtau, \_sigma(d,d)= N\_sigmad\textasciicircum{}(2sigma).}
\expandafter\gdef\csname CGTip@eq:diag\endcsname{Equation: sum\_(d<= N)mu(d)\textasciicircum{}2 \_sigma(d,d)<= N\_sigma zeta(2sigma) (sigma>1/2),}
\expandafter\gdef\csname CGTip@eq:gramproof\endcsname{Equation: \_sigma(d,e)=int\_( R)g\_dg\_e dv=1/(2pi)int\_(-infinity)\textasciicircum{}infinity d\textasciicircum{}(-sigma-itau) e\textasciicircum{}(-sigma-itau) |Fhat(sigma+itau)|\textasciicircum{}2 dtau=1/((de)\textasciicircum{}sigma) \_sigma(log ed),}
\expandafter\gdef\csname CGTip@thm:energy\endcsname{Theorem - A sufficient condition for RH: If for every sigma>1/2 one has \_(N to infinity) E\_N(sigma)<infinity, then RH holds. In particular \_(N) E\_N(sigma)<infinity for every sigma>1/2 implies RH.}
\expandafter\gdef\csname CGTip@eq:energylower\endcsname{Equation: E\_N(sigma)>=int\_(1/N)\textasciicircum{}1|sum\_(d<= N)mu(d) k(dt)|\textasciicircum{}2t\textasciicircum{}(2sigma-1) dt=int\_1\textasciicircum{}N|H(x)|\textasciicircum{}2x\textasciicircum{}(-2sigma-1) dx.}
\expandafter\gdef\csname CGTip@eq:target2\endcsname{Equation: \_(N>=1)sum\_d,e<= N d != e(mu(d)mu(e))/((de)\textasciicircum{}sigma) \_sigma(log ed)<infinity for every sigma>1/2}
\expandafter\gdef\csname CGTip@sec:nb\endcsname{Section: A weighted strong Nyman--Beurling criterion for the log-sine kernel}
\expandafter\gdef\csname CGTip@eq:Pk\endcsname{Equation: (P k)(t)=sum\_(n>=1) k(nt)-1/tint\_0\textasciicircum{}infinity k=R(1/t) (0<t<=1), (P k)(t)=(C-log2)/t (t>1),}
\expandafter\gdef\csname CGTip@lem:convex\endcsname{Lemma - Uniform convexity bound: Let 1/2<sigma\_0<1 and eps>0. Then |zeta(sigma'+itau)| << \_(sigma\_0,eps)|tau|\textasciicircum{}((1-sigma\_0)/2+eps) uniformly for sigma\_0<=sigma'<=2, |tau|>=2.}
\expandafter\gdef\csname CGTip@lem:Rweight\endcsname{Lemma - R in the weighted spaces: (i) int\_1\textasciicircum{}X R(x)\textasciicircum{}2dx<=(pi\textasciicircum{}2)/(12)X\textasciicircum{}3+2(log2-C)\textasciicircum{}2X\textasciicircum{}3 for X>=1, so that R(s)=int\_1\textasciicircum{}infinity R(x)x\textasciicircum{}(-s-1)dx converges absolutely for Re s>1, and there R(s)=zeta(s)Fhat(s)-(Fhat(1))/(s-1).}
\expandafter\gdef\csname CGTip@eq:Rhat\endcsname{Equation: R(s)=zeta(s)Fhat(s)-(Fhat(1))/(s-1).}
\expandafter\gdef\csname CGTip@eq:Rhatbound\endcsname{Equation: | R(sigma'+itau)| << |tau|\textasciicircum{}(-sigma\_0+2eps)+|tau|\textasciicircum{}(-(1+sigma\_0)/2+eps)log|tau|+|tau|\textasciicircum{}(-1) << |tau|\textasciicircum{}(-sigma\_0+2eps),}
\expandafter\gdef\csname CGTip@eq:varrho\endcsname{Equation: \_n(t)=(P k)(nt)-1/n(P k)(t) (t>0).}
\expandafter\gdef\csname CGTip@eq:approxmellin\endcsname{Equation: int\_0\textasciicircum{}infinity( k(t)-sum\_(n<= N)c\_n\_n(t))t\textasciicircum{}(s-1)dt=Fhat(s)(1-zeta(s) D(s)) (Re s=sigma),}
\expandafter\gdef\csname CGTip@thm:nb\endcsname{Theorem - Weighted strong criterion: Let 1/2<sigma<1. (i) If k lies in the H\_sigma-closure of the linear span of (\_n:n>=1), then zeta(s) != 0 for Re s>sigma. (ii) Assume RH. Then for every eps>0, \| k-sum\_(n<= N)mu(n)\_n\|\_sigma << \_(sigma,eps)N\textasciicircum{}(-(sigma-1/2}
\expandafter\gdef\csname CGTip@eq:nbrate\endcsname{Equation: \| k-sum\_(n<= N)mu(n)\_n\|\_sigma << \_(sigma,eps)N\textasciicircum{}(-(sigma-1/2)+eps) (N>=2).}
\expandafter\gdef\csname CGTip@eq:nbpoint\endcsname{Equation: |Phi\_D(rho)|\textasciicircum{}2<=(eta\textasciicircum{}2)/(2(Re rho-sigma)).}
\expandafter\gdef\csname CGTip@eq:RHbounds\endcsname{Equation: |zeta(u)| << \_(delta,eps)|Im u|\textasciicircum{}eps, |1/zeta(u)| << \_(delta,eps)|Im u|\textasciicircum{}eps:}
\expandafter\gdef\csname CGTip@eq:perron\endcsname{Equation: D\_N(s):=sum\_(n<= N)(mu(n))/(n\textasciicircum{}s)=1/(2pi i)int\_(c-iT)\textasciicircum{}(c+iT)(x\textasciicircum{}w)/(zeta(s+w))(dw)/w+O((x\textasciicircum{}c)/(T))+O(x\textasciicircum{}(1-sigma)log xT),}
\expandafter\gdef\csname CGTip@eq:perronvert\endcsname{Equation: |int(x\textasciicircum{}w)/(zeta(s+w))(dw)/w| << x\textasciicircum{}(-eta\_0)int\_(-T)\textasciicircum{}T((2+|tau|+|v|)\textasciicircum{}eps)/(eta\_0+|v|) dv << \_(eta\_0)N\textasciicircum{}(-eta\_0)(N\textasciicircum{}2+|tau|)\textasciicircum{}epslog N.}
\expandafter\gdef\csname CGTip@eq:DNbound\endcsname{Equation: |D\_N(s)-1/(zeta(s))| << \_(sigma,delta,eps)N\textasciicircum{}(-eta\_0+3eps)(2+|tau|)\textasciicircum{}eps (Re s=sigma).}
\expandafter\gdef\csname CGTip@eq:oneminus\endcsname{Equation: |1-zeta(s)D\_N(s)| << N\textasciicircum{}(-eta\_0+3eps)(2+|tau|)\textasciicircum{}(2eps) (Re s=sigma).}
\expandafter\gdef\csname CGTip@eq:nbnorm\endcsname{Equation: \| k-sum\_(n<= N)mu(n)\_n\|\_sigma<=(1/(2pi)int|Fhat(sigma+itau)|\textasciicircum{}2|1-zeta D\_N|\textasciicircum{}2dtau)\textasciicircum{}(1/2)+|m(N)| \|P k\|\_sigma,}
\expandafter\gdef\csname CGTip@prop:primecut\endcsname{Proposition - Finite prime cut: For Y>=2 let S\_Y be the set of squarefree integers all of whose prime factors are <= Y, and \_Y(s)=prod\_(p<= Y)(1-p\textasciicircum{}(-s))=sum\_(n in S\_Y)mu(n)n\textasciicircum{}(-s), zeta\_(>Y)(s)=zeta(s)\_Y(s),}
\expandafter\gdef\csname CGTip@eq:zetaY\endcsname{Equation: \_Y(s)=prod\_(p<= Y)(1-p\textasciicircum{}(-s))=sum\_(n in S\_Y)mu(n)n\textasciicircum{}(-s), zeta\_(>Y)(s)=zeta(s)\_Y(s),}
\expandafter\gdef\csname CGTip@eq:primecut\endcsname{Equation: \_(Y to infinity)int\_(-infinity)\textasciicircum{}infinity|Fhat(sigma+itau)|\textasciicircum{}2 |1-zeta\_(>Y)(sigma+itau)|\textasciicircum{}2 dtau=0,}
\expandafter\gdef\csname CGTip@eq:primecutnorm\endcsname{Equation: \| k-sum\_(n in S\_Y)mu(n)\_n\|\_sigma<=(1/(2pi)int|Fhat(sigma+itau)|\textasciicircum{}2|1-zeta\_(>Y)(sigma+itau)|\textasciicircum{}2dtau)\textasciicircum{}(1/2)+(\|P k\|\_sigma)/(log Y),}
\expandafter\gdef\csname CGTip@sec:liouville\endcsname{Section: The corner as a weighted Liouville sum, and the Clausen kernel}
\expandafter\gdef\csname CGTip@eq:aU\endcsname{Equation: Corner(X)=sum\_(U<n<= X)a\_U(n) n E\_G( Xn), a\_U(n)=sum\_dj=n d<= U, j<= Umu(d)mu(j).}
\expandafter\gdef\csname CGTip@prop:liouville\endcsname{Proposition - Liouville form of the corner: For n>=1 and U>=1, a\_U(n)=lambda(n) b\_U(n) with b\_U(n)>=0 the number of factorizations n=dj with d,j<= U both squarefree. Explicitly, b\_U(n)=0 unless n=ab\textasciicircum{}2 with a,b squarefree and (a,b)=1, and then b\_U(n)=\(t| a: (ab)/U}
\expandafter\gdef\csname CGTip@eq:bU\endcsname{Equation: b\_U(n)=\(t| a: (ab)/U<= t<= Ub).}
\expandafter\gdef\csname CGTip@eq:Cliouville\endcsname{Equation: Corner(X)=sum\_(sqrt X<n<= X)lambda(n) b\_(sqrt X)(n) n E\_G( Xn),}
\expandafter\gdef\csname CGTip@prop:clausenmuntz\endcsname{Proposition - The Clausen kernel and its M\"untz transform: Let c(t)=t Cl2(2pi/t) for 0<t<=1 and c(t)=0 for t>1. Then int\_0\textasciicircum{}infinity c(t)t\textasciicircum{}(s-1)dt=Theta(s) for Re s>0, int\_0\textasciicircum{}infinity c=a, and for Y>=1, Y, (E\_G(Y))/Y=sum\_(l>=1) c( Y)-Yint\_0\textasciicircum{}infinity c=(P c)(1/Y),}
\expandafter\gdef\csname CGTip@eq:EGmuntz\endcsname{Equation: (E\_G(Y))/Y=sum\_(l>=1) c( Y)-Yint\_0\textasciicircum{}infinity c=(P c)(1/Y),}
\expandafter\gdef\csname CGTip@eq:Cmuntz\endcsname{Equation: Corner(X)=Xsum\_(sqrt X<n<= X)lambda(n) b\_(sqrt X)(n) (P c)( nX),}
\expandafter\gdef\csname CGTip@eq:LU\endcsname{Equation: L\_U(s)=sum\_(U<n<= X)lambda(n)b\_U(n)n\textasciicircum{}(-s)=P\_U(s)\textasciicircum{}2-sum\_(n<= U)(mu*mu)(n)n\textasciicircum{}(-s),}
\expandafter\gdef\csname CGTip@eq:Cspectral\endcsname{Equation: Corner(X)=1/(2pi i)int\_((5/2))zeta(s-1)Theta(s) L\_U(s-1) X\textasciicircum{}s ds-aX\textasciicircum{}2sum\_(U<n<= X)(lambda(n)b\_U(n))/n.}
\expandafter\gdef\csname CGTip@prop:lambdacrit\endcsname{Proposition - One-sided Liouville--Clausen criterion: Let Lambda\_G(X)=sum\_(n<= X)lambda(n) n E\_G(X/n) for X>=1. Suppose that for every eps>0 there is K\_eps such that either Lambda\_G(X)<= K\_eps X\textasciicircum{}(3/2+eps) for all X>=1, or Lambda\_G(X)>=-K\_eps X\textasciicircum{}(3/2+eps) for all X>=1. Then R}
\expandafter\gdef\csname CGTip@eq:lambdatrivial\endcsname{Equation: |Lambda\_G(X)|<=sum\_(n<= X)n c\_(1/2)( Xn)\textasciicircum{}(3/2)<= c\_(1/2)X\textasciicircum{}(3/2)sum\_(n<= X)n\textasciicircum{}(-1/2)<=2c\_(1/2)X\textasciicircum{}2,}
\expandafter\gdef\csname CGTip@eq:lambdamellin\endcsname{Equation: \_G(s)=(zeta(s-1)Theta(s)-a/(s-2))(zeta(2s-2))/(zeta(s-1)) (Re s>2).}
\expandafter\gdef\csname CGTip@rem:liouvillenum\endcsname{Remark - Measurements on the Liouville form: At X=30000.5, 120000.5, 380000.5 (liouville\_weight.py): b\_(sqrt X)(n)>0 for 14.2\, 13.9\, 13.8\ of the n in (sqrt X,X], with mean 2.48, 2.61, 2.66 on its support and maximum 16, 22, 26. The contributions of the n with a }
\expandafter\gdef\csname CGTip@rem:musqcorner\endcsname{Remark - The mu\textasciicircum{}2-corner: Let Corner\_(mu\textasciicircum{}2)(X) be eq:Rest with mu(d)mu(j) replaced by mu(d)\textasciicircum{}2mu(j)\textasciicircum{}2. Replacing the sum over squarefree d=Us, j=Ut by the integral with density 6/pi\textasciicircum{}2 in each variable suggests Corner\_(mu\textasciicircum{}2)(X) X\textasciicircum{}2 with kappa =(6/(}
\expandafter\gdef\csname CGTip@eq:kappa\endcsname{Equation: kappa =(6/(pi\textasciicircum{}2))\textasciicircum{}2int\_0\textasciicircum{}1 int\_0\textasciicircum{}1st E\_G(1/(st))ds dt=(6/(pi\textasciicircum{}2))\textasciicircum{}2int\_1\textasciicircum{}infinity(E\_G(y)log y)/(y\textasciicircum{}3) dy, int\_1\textasciicircum{}infinity(E\_G(y)log y)/(y\textasciicircum{}3) dy =- E\_G'(2)=-(Theta''(2))/2-gamma\_ETheta'(2)+gamma\_1a,}
\expandafter\gdef\csname CGTip@sec:centred\endcsname{Section: The value of the transform of E\_G at the zeros, and the centred primitive of the corner}
\expandafter\gdef\csname CGTip@prop:EGrho\endcsname{Proposition - E\_G at 1+rho: The integral E\_G(s)=int\_1\textasciicircum{}infinity E\_G(Y)Y\textasciicircum{}(-s-1)dY converges absolutely for Re s>1, and there E\_G(s)=zeta(s-1)Theta(s)-a/(s-2),}
\expandafter\gdef\csname CGTip@eq:EGhat\endcsname{Equation: E\_G(s)=zeta(s-1)Theta(s)-a/(s-2),}
\expandafter\gdef\csname CGTip@eq:EGrho\endcsname{Equation: E\_G(1+rho)=-a/(rho-1).}
\expandafter\gdef\csname CGTip@rem:pairs\endcsname{Remark - Bands and zeros: Proposition prop:EGrho is the exact form of the following computation. Replacing mu(d) formally by sum\_rho d\textasciicircum{}(rho-1)/zeta'(rho) in the corner, the substitution d=sqrt X x, j=sqrt X y, p=xy, q=x/y shows that the pair (rho}
\expandafter\gdef\csname CGTip@eq:Jbands\endcsname{Equation: J\_k(rho) =1gammaint\_(2\textasciicircum{}k)\textasciicircum{}2\textasciicircum{}(k+1)E\_G(Y) Y\textasciicircum{}(-beta-2)sin(gammalog Y) dY, J(rho) =sum\_(k>=0)J\_k(rho)=-Im E\_G(1+rho)gamma=-a/(|rho-1|\textasciicircum{}2).}
\expandafter\gdef\csname CGTip@eq:DUdef\endcsname{Equation: D\_U(n)=sum\_d,j<= U dj>n(mu(d)mu(j))/(dj)=m(U)\textasciicircum{}2-sum\_d,j<= U dj<= n(mu(d)mu(j))/(dj)=sum\_(n<k<= U\textasciicircum{}2)(lambda(k) b\_U(k))/k,}
\expandafter\gdef\csname CGTip@lem:Pbound\endcsname{Lemma: Assume RH and let eps>0. Then, uniformly in U>=1 and t in R, |P\_U(1/2+it)|=|sum\_(d<= U)mu(d)d\textasciicircum{}(1/2+it)| << \_eps U\textasciicircum{}eps(2+|t|)\textasciicircum{}eps .}
\expandafter\gdef\csname CGTip@eq:Pbound\endcsname{Equation: |P\_U(1/2+it)|=|sum\_(d<= U)mu(d)d\textasciicircum{}(1/2+it)| << \_eps U\textasciicircum{}eps(2+|t|)\textasciicircum{}eps .}
\expandafter\gdef\csname CGTip@eq:Mt\endcsname{Equation: M(x t)=1/(2pi i)int\_(c-iT)\textasciicircum{}(c+iT)(x\textasciicircum{}w)/(zeta(w+it))(dw)/w+O((log x)/x).}
\expandafter\gdef\csname CGTip@thm:DU\endcsname{Theorem - Centred primitive criterion: The following are equivalent. (a) RH. (b) For every eps>0 there is K\_eps such that, for all U>=1 and all integers U<= n<U\textasciicircum{}2, |D\_U(n)|<= K\_eps (U\textasciicircum{}eps)/(sqrt n).}
\expandafter\gdef\csname CGTip@eq:DUcrit\endcsname{Equation: |D\_U(n)|<= K\_eps (U\textasciicircum{}eps)/(sqrt n).}
\expandafter\gdef\csname CGTip@eq:DUperron\endcsname{Equation: sum\_(k<= n)(a\_U(k))/k=1/(2pi i)int\_(c-iT)\textasciicircum{}(c+iT)P\_U(1+w)\textasciicircum{}2(x\textasciicircum{}w)/w dw+O\_eps((log\textasciicircum{}2x)/(T)+(x\textasciicircum{}epslog x)/T).}
\expandafter\gdef\csname CGTip@eq:DUshift\endcsname{Equation: D\_U(n)=-1/(2pi i)int\_(-1/2-iT)\textasciicircum{}(-1/2+iT)P\_U(1/2+iv)\textasciicircum{}2(x\textasciicircum{}w)/w dw-1/(2pi i)(int\_(-1/2+iT)\textasciicircum{}(c+iT)-int\_(-1/2-iT)\textasciicircum{}(c-iT))P\_U(1+w)\textasciicircum{}2(x\textasciicircum{}w)/w dw+O\_eps(n\textasciicircum{}(-2+eps)).}
\expandafter\gdef\csname CGTip@eq:DUabel\endcsname{Equation: Corner(X)=D\_U(U) V(U+1)+sum\_(U<k<U\textasciicircum{}2)D\_U(k)( V(k+1)- V(k)).}
\expandafter\gdef\csname CGTip@eq:hprime\endcsname{Equation: V'(xi)=2xi E\_G(X/xi)+2pi X R(X/xi) for almost every xi,}
\expandafter\gdef\csname CGTip@eq:DUestimate\endcsname{Equation: |Corner(X)|<= K\_eps U\textasciicircum{}eps(U\textasciicircum{}(-1/2)| V(U+1)|+sqrt2int\_U\textasciicircum{}Xxi\textasciicircum{}(-1/2)(2xi|E\_G(X/xi)|+2pi X|R(X/xi)|)dxi).}
\expandafter\gdef\csname CGTip@rem:DUnum\endcsname{Remark: Two comments on eq:DUcrit. First, bounding the terms of D\_U(n) one by one does not give it even under RH: for U<= n<U\textasciicircum{}2, D\_U(n)=m(U)\textasciicircum{}2-sum\_(k<= n)(mu*mu)(k)/k+2sum\_(j<n/U)mu(j)j\textasciicircum{}(-1)(m(n/j)-m(U)), and inserting m(y) << y}
\expandafter\gdef\csname CGTip@eq:shifts\endcsname{Equation: m(U)\textasciicircum{}2=sum\_(|h|<U) sum\_d<= U 1<= d+h<= U(mu(d)mu(d+h))/(d(d+h)),}
\expandafter\gdef\csname CGTip@sec:attempts\endcsname{Section: Attempts on the corner and where they stopped}
\expandafter\gdef\csname CGTip@sec:num\endcsname{Section: Numerical observations}
\expandafter\gdef\csname CGTip@tab:corner\endcsname{Table: The exact decomposition eq:decomp, in units of X\textasciicircum{}(3/2) in parentheses the sum of the absolute values of the terms. The last column is the residual of eq:decomp.}
\expandafter\gdef\csname CGTip@tab:bands\endcsname{Table: The exact decomposition eq:decomp, in units of X\textasciicircum{}(3/2) in parentheses the sum of the absolute values of the terms. The last column is the residual of eq:decomp. lrrrrrr to prule X W 2W\textasciicircum{}b -aX\textasciicircum{}2m\textasciicircum{}2 -Rest -Corner residual |}
\expandafter\gdef\csname CGTip@tab:res\endcsname{Table: The exact decomposition eq:decomp, in units of X\textasciicircum{}(3/2) in parentheses the sum of the absolute values of the terms. The last column is the residual of eq:decomp. lrrrrrr to prule X W 2W\textasciicircum{}b -aX\textasciicircum{}2m\textasciicircum{}2 -Rest -Corner residual |}
\expandafter\gdef\csname CGTip@sec:remaining\endcsname{Section: What the route reduces the Riemann hypothesis to}
\expandafter\gdef\csname CGTip@eq:target\endcsname{Equation: sum\_d,j<=sqrt X dj>sqrt Xmu(d)mu(j) dj E\_G(X/(dj))>=-K\_eps X\textasciicircum{}(3/2+eps) (X>=1),}
\expandafter\gdef\csname CGTip@app:code\endcsname{Section: Code and data}
\makeatother

\hypersetup{
  colorlinks=false,
  pdfborder={0 0 1},
  linkbordercolor={1 0 0},
  citebordercolor={0 0.3 1},
  urlbordercolor={0 0.6 0.6},
  filebordercolor={0 0.6 0.6},
  linktoc=all,
  bookmarksnumbered=true,
  bookmarksopen=true,
  bookmarksopenlevel=1,
  pdfstartview=FitH,
  pdftitle={The log-sine Moebius route to the Riemann hypothesis},
  pdfauthor={V\'ictor Casa\~na-Giner},
  pdfsubject={Riemann hypothesis; Moebius function; log-sine and Clausen functions},
  pdfkeywords={Riemann hypothesis, Moebius function, Clausen function, Mellin transform, Mertens function}
}
% Back-references in the bibliography.
\renewcommand*{\backref}[1]{}
\renewcommand*{\backrefalt}[4]{%
  \ifcase #1 %
    \relax
  \or
    \ (Cited on p.~#2.)%
  \else
    \ (Cited on pp.~#2.)%
  \fi}

% In \cite[note]{key} the note (e.g. a section number) is linked too.
\makeatletter
\let\CG@origcite\cite
\renewcommand*{\cite}[2][]{\def\CG@citekey{#2}%
  \ifx\relax#1\relax\CG@origcite{#2}\else\CG@origcite[#1]{#2}\fi}
\def\@cite#1#2{{\@citestyle[\citeform{#1}\if@tempswa, \hyper@link{cite}{cite.\CG@citekey}{#2}\fi]}}
\makeatother

% Keep the review annotations entirely inside the one-inch outer margin.
\geometry{margin=1in}
\setlength{\marginparwidth}{0.72in}
\setlength{\marginparsep}{4pt}
\numberwithin{equation}{section}
\setcounter{tocdepth}{3}
\setlength{\emergencystretch}{1.5em}

\theoremstyle{plain}
\newtheorem{theorem}{Theorem}[section]
\newtheorem{proposition}[theorem]{Proposition}
\newtheorem{lemma}[theorem]{Lemma}
\newtheorem{corollary}[theorem]{Corollary}
\theoremstyle{definition}
\newtheorem{definition}[theorem]{Definition}
\newtheorem{hypothesis}[theorem]{Hypothesis}
\newtheorem{heuristic}[theorem]{Heuristic}
\theoremstyle{remark}
\newtheorem{remark}[theorem]{Remark}

\newcommand{\PP}{\mathbb{P}}
\newcommand{\EE}{\mathbb{E}}
\newcommand{\Var}{\operatorname{Var}}


% CorrectVCG review marks: every statement reviewed by Claude carries \VCGreviewed{n},
% n = how many times that statement has been reviewed.  \VCGshowreviewsfalse hides the marks.
\makeatletter
\@ifundefined{VCGreviewed}{%
\newcounter{VCGreviewcount}
\newif\ifVCGshowreviews \VCGshowreviewstrue
\definecolor{VCGgreen}{rgb}{0,0.45,0.15}
\newcommand{\VCGreviewed}[1]{\stepcounter{VCGreviewcount}\ifVCGshowreviews\marginpar[{\raggedleft\tiny\textcolor{VCGgreen}{Reviewed\\ by Claude\\ (review~#1)}}]{\raggedright\tiny\textcolor{VCGgreen}{Reviewed\\ by Claude\\ (review~#1)}}\fi}
\AtEndDocument{\ifVCGshowreviews\par\medskip\noindent{\small\textcolor{VCGgreen}{Statements reviewed by Claude: \arabic{VCGreviewcount}.}}\fi}
}{}
\makeatother

\begin{document}

\title[The log-sine M\"obius route]{The log-sine M\"obius route to the Riemann hypothesis:\\ what is proved, what is measured, and the one inequality that remains}
\author{V\'ictor Casa\~na-Giner}
\email{vcasana1@alumno.uned.es}
\date{}

\begin{abstract}
Let $F(u)=\log|2\sin\pi u|$ and $H(x)=\sum_{d\le x}\mu(d)F(x/d)$. We record, with the exact logical status of every statement, the route from $H$ to the Riemann hypothesis followed so far. Proved here: the Mellin transform of $F$ on $[1,\infty)$ is $(\pi/s)\cot(\pi s/2)\zeta(s)+A(s)$ with an explicit $A$, and $A$ has no zeros in $0<\operatorname{Re}s<1$; consequently the Riemann hypothesis is equivalent to $\int_1^XH(x)^2\,dx\ll X^{2+\varepsilon}$, and to $W(X)\ll X^{3/2+\varepsilon}$ for the Clausen sum $W(X)=\sum_{d\le X}\mu(d)\,d\,\mathrm{Cl}_2(2\pi X/d)=-2\pi\int_1^XH$, and also to each of the one-sided bounds $W(X)\le K_\varepsilon X^{3/2+\varepsilon}$ and $W(X)\ge-K_\varepsilon X^{3/2+\varepsilon}$. Unconditionally $W(X)/X^{3/2}$ has $\limsup\ge0.0821$ and $\liminf\le-0.0821$, the constant being computed numerically from the first zero of $\zeta$. The identity $\mu=2\mu_U-\mu_U*\mu_U*1$ on $[1,U^2]$ with $U=\sqrt X$ splits $W$ exactly into a part of size $O(X)$, a part $-aX^2m(\sqrt X)^2\le0$ with $a=\pi(C-\log2)>0$, a part $O(X^{3/2+\varepsilon})$, and a corner sum $\mathcal C(X)$ that depends only on $\mu$ restricted to $[1,\sqrt X]$. Hence the Riemann hypothesis is equivalent to the single one-sided inequality $\mathcal C(X)\ge-K_\varepsilon X^{3/2+\varepsilon}$. That inequality is not proved. The M\"untz transform of the kernel $\mathfrak k(t)=F(1/t)\mathbf 1_{(0,1]}(t)$ is $R(1/t)$, with $R=-E_G'/(2\pi)$, its M\"obius dilation sum is $H(1/t)$, and the energy $\int|\sum_{d\le N}\mu(d)\mathfrak k(dt)|^2t^{2\sigma-1}dt$ is an explicit positive semidefinite quadratic form in $\mu$ whose uniform boundedness for every $\sigma>\frac12$ implies RH (converse not proved). With $\varrho_n=(P\mathfrak k)(n\,\cdot)-(P\mathfrak k)/n$ and $\mathcal H_\sigma=L^2(t^{2\sigma-1}dt)$, RH is equivalent to $\mathfrak k$ lying in the $\mathcal H_\sigma$-closure of $\operatorname{span}\{\varrho_n\}$ for every $\sigma\in(\frac12,1)$, a weighted strong Nyman--Beurling criterion for a kernel that is not continuously differentiable. Finally, RH is equivalent to $|D_U(n)|\ll U^\varepsilon n^{-1/2}$ for $U\le n<U^2$, where $D_U(n)=\sum_{d,j\le U,\,dj>n}\mu(d)\mu(j)/(dj)$ is the centred primitive of the corner. We also record what was measured numerically and which methods were tried and why they do not reach the corner inequality.
\end{abstract}

\subjclass[2020]{11M26; 11N37; 11M06; 33E20}
\keywords{Riemann hypothesis; M\"obius function; log-sine function; Clausen function; Mellin transform; Mertens function}

\maketitle
\tableofcontents

\section{Introduction}\label{sec:intro}

This document records the state of one line of work. Its purpose is to say exactly what has been proved, in which direction, and under which hypothesis, and to separate it from what has only been measured. Nothing in it is used in a direction that has not been proved. The Riemann hypothesis (RH) is used in the form: $\zeta(s)\neq0$ for $\operatorname{Re}s>\tfrac12$.

Throughout, $L^1$ and $L^2$ are the Lebesgue spaces, $F(u)=\log|2\sin\pi u|$ for $u\in\mathbb R\setminus\mathbb Z$, $\mu$ is the M\"obius function, $M(x)=\sum_{n\le x}\mu(n)$, $m(x)=\sum_{n\le x}\mu(n)/n$, and $\mathrm{Cl}_2$ is the Clausen function \CGeqref{eq:Cl2}. The constant
\begin{equation}\label{eq:Cdef}
C=-\int_1^\infty\log|\sin\pi u|\,u^{-2}\,du=0.74379486183940054\ldots
\end{equation}
and $a=\pi(C-\log2)=0.15911438342889\ldots$ (Corollary~\CGref{cor:B}).

The table below lists every item and its status. ``Proved'' means proved in this document, with the classical inputs cited at the step where they are used. ``Measured'' means a floating-point computation, reported with its range, with no claim beyond it. ``Not proved'' means exactly that.

\begin{center}
\small
\begin{tabular}{p{0.62\textwidth}p{0.28\textwidth}}
\toprule
Item & Status \\
\midrule
Mellin transform of $F$, Theorem~\CGref{thm:mellin} & proved \\
$A(s)\neq0$ for $0<\operatorname{Re}s<1$, Theorem~\CGref{thm:Anonzero} & proved \\
Local structure of $H$ and the two identities, Propositions~\CGref{prop:Hlocal}, \CGref{prop:Hid} & proved \\
RH $\iff\int_1^XH^2\ll X^{2+\varepsilon}$, Theorem~\CGref{thm:L2} & proved, both directions \\
RH $\iff W\ll X^{3/2+\varepsilon}\iff$ either one-sided bound, Theorem~\CGref{thm:Wcrit} & proved, all directions \\
$\Omega_\pm$ bounds for $W/X^{3/2}$, Theorem~\CGref{thm:omega} & proved; the constant $0.0821$ is computed numerically \\
Exact $d=2$ decomposition of $W$, Proposition~\CGref{prop:decomp} & proved \\
$E_G(Y)\ll_\delta Y^{1+\delta}$, Lemma~\CGref{lem:EG} & proved \\
RH $\iff\mathcal C(X)\ge-K_\varepsilon X^{3/2+\varepsilon}\iff|\mathcal C(X)|\ll X^{3/2+\varepsilon}$, Theorem~\CGref{thm:corner} & proved, all directions \\
Other exact identities (lcm form, recursion over large primes, $d=3$), Section~\CGref{sec:other} & proved \\
$W_{\mu^2}(X)=(6/\pi^2)aX^2+O(X^{3/2+\delta})$, Proposition~\CGref{prop:musq} & proved \\
M\"untz form of $R$ and $H$, energy and Gram matrix, Propositions~\CGref{prop:muntz}, \CGref{prop:energy} & proved \\
$\liminf_N\mathcal E_N(\sigma)<\infty\ \forall\sigma>\frac12\Rightarrow$ RH, Theorem~\CGref{thm:energy} & proved, one direction only \\
Weighted strong Nyman--Beurling criterion, Theorem~\CGref{thm:nb} & proved, both directions \\
Finite prime cut $\Rightarrow$ RH, Proposition~\CGref{prop:primecut} & proved, one direction; converse \textbf{not proved} \\
Liouville form of the corner and Clausen--M\"untz structure, Propositions~\CGref{prop:liouville}, \CGref{prop:clausenmuntz} & proved \\
One-sided Liouville--Clausen criterion, Proposition~\CGref{prop:lambdacrit} & proved, one direction \\
$\mathcal C_{\mu^2}\sim\kappa X^2$, Remark~\CGref{rem:musqcorner} & constant computed; asymptotic \textbf{not proved}, not confirmed by the data \\
$\widehat E_G(1+\rho)=-a/(\rho-1)$, Proposition~\CGref{prop:EGrho} & proved \\
RH $\iff|D_U(n)|\ll U^\varepsilon n^{-1/2}$ ($U\le n<U^2$), Theorem~\CGref{thm:DU} & proved, both directions \\
Corner inequality \CGeqref{eq:target} & \textbf{not proved} \\
Centred primitive bound \CGeqref{eq:DUcrit} & \textbf{not proved} \\
Off-diagonal bound \CGeqref{eq:target2} & \textbf{not proved}; sufficient for RH, not known to be necessary \\
Localization inequality of Remark~\CGref{rem:local} & \textbf{not proved}, not used \\
Tables of Section~\CGref{sec:num} & measured \\
\bottomrule
\end{tabular}
\end{center}

Every criterion below is a reformulation of RH, and it is stated with both directions proved. The reduction of Section~\CGref{sec:corner} replaces the two-sided bound for $W$ by a one-sided bound for a sum that involves $\mu$ only on $[1,\sqrt X]$. The direction used to reach RH, corner bound $\Rightarrow$ RH, needs no information on $m(\sqrt X)$ or on any tail of $m$: the term $-aX^2m(\sqrt X)^2$ has a sign and is simply dropped. Information on $m$ is used only in the converse direction, RH $\Rightarrow$ corner bound.

\section{The Mellin transform of the log-sine function}\label{sec:mellin}

For $\theta\in\mathbb R$ let
\begin{equation}\label{eq:Cl2}
\mathrm{Cl}_2(\theta)=-\int_0^\theta\log\Bigl|2\sin\frac\psi2\Bigr|\,d\psi=\sum_{\nu\ge1}\frac{\sin\nu\theta}{\nu^2}.
\end{equation}
The second equality is \CGeqref{eq:Phi} below with $\theta=2\pi u$. The function $\mathrm{Cl}_2$ is odd, $2\pi$-periodic and continuous. Its derivative $-\log|2\sin(\theta/2)|$ vanishes on $(0,\pi)$ only at $\theta=\pi/3$, is positive on $(0,\pi/3)$ and negative on $(\pi/3,5\pi/3)$. Hence
\begin{equation}\label{eq:Clmax}
|\mathrm{Cl}_2(\theta)|\le\mathrm{Cl}_2(\pi/3)=1.01494160640965\ldots\qquad(\theta\in\mathbb R).
\end{equation}

\begin{lemma}[Euler integrals]\label{lem:euler}\VCGreviewed{1}
Let $\alpha>0$. For $-1<\operatorname{Re}w<0$ the integral below converges absolutely and
\begin{equation}\label{eq:euler}
\int_0^\infty u^{w-1}\sin(\alpha u)\,du=\alpha^{-w}\,\Gamma(w)\sin\frac{\pi w}2.
\end{equation}
\end{lemma}

\begin{proof}
First let $0<\operatorname{Re}z<1$. We show that
\begin{equation}\label{eq:eulerz}
\lim_{R\to\infty}\int_0^Ru^{z-1}e^{\pm iu}\,du=\Gamma(z)\,e^{\pm i\pi z/2}.
\end{equation}
Take the upper sign. Integrate $\xi^{z-1}e^{i\xi}$, with the principal branch, over the boundary of $\{\eta\le|\xi|\le R,\ 0\le\arg\xi\le\pi/2\}$, where $0<\eta<R$. Cauchy's theorem gives $0$. On the arc $\xi=Re^{i\theta}$ we have $|e^{i\xi}|=e^{-R\sin\theta}\le e^{-2R\theta/\pi}$ and $|\xi^{z-1}|\le R^{\operatorname{Re}z-1}e^{\pi|\operatorname{Im}z|/2}$, so the arc contributes at most
\begin{equation}\label{eq:arc}
e^{\pi|\operatorname{Im}z|/2}R^{\operatorname{Re}z}\int_0^{\pi/2}e^{-2R\theta/\pi}\,d\theta\le e^{\pi|\operatorname{Im}z|/2}\,\frac\pi2\,R^{\operatorname{Re}z-1},
\end{equation}
which tends to $0$ because $\operatorname{Re}z<1$. The arc of radius $\eta$ contributes at most $e^{\pi|\operatorname{Im}z|/2}\frac\pi2\eta^{\operatorname{Re}z}$, which tends to $0$ as $\eta\to0$ because $\operatorname{Re}z>0$. On the imaginary axis $\xi=iy$ the integral is $i^z\int_0^\infty y^{z-1}e^{-y}\,dy=e^{i\pi z/2}\Gamma(z)$. This proves \CGeqref{eq:eulerz} with the upper sign. For the lower sign integrate $\xi^{z-1}e^{-i\xi}$ over the boundary of $\{\eta\le|\xi|\le R,\ -\pi/2\le\arg\xi\le0\}$. On the arc $\xi=Re^{i\theta}$, $-\pi/2\le\theta\le0$, we have $|e^{-i\xi}|=e^{R\sin\theta}\le e^{-2R|\theta|/\pi}$, so both arcs obey the bounds above, and on the axis $\xi=-iy$ the integral is $(-i)^z\Gamma(z)=e^{-i\pi z/2}\Gamma(z)$. Half the sum and half the difference divided by $i$ give $\int_0^\infty u^{z-1}\cos u\,du=\Gamma(z)\cos(\pi z/2)$ and $\int_0^\infty u^{z-1}\sin u\,du=\Gamma(z)\sin(\pi z/2)$ as improper integrals.

Now let $-1<\operatorname{Re}w<0$. The integrand in \CGeqref{eq:euler} with $\alpha=1$ is $O(u^{\operatorname{Re}w})$ near $0$ and $O(u^{\operatorname{Re}w-1})$ at infinity, so it is absolutely integrable. Integration by parts gives
\begin{equation}\label{eq:parts}
\int_0^\infty u^{w-1}\sin u\,du=\Bigl[\frac{u^w\sin u}{w}\Bigr]_0^\infty-\frac1w\int_0^\infty u^{w}\cos u\,du=-\frac1w\,\Gamma(w+1)\cos\frac{\pi(w+1)}2,
\end{equation}
where $w\neq0$ because $\operatorname{Re}w<0$, the boundary terms vanish because $-1<\operatorname{Re}w<0$, and the last integral is the cosine case with $z=w+1$, $0<\operatorname{Re}z<1$. Since $\Gamma(w+1)=w\Gamma(w)$ and $\cos(\pi(w+1)/2)=-\sin(\pi w/2)$, this is $\Gamma(w)\sin(\pi w/2)$. The substitution $u\mapsto\alpha u$ gives the factor $\alpha^{-w}$.
\end{proof}

For $\operatorname{Re}s<2$, $s\neq0$, let
\begin{equation}\label{eq:Adef}
A(s)=\frac{\log(2\pi)}s+\frac1{s^2}+\sum_{k\ge1}\frac{\zeta(2k)}{k\,(2k-s)}.
\end{equation}
The series converges absolutely and locally uniformly on $\operatorname{Re}s<2$, because $\zeta(2k)\le\zeta(2)$ and $|2k-s|\ge2k-\operatorname{Re}s\ge2(k-1)+(2-\operatorname{Re}s)$; so $A$ is holomorphic on $\{\operatorname{Re}s<2\}\setminus\{0\}$.

\begin{theorem}[Mellin transform of $F$]\label{thm:mellin}\VCGreviewed{1}
For $\operatorname{Re}s>0$ the integral
\begin{equation}\label{eq:Fhatdef}
\widehat F(s)=\int_1^\infty F(u)\,u^{-s-1}\,du
\end{equation}
converges absolutely. The function $\widehat F$ extends holomorphically to $\operatorname{Re}s>-1$, and for $0<\operatorname{Re}s<2$
\begin{equation}\label{eq:Fhat}
\widehat F(s)=\frac\pi s\cot\Bigl(\frac{\pi s}2\Bigr)\zeta(s)+A(s),
\end{equation}
where at $s=1$ the first term is defined by continuity, with value $-\pi^2/2$.
\end{theorem}

\begin{proof}
Absolute convergence for $\operatorname{Re}s>0$ holds because $F$ is $1$-periodic and integrable over a period. For $0<\lambda<1$ and real $v$ one has $\log|1-\lambda e^{2\pi iv}|=-\sum_{\nu\ge1}\lambda^\nu\cos(2\pi\nu v)/\nu$. As $\lambda\to1^-$ the left side tends to $F(v)$ in $L^2(0,1)$ by dominated convergence, and the Fourier coefficients on the right tend to $-1/(2\nu)$. Hence
\begin{equation}\label{eq:fourier}
F(v)=-\sum_{\nu\ge1}\frac{\cos2\pi\nu v}{\nu}\qquad\text{in }L^2(0,1).
\end{equation}
In particular $\int_0^1F=0$ and, by Parseval, $\int_0^1F^2=\sum_\nu1/(2\nu^2)=\pi^2/12$. Integrating \CGeqref{eq:fourier} over $[0,u]$ and using periodicity, the function $\Phi(u)=\int_0^uF(v)\,dv$ satisfies
\begin{equation}\label{eq:Phi}
\Phi(u)=-\frac1{2\pi}\sum_{\nu\ge1}\frac{\sin2\pi\nu u}{\nu^2}=-\frac{\mathrm{Cl}_2(2\pi u)}{2\pi}\qquad(u\ge0),
\end{equation}
with an absolutely and uniformly convergent series. The second equality is the definition \CGeqref{eq:Cl2} after the substitution $\psi=2\pi v$. Thus $\Phi$ is bounded, $\Phi(n)=0$ for every integer $n\ge0$, and, since $2\sin\pi v\ge4v$ on $[0,\tfrac12]$,
\begin{equation}\label{eq:Phi0}
|\Phi(u)|\le\int_0^u\log\frac1{4v}\,dv=u\Bigl(\log\frac1{4u}+1\Bigr)\qquad(0<u\le\tfrac16).
\end{equation}

By parts, using $\Phi(1)=0$ and the boundedness of $\Phi$, for $\operatorname{Re}s>0$
\begin{equation}\label{eq:Gparts}
\widehat F(s)=(s+1)\int_1^\infty\Phi(u)\,u^{-s-2}\,du,
\end{equation}
and the right-hand side converges absolutely and is holomorphic for $\operatorname{Re}s>-1$. This gives the extension.

Let $\sigma=\operatorname{Re}s\in(-1,0)$. Integration by parts, with boundary terms vanishing by \CGeqref{eq:Phi0} at $0$ and by boundedness of $\Phi$ at infinity, gives the convergent improper integral
\begin{equation}\label{eq:L}
\mathcal L(s):=\int_0^\infty F(u)\,u^{-s-1}\,du=(s+1)\int_0^\infty\Phi(u)\,u^{-s-2}\,du.
\end{equation}
Substituting \CGeqref{eq:Phi} and integrating termwise is justified by Fubini's theorem, because the series of absolute values converges by Tonelli's theorem:
\begin{equation}\label{eq:tonelli}
\sum_{\nu\ge1}\frac1{\nu^2}\int_0^\infty|\sin2\pi\nu u|\,u^{-\sigma-2}\,du=b_\sigma(2\pi)^{\sigma+1}\sum_{\nu\ge1}\nu^{\sigma-1}<\infty,
\end{equation}
where $b_\sigma=\int_0^\infty|\sin v|\,v^{-\sigma-2}\,dv$ is finite because $-1<\sigma<0$. Lemma~\CGref{lem:euler} with $w=-s-1$ and $\alpha=2\pi\nu$, together with $(s+1)\Gamma(-s-1)=-\Gamma(-s)$ and $\sin(-\pi(s+1)/2)=-\cos(\pi s/2)$, gives
\begin{equation}\label{eq:Lzeta}
\mathcal L(s)=-(2\pi)^s\,\Gamma(-s)\cos\Bigl(\frac{\pi s}2\Bigr)\zeta(1-s)=\frac\pi s\cot\Bigl(\frac{\pi s}2\Bigr)\zeta(s)\qquad(-1<\sigma<0).
\end{equation}
The second equality uses the functional equation $\zeta(1-s)=2(2\pi)^{-s}\cos(\pi s/2)\Gamma(s)\zeta(s)$ \cite[(2.1.8), p.~16]{TitchmarshHB1986} and Euler's reflection formula, which gives $\Gamma(s)\Gamma(-s)=-\pi/(s\sin\pi s)$.

Split $\mathcal L(s)=\widehat F(s)+\mathcal I(s)$ according to $u\ge1$ and $u<1$; here $\widehat F$ is given by \CGeqref{eq:Gparts}, and the split is legitimate because \CGeqref{eq:L} and \CGeqref{eq:Gparts} are obtained by the same integration by parts on $[1,\infty)$. For $\operatorname{Re}s<0$, writing $F(u)=\log(2\pi u)+\log(\sin\pi u/(\pi u))$ on $(0,1)$,
\begin{equation}\label{eq:I}
\mathcal I(s):=\int_0^1F(u)\,u^{-s-1}\,du=-\frac{\log(2\pi)}s-\frac1{s^2}+\mathcal J(s),
\end{equation}
\begin{equation}\label{eq:J}
\mathcal J(s):=\int_0^1\log\Bigl(\frac{\sin\pi u}{\pi u}\Bigr)u^{-s-1}\,du=-\sum_{k\ge1}\frac{\zeta(2k)}{k\,(2k-s)}\qquad(\operatorname{Re}s<2).
\end{equation}
For \CGeqref{eq:J} we use Euler's product for the sine, which gives, for $0\le u<1$,
\begin{equation}\label{eq:sinprod}
\log\frac{\sin\pi u}{\pi u}=\sum_{\nu\ge1}\log\Bigl(1-\frac{u^2}{\nu^2}\Bigr)=-\sum_{k\ge1}\frac{\zeta(2k)}{k}\,u^{2k}.
\end{equation}
All terms have the same sign, and $\sum_k\zeta(2k)/(k(2k-\operatorname{Re}s))<\infty$ for $\operatorname{Re}s<2$, so termwise integration is justified by Fubini's theorem, the absolute series being bounded through Tonelli's theorem, and $\mathcal J$ is holomorphic for $\operatorname{Re}s<2$.

Combining \CGeqref{eq:Lzeta} and \CGeqref{eq:I}, for $-1<\operatorname{Re}s<0$,
\begin{equation}\label{eq:Gid}
\widehat F(s)=\frac\pi s\cot\Bigl(\frac{\pi s}2\Bigr)\zeta(s)+\frac{\log(2\pi)}s+\frac1{s^2}-\mathcal J(s).
\end{equation}
By \CGeqref{eq:Adef} and \CGeqref{eq:J}, the last three terms are $A(s)$. Both sides are meromorphic on the connected strip $-1<\operatorname{Re}s<2$, so \CGeqref{eq:Gid} holds there by the identity theorem, and this is \CGeqref{eq:Fhat}. At $s=1$, $\cot(\pi s/2)=-\tan(\pi(s-1)/2)$ and $(s-1)\zeta(s)\to1$ \cite[(2.1.16), p.~16]{TitchmarshHB1986}, so the first term tends to $-\pi^2/2$.
\end{proof}

\begin{corollary}[Two forms of $sA(s)$, and the constants]\label{cor:B}\VCGreviewed{1}
Let $B(s)=sA(s)$. For $\operatorname{Re}s<2$, $s\neq0$,
\begin{equation}\label{eq:B1}
B(s)=\log(2\pi)+\frac1s+\sum_{k\ge1}\zeta(2k)\Bigl(\frac2{2k-s}-\frac1k\Bigr),
\end{equation}
\begin{equation}\label{eq:B2}
B(s)=\log\pi+\frac1s+\sum_{k\ge1}\frac{s/k+2(\zeta(2k)-1)}{2k-s},
\end{equation}
with absolutely convergent series. Moreover $\widehat F(1)=\log2-C$, and
\begin{equation}\label{eq:aconst}
a:=\pi(C-\log2)=-\pi\widehat F(1)=0.15911438342889\ldots>0.
\end{equation}
\end{corollary}

\begin{proof}
For $k\ge1$, $s/(k(2k-s))=2/(2k-s)-1/k$; multiplying \CGeqref{eq:Adef} by $s$ gives \CGeqref{eq:B1}, whose terms are $\zeta(2k)s/(k(2k-s))=O(k^{-2})$. Next, $\sum_{k\ge1}(\zeta(2k)-1)/k=\sum_{\nu\ge2}-\log(1-\nu^{-2})$ by Tonelli's theorem, and $\prod_{\nu=2}^N\nu^2/((\nu-1)(\nu+1))=2N/(N+1)\to2$ as $N\to\infty$, so
\begin{equation}\label{eq:log2}
\sum_{k\ge1}\frac{\zeta(2k)-1}k=\log2.
\end{equation}
Writing $\zeta(2k)=1+(\zeta(2k)-1)$ in \CGeqref{eq:B1}, the part with $1$ is $\sum_k(2/(2k-s)-1/k)=\sum_k(s/k)/(2k-s)$, and the part with $\zeta(2k)-1$ is $\sum_k2(\zeta(2k)-1)/(2k-s)-\log2$ by \CGeqref{eq:log2}. This gives \CGeqref{eq:B2}. Finally, $-\log|\sin\pi u|=\log2-F(u)$ and $\int_1^\infty u^{-2}\,du=1$, so \CGeqref{eq:Cdef} gives $C=\log2-\widehat F(1)$. The numerical values of $C$ and $a$ are computed from \CGeqref{eq:Fhat} at $s=1$, which gives $\widehat F(1)=-\pi^2/2+A(1)$, by the script \texttt{route\_constants.py}; the same script checks \CGeqref{eq:Fhat} against direct numerical integration of \CGeqref{eq:Gparts} at $s=\tfrac12+3i$ and $s=\tfrac32+2i$ (agreement to $10^{-8}$, the size of the truncation of the integral).
\end{proof}

\section{The function \texorpdfstring{$A$}{A} has no zeros in the critical strip}\label{sec:Anonzero}

\begin{theorem}\label{thm:Anonzero}\VCGreviewed{1}
$A(s)\neq0$ for $0<\operatorname{Re}s<1$. More precisely, with $s=\sigma+it$, $0<\sigma<1$,
\begin{equation}\label{eq:caseI}
\operatorname{Re}B(s)>\log\pi-\frac{\zeta(3)}2=0.54370\ldots\qquad(|t|\le\sqrt2),
\end{equation}
\begin{equation}\label{eq:caseII}
\frac{\operatorname{Im}B(s)}t=-\frac1{\sigma^2+t^2}+\sum_{k\ge1}\frac{2\zeta(2k)}{(2k-\sigma)^2+t^2}>0\qquad(|t|\ge\sqrt2).
\end{equation}
\end{theorem}

\begin{proof}
Since $A=B/s$ and $s\neq0$, $A(s)=0$ if and only if $B(s)=0$. We show $B(s)\neq0$.

Case $|t|\le\sqrt2$. Use \CGeqref{eq:B2}. For $k\ge1$ the term $2(\zeta(2k)-1)/(2k-s)$ has real part $2(\zeta(2k)-1)(2k-\sigma)/|2k-s|^2\ge0$, because $\zeta(2k)>1$ and $2k-\sigma>0$. Also $\operatorname{Re}(1/s)=\sigma/|s|^2>0$. For the term $s/(k(2k-s))$ we use the identity, valid for all real $\sigma,t$ and $k\ge1$,
\begin{equation}\label{eq:caseIid}
\operatorname{Re}\frac{s}{2k-s}+\frac{t^2}{4k^2+t^2}=\frac{2k\sigma\bigl(2k(2k-\sigma)-t^2\bigr)}{\bigl((2k-\sigma)^2+t^2\bigr)\bigl(4k^2+t^2\bigr)}.
\end{equation}
To check it, $\operatorname{Re}(s/(2k-s))=(2k\sigma-\sigma^2-t^2)/D_1$ with $D_1=(2k-\sigma)^2+t^2$; with $D_2=4k^2+t^2$ the numerator of the left side over $D_1D_2$ is $(2k\sigma-\sigma^2)D_2-t^2D_2+t^2D_1=\sigma(2k-\sigma)(4k^2+t^2)-\sigma t^2(4k-\sigma)=2k\sigma(2k(2k-\sigma)-t^2)$, since $D_1-D_2=-4k\sigma+\sigma^2$. For $0<\sigma<1$, $k\ge1$ and $t^2\le2$ we have $2k(2k-\sigma)>2k(2k-1)\ge2\ge t^2$, so the right side of \CGeqref{eq:caseIid} is positive and
\begin{equation}\label{eq:caseIbound}
\operatorname{Re}\frac{s}{k(2k-s)}>-\frac{t^2}{k(4k^2+t^2)}\ge-\frac{t^2}{4k^3}\ge-\frac1{2k^3}.
\end{equation}
Summing over $k$ gives \CGeqref{eq:caseI}, since $\log\pi-\zeta(3)/2=1.14473\ldots-0.60103\ldots=0.54370\ldots$.

Case $|t|\ge\sqrt2$. Use \CGeqref{eq:B1}. The terms $\log(2\pi)$ and $-\zeta(2k)/k$ are real, $\operatorname{Im}(1/s)=-t/(\sigma^2+t^2)$ and $\operatorname{Im}(2/(2k-s))=2t/((2k-\sigma)^2+t^2)$, which gives the equality in \CGeqref{eq:caseII}. All terms of the sum are positive, and the term $k=1$ alone exceeds $1/(\sigma^2+t^2)$, because $2\zeta(2)=\pi^2/3$ and
\begin{equation}\label{eq:caseIIbound}
\frac{\pi^2}3(\sigma^2+t^2)-(2-\sigma)^2-t^2\ge\Bigl(\frac{\pi^2}3-1\Bigr)t^2-(2-\sigma)^2\ge2\Bigl(\frac{\pi^2}3-1\Bigr)-4=\frac{2\pi^2}3-6=0.5797\ldots>0
\end{equation}
for $t^2\ge2$ and $0<\sigma<1$. Hence $\operatorname{Im}B(s)\neq0$.
\end{proof}

\begin{remark}\label{rem:Anum}\VCGreviewed{1}
The identities \CGeqref{eq:caseIid} and \CGeqref{eq:caseII}, and the expansion behind \CGeqref{eq:caseIIbound}, are also checked symbolically by \texttt{route\_constants.py}. At the first zero $\rho_1=\tfrac12+i\,14.1347\ldots$ of $\zeta$ one finds numerically $A(\rho_1)=0.10741008+0.10085232\,i$, $|A(\rho_1)|=0.147337$.
\end{remark}

\section{The M\"obius log-sine function \texorpdfstring{$H$}{H}}\label{sec:H}

For $x\ge1$, $x\notin\mathbb Z$, let
\begin{equation}\label{eq:Hdef}
H(x)=\sum_{d\le x}\mu(d)\,F\Bigl(\frac xd\Bigr).
\end{equation}
Each $x/d$ is then not an integer, so every term is finite.

\begin{proposition}[Local structure]\label{prop:Hlocal}\VCGreviewed{1}
\textup{(i)} For $X\ge1$,
\begin{equation}\label{eq:Hcrude}
\int_1^X|H(x)|^2\,dx\le\frac{\pi^2}{12}X^3,\qquad\int_1^X|H(x)|\,dx\le\frac{\pi}{2\sqrt3}X^2.
\end{equation}
\textup{(ii)} Let $n\ge2$ be an integer. As $x\to n^-$, $H(x)+\mu(n)\log(n-x)$ has a finite limit. As $x\to n^+$, $H(x)$ has the finite limit
\begin{equation}\label{eq:Hplus}
H(n^+)=\Lambda(n)+\sum_{d\le n,\ d\nmid n}\mu(d)\,F\Bigl(\frac nd\Bigr).
\end{equation}
\end{proposition}

\begin{proof}
(i) For $1\le d\le X$, by periodicity, $\int_d^XF(x/d)^2\,dx=d\int_1^{X/d}F^2\le d\,\lceil X/d-1\rceil\int_0^1F^2\le X\pi^2/12$, because $[1,X/d]$ is covered by $\lceil X/d-1\rceil\le X/d$ intervals of length $1$. By Cauchy--Schwarz over the at most $\lfloor x\rfloor\le X$ terms, $|H(x)|^2\le X\sum_{d\le x}F(x/d)^2$. Integrating over $[1,X]$ and summing the bound above over the at most $X$ values of $d$ gives $\int_1^X|H|^2\le X\cdot X\cdot X\pi^2/12$, the first bound. For the second, $\int_d^X|F(x/d)|\,dx\le d\lceil X/d-1\rceil\int_0^1|F|\le X(\pi^2/12)^{1/2}$, and summing over $d\le X$ gives the second bound, since $(\pi^2/12)^{1/2}=\pi/(2\sqrt3)$.

(ii) Let $x\in(n-1,n)$. The terms of \CGeqref{eq:Hdef} are those with $d\le n-1$. A term with $d\nmid n$ is continuous at $x=n$. For $d\mid n$, $d<n$, we have $F(x/d)=\log(2\pi(n-x)/d)+o(1)$ as $x\to n^-$, because $|\sin(\pi x/d)|=|\sin(\pi(n-x)/d)|$. Since $\sum_{d\mid n,\,d<n}\mu(d)=-\mu(n)$ for $n\ge2$, the singular part is $-\mu(n)\log(n-x)$. For $x\in(n,n+1)$ the term $d=n$ is present too, the singular parts of the $d\mid n$ terms add up to $\sum_{d\mid n}\mu(d)\log(2\pi(x-n))=0$, and the remaining constants are $-\sum_{d\mid n}\mu(d)\log d=\Lambda(n)$.
\end{proof}

\begin{proposition}[Two identities]\label{prop:Hid}\VCGreviewed{1}
Let $x>1$, $x\notin\mathbb Z$, and let $\log(1-e^{2\pi iv})$ denote the principal branch. Then
\begin{equation}\label{eq:Hcomplex}
\sum_{d\le x}\mu(d)\log\bigl(1-e^{2\pi ix/d}\bigr)=H(x)+i\pi\Bigl(x\,m(x)-1-\frac{M(x)}2\Bigr).
\end{equation}
With $\operatorname{sinc}v=\sin(\pi v)/(\pi v)$,
\begin{equation}\label{eq:Hsinc}
H(x)=\log|x^2-1|+\log(2\pi)M(x)+\int_1^x\frac{M(u)}u\,du-\sum_{d>x}\mu(d)\log\operatorname{sinc}\frac xd,
\end{equation}
where the last series converges absolutely.
\end{proposition}

\begin{proof}
For $v>0$, $v\notin\mathbb Z$, $\log(1-e^{2\pi iv})=F(v)+i\pi(\{v\}-\tfrac12)$. Summing with weights $\mu(d)$ at $v=x/d$, and using $\{x/d\}=x/d-\lfloor x/d\rfloor$ and $\sum_{d\le x}\mu(d)\lfloor x/d\rfloor=1$ for $x\ge1$, gives \CGeqref{eq:Hcomplex}.

For \CGeqref{eq:Hsinc}, Euler's product gives $\log|\operatorname{sinc}v|=\sum_{\nu\ge1}\log|1-v^2/\nu^2|$ for $v\notin\mathbb Z$. For fixed $x$, $|\log|1-x^2/n^2||\le2x^2/n^2$ when $n\ge2x$, and the pairs $(d,\nu)$ with $d\nu<2x$ are finitely many, so $\sum_{d,\nu}|\log|1-x^2/(d\nu)^2||\le\sum_n\tau(n)|\log|1-x^2/n^2||<\infty$, where $\tau(n)$ is the number of divisors of $n$. Grouping by $n=d\nu$ and using $\sum_{d\mid n}\mu(d)=[n=1]$,
\begin{equation}\label{eq:sincid}
\sum_{d\ge1}\mu(d)\log\Bigl|\operatorname{sinc}\frac xd\Bigr|=\log|1-x^2|.
\end{equation}
For $d\le x$, $F(x/d)=\log(2\pi x/d)+\log|\operatorname{sinc}(x/d)|$, and $\sum_{d\le x}\mu(d)\log(x/d)=\int_1^xM(u)\,du/u$ by partial summation. For $d>x$ we have $0<x/d<1$, so $\operatorname{sinc}(x/d)>0$. Inserting \CGeqref{eq:sincid} gives \CGeqref{eq:Hsinc}.
\end{proof}

\begin{proposition}[Mellin transform of $H$]\label{prop:HMellin}\VCGreviewed{1}
For $\operatorname{Re}s>1$ the integral $\widehat H(s)=\int_1^\infty H(x)x^{-s-1}\,dx$ converges absolutely and
\begin{equation}\label{eq:HMellin}
\widehat H(s)=\frac{\widehat F(s)}{\zeta(s)},\qquad\widehat H(s)=\frac\pi s\cot\Bigl(\frac{\pi s}2\Bigr)+\frac{A(s)}{\zeta(s)}\quad(1<\operatorname{Re}s<2).
\end{equation}
\end{proposition}

\begin{proof}
By \CGeqref{eq:Hcrude} and dyadic decomposition, $\int_1^\infty|H|x^{-\sigma-1}<\infty$ for $\sigma>1$. Moreover, by the substitution $x=du$, the sum over $d$ of $\int_d^\infty|F(x/d)|x^{-\sigma-1}dx$ equals $\zeta(\sigma)\int_1^\infty|F(u)|u^{-\sigma-1}du<\infty$, so Fubini's theorem applies, and $\int_d^\infty F(x/d)x^{-s-1}\,dx=d^{-s}\widehat F(s)$ gives $\widehat H=\widehat F\sum_d\mu(d)d^{-s}=\widehat F/\zeta$. The second form is \CGeqref{eq:Fhat}; it is stated for $1<\operatorname{Re}s<2$, where $A$ is defined by \CGeqref{eq:Adef}, and only this range is used below.
\end{proof}

\begin{theorem}[$L^2$ criterion]\label{thm:L2}\VCGreviewed{1}
The following are equivalent.
\textup{(a)} RH.
\textup{(b)} $\int_1^\infty|H(x)|^2x^{-2\sigma-1}\,dx<\infty$ for every $\sigma>\tfrac12$.
\textup{(c)} For every $\varepsilon>0$ there is $K_\varepsilon$ with $\int_1^XH(x)^2\,dx\le K_\varepsilon X^{2+\varepsilon}$ for all $X\ge1$.
\end{theorem}

\begin{proof}
(b)$\Rightarrow$(c). With $\sigma=\tfrac12+\varepsilon/2$, $\int_1^XH^2\le X^{2\sigma+1}\int_1^X|H|^2x^{-2\sigma-1}\le K_\varepsilon X^{2+\varepsilon}$.

(c)$\Rightarrow$(a). Let $\sigma>\tfrac12$ and $\varepsilon=\sigma-\tfrac12$. By Cauchy--Schwarz on $[2^q,2^{q+1}]$, $q\ge0$ an integer,
\begin{equation}\label{eq:dyadic}
\int_{2^q}^{2^{q+1}}|H|x^{-\sigma-1}\,dx\le2^{-q(\sigma+1)}2^{q/2}\bigl(K_\varepsilon2^{(q+1)(2+\varepsilon)}\bigr)^{1/2}\ll_\sigma2^{-q(\sigma-1/2-\varepsilon/2)}=2^{-q\varepsilon/2},
\end{equation}
which is summable. Hence $\widehat H(s)$ converges absolutely and locally uniformly on $\operatorname{Re}s>\tfrac12$ and is holomorphic there. By \CGeqref{eq:HMellin} and uniqueness of continuation, $\frac\pi s\cot(\frac{\pi s}2)+A(s)/\zeta(s)$ extends holomorphically to $\tfrac12<\operatorname{Re}s<2$. The first term is holomorphic on $0<\operatorname{Re}s<2$. So $A/\zeta$ is holomorphic on $\tfrac12<\operatorname{Re}s<1$. If $\zeta(\rho)=0$ with $\tfrac12<\operatorname{Re}\rho<1$, then $A(\rho)\neq0$ by Theorem~\CGref{thm:Anonzero}, and $A/\zeta$ has a pole at $\rho$, a contradiction. Since $\zeta(s)\neq0$ for $\operatorname{Re}s>1$ (Euler product) and on $\operatorname{Re}s=1$ \cite[\S3.1, p.~45]{TitchmarshHB1986}, RH holds.

(a)$\Rightarrow$(b). Assume RH and fix $\sigma_0\in(\tfrac12,1)$. By \CGeqref{eq:HMellin}, $\widehat H$ continues holomorphically to $\tfrac12<\operatorname{Re}s<2$: the cotangent term has no poles there, $\zeta$ has no zeros there, and at $s=1$ the pole of $\zeta$ makes $A/\zeta$ vanish. Let $s=\sigma+it$ with $\sigma_0\le\sigma\le\tfrac32$ and $|t|\ge2$. Then $|\cot(\pi s/2)|\le\coth(\pi|t|/2)\le\coth\pi$, since $|\cot z|^2=(\cos^2\operatorname{Re}z+\sinh^2\operatorname{Im}z)/(\sin^2\operatorname{Re}z+\sinh^2\operatorname{Im}z)$. Next, $|A(s)|\ll\log|t|/|t|$: indeed $|2k-s|\ge|t|$ for all $k$, $|2k-s|\ge2k-\tfrac32\ge k/2$ for $k\ge1$, so $\sum_{k\le|t|}\zeta(2)/(k|t|)+\sum_{k>|t|}2\zeta(2)/k^2\ll\log|t|/|t|$. Finally $1/\zeta(s)\ll_\varepsilon|t|^\varepsilon$ uniformly for $\sigma_0\le\sigma\le1$ by \cite[Theorem~14.2, (14.2.1), p.~336, and (14.2.6), p.~337]{TitchmarshHB1986}, and $1/\zeta(s)\ll\log^7|t|$ for $1\le\sigma\le\frac32$ by \cite[(3.6.5), p.~51]{TitchmarshHB1986}. Hence
\begin{equation}\label{eq:Hhatbound}
|\widehat H(\sigma+it)|\ll_\varepsilon|t|^{-1+\varepsilon}\qquad(\sigma_0\le\sigma\le\tfrac32,\ |t|\ge2),
\end{equation}
and $\widehat H$ is bounded on the compact rectangle $[\sigma_0,\frac32]\times[-2,2]$. So $\sup_{\sigma_0<\sigma\le3/2}\int|\widehat H(\sigma+it)|^2dt<\infty$. For $\sigma\ge\frac32$ the function $\Upsilon_\sigma(v)=H(e^v)e^{-\sigma v}$ lies in $L^1\cap L^2(0,\infty)$ by \CGeqref{eq:Hcrude}, $\widehat H(\sigma+it)=\int_0^\infty\Upsilon_\sigma(v)e^{-itv}dv$, and Plancherel's theorem \cite[Theorem~9.13, p.~186]{Rudin1987} gives $\int|\widehat H(\sigma+it)|^2dt=2\pi\int_0^\infty|\Upsilon_\sigma|^2\le2\pi\int_0^\infty|\Upsilon_{3/2}|^2$. Thus
\begin{equation}\label{eq:H2sup}
\sup_{\sigma>\sigma_0}\frac1{2\pi}\int_{-\infty}^\infty|\widehat H(\sigma+it)|^2\,dt<\infty.
\end{equation}
Put $\mathcal P(z)=\widehat H(\sigma_0-iz)$ for $\operatorname{Im}z>0$, so that $\mathcal P(\xi+i\upsilon)=\widehat H(\sigma_0+\upsilon-i\xi)$. By \CGeqref{eq:H2sup} and the Paley--Wiener theorem \cite[Theorem~19.2, pp.~372--373]{Rudin1987} there is $\Psi\in L^2(0,\infty)$ with $\mathcal P(z)=\int_0^\infty\Psi(v)e^{ivz}\,dv$ for $\operatorname{Im}z>0$. For $\upsilon>\frac32-\sigma_0$ also $\mathcal P(\xi+i\upsilon)=\int_0^\infty H(e^v)e^{-\sigma_0v}e^{-\upsilon v}e^{i\xi v}\,dv$. Both $\Psi(v)e^{-\upsilon v}$ and $H(e^v)e^{-(\sigma_0+\upsilon)v}$ lie in $L^1(0,\infty)$ and have the same Fourier transform, so they agree almost everywhere \cite[Theorem~9.12, p.~185]{Rudin1987}. Hence $\int_0^\infty|H(e^v)|^2e^{-2\sigma_0v}dv=\int_0^\infty|\Psi|^2<\infty$, which is $\int_1^\infty|H(x)|^2x^{-2\sigma_0-1}dx<\infty$. Since $\sigma_0\in(\frac12,1)$ was arbitrary and the integral decreases in $\sigma$, (b) holds.
\end{proof}

\begin{remark}\label{rem:L2}\VCGreviewed{1}
A pointwise bound $H(x)\ll x^{1/2+\varepsilon}$ fails: by Proposition~\CGref{prop:Hlocal}(ii), $H(x)+\mu(n)\log(n-x)$ has a finite limit as $x\to n^-$, so $H$ is unbounded near every squarefree $n\ge2$. This is why the pointwise criterion is stated for the primitive $W$ in the next section.
\end{remark}

\section{The Clausen sum \texorpdfstring{$W$}{W}}\label{sec:W}

For $X\ge1$ let
\begin{equation}\label{eq:Wdef}
W(X)=\sum_{d\le X}\mu(d)\,d\,\mathrm{Cl}_2\Bigl(\frac{2\pi X}d\Bigr).
\end{equation}
The term $d$ vanishes at $X=d$ because $\mathrm{Cl}_2(2\pi)=0$, so $W$ is continuous, $W(1)=0$, and by \CGeqref{eq:Clmax}, $|W(X)|\le\mathrm{Cl}_2(\pi/3)X(X+1)/2$.

\begin{proposition}\label{prop:Wbasic}\VCGreviewed{1}
\textup{(i)} $W(X)=-2\pi\int_1^XH(x)\,dx$ for $X\ge1$.
\textup{(ii)} With $\mathcal K(v)=\mathrm{Cl}_2(2\pi v)+2\pi v\,F(v)$ for $v\notin\mathbb Z$,
\begin{equation}\label{eq:WMK}
W(X)=-\int_1^XM(y)\,\mathcal K\Bigl(\frac Xy\Bigr)\,dy.
\end{equation}
\textup{(iii)} For $\operatorname{Re}s>2$ the integral $\widehat W(s)=\int_1^\infty W(X)X^{-s-1}dX$ converges absolutely and
\begin{equation}\label{eq:WMellin}
\widehat W(s)=-\frac{2\pi}s\,\widehat H(s-1)=-\frac{2\pi}s\,\frac{\widehat F(s-1)}{\zeta(s-1)}.
\end{equation}
The right-hand side is meromorphic on $\operatorname{Re}s>0$ and holomorphic in a neighbourhood of the half-line $(1,\infty)$.
\end{proposition}

\begin{proof}
(i) For $X\notin\mathbb Z$, $\frac{d}{dX}\bigl(d\,\mathrm{Cl}_2(2\pi X/d)\bigr)=-2\pi F(X/d)$ by \CGeqref{eq:Cl2}. On each open interval $(n,n+1)$, $n\ge1$ an integer, every $x/d$ is a non-integer, so $W$ is continuously differentiable there with $W'=-2\pi H$. Since $W$ is continuous on $[n,n+1]$ and $H$ is integrable there by \CGeqref{eq:Hcrude}, $W(n+1)-W(n)=-2\pi\int_n^{n+1}H$, and also for intermediate endpoints; with $W(1)=0$ this gives (i).

(ii) Write $d\,\mathrm{Cl}_2(2\pi X/d)=X\mathcal Q(X/d)$ with $\mathcal Q(v)=\mathrm{Cl}_2(2\pi v)/v$. Partial summation, with $M(1^-)=0$ and $\mathcal Q(1)=0$, gives $\sum_{d\le X}\mu(d)\mathcal Q(X/d)=\int_1^XM(y)\,(X/y^2)\mathcal Q'(X/y)\,dy$. Since $v^2\mathcal Q'(v)=2\pi v\,\mathrm{Cl}_2'(2\pi v)-\mathrm{Cl}_2(2\pi v)=-\mathcal K(v)$, (ii) follows.

(iii) Absolute convergence follows from $|W(X)|\ll X^2$. Integrating by parts with (i), the boundary terms vanish for $\operatorname{Re}s>2$, and $\widehat W(s)=\frac1s\int_1^\infty W'(X)X^{-s}dX=-\frac{2\pi}s\widehat H(s-1)$. By Theorem~\CGref{thm:mellin}, $\widehat F$ is holomorphic on $\operatorname{Re}s>-1$, so $\widehat F(s-1)/\zeta(s-1)$ is meromorphic on $\operatorname{Re}s>0$. On the real half-line $s-1>0$, $\zeta(s-1)\neq0$: for $0<\sigma<1$, $(1-2^{1-\sigma})\zeta(\sigma)=\sum_n(-1)^{n-1}n^{-\sigma}>0$, and $\zeta(\sigma)>0$ for $\sigma>1$; at $s-1=1$, $1/\zeta$ is holomorphic with a zero. Since the zeros of $\zeta$ are isolated, the right-hand side of \CGeqref{eq:WMellin} is holomorphic near $(1,\infty)$.
\end{proof}

\begin{lemma}[Landau]\label{lem:landau}\VCGreviewed{1}
Let $f$ be real, locally bounded and measurable on $[1,\infty)$, with $f(X)\ge0$ for $X\ge X_0$. Suppose $g(s)=\int_1^\infty f(X)X^{-s-1}dX$ converges absolutely for $\operatorname{Re}s>\sigma_1$, and that $g$ extends holomorphically to a neighbourhood of the real segment $[\sigma_2,\sigma_1]$, where $\sigma_2<\sigma_1$. Then the integral converges absolutely for $\operatorname{Re}s>\sigma_2$.
\end{lemma}

\begin{proof}
Let $\sigma_c$ be the infimum of the real $\sigma$ with $\int_1^\infty|f|X^{-\sigma-1}<\infty$; then $\sigma_c\le\sigma_1$, and $g$ is holomorphic on $\operatorname{Re}s>\sigma_c$. Suppose $\sigma_c>\sigma_2$. Then $g$ is holomorphic on a neighbourhood $U$ of $[\sigma_c,\infty)$, the union of the half-plane and the given neighbourhood. Choose $h_0>0$ so small that the closed disc of radius $2h_0$ about $\sigma_3=\sigma_c+h_0$ lies in $U$. Splitting off $[1,X_0]$, where $f$ is bounded, we may assume $f\ge0$ on $[1,\infty)$. The extension agrees with the integral on the part of the disc in $\operatorname{Re}s>\sigma_c$, because that part is convex and contains real points of $(\sigma_c,\infty)$, where both are equal by uniqueness of continuation. For $k\ge0$, $g^{(k)}(\sigma_3)=\int_1^\infty f(X)(-\log X)^kX^{-\sigma_3-1}dX$ by differentiation under the integral. Taylor's theorem at $\sigma_3$ gives, for $h_0<h<2h_0$,
\begin{equation}\label{eq:landau}
g(\sigma_3-h)=\sum_{k\ge0}\frac{h^k}{k!}\int_1^\infty f(X)(\log X)^kX^{-\sigma_3-1}\,dX=\int_1^\infty f(X)X^{-(\sigma_3-h)-1}\,dX,
\end{equation}
where Tonelli's theorem applies because every term is nonnegative ($X\ge1$). The left side is finite, so the integral converges at the real point $\sigma_3-h<\sigma_c$, a contradiction. Hence $\sigma_c\le\sigma_2$.
\end{proof}

\begin{theorem}[Pointwise and one-sided criteria]\label{thm:Wcrit}\VCGreviewed{1}
The following are equivalent.
\textup{(a)} RH.
\textup{(b)} For every $\varepsilon>0$ there is $K_\varepsilon$ with $|W(X)|\le K_\varepsilon X^{3/2+\varepsilon}$ for $X\ge1$.
\textup{(c)} For every $\varepsilon>0$ there is $K_\varepsilon$ with $W(X)\le K_\varepsilon X^{3/2+\varepsilon}$ for $X\ge1$.
\textup{(d)} For every $\varepsilon>0$ there is $K_\varepsilon$ with $W(X)\ge-K_\varepsilon X^{3/2+\varepsilon}$ for $X\ge1$.
\end{theorem}

\begin{proof}
(a)$\Rightarrow$(b). Under RH, $M(y)\ll_\varepsilon y^{1/2+\varepsilon}$ \cite[Theorem~14.25(C), p.~370]{TitchmarshHB1986}. In \CGeqref{eq:WMK} substitute $y=X/v$:
\begin{equation}\label{eq:WMKbound}
|W(X)|\le\int_1^X\Bigl|M\Bigl(\frac Xv\Bigr)\Bigr||\mathcal K(v)|\frac X{v^2}\,dv\ll_\varepsilon X^{3/2+\varepsilon}\int_1^\infty v^{-5/2}|\mathcal K(v)|\,dv.
\end{equation}
Since $|\mathcal K(v)|\le1.015+2\pi v|F(v)|$ and $\int_n^{n+1}|F|=\int_0^1|F|$, the last integral is finite.

(b)$\Rightarrow$(c) and (b)$\Rightarrow$(d) are trivial.

(c)$\Rightarrow$(a). Fix $\varepsilon>0$ and put $\beta=\frac32+\varepsilon$ and $f(X)=K_\varepsilon X^\beta-W(X)\ge0$. Its Mellin integral $g(s)=K_\varepsilon/(s-\beta)-\widehat W(s)$ converges absolutely for $\operatorname{Re}s>2$, and by Proposition~\CGref{prop:Wbasic}(iii) it extends holomorphically to a neighbourhood of every real segment $[\sigma_2,2]$ with $\beta<\sigma_2$. By Lemma~\CGref{lem:landau} the integral of $f$ converges absolutely for $\operatorname{Re}s>\beta$, hence so does $\widehat W$, since $|W|\le f+K_\varepsilon X^\beta$. Thus $-\frac{2\pi}s\widehat F(s-1)/\zeta(s-1)$ is holomorphic on $\operatorname{Re}s>\beta$, i.e.\ $\widehat F(w)/\zeta(w)=\frac\pi w\cot\frac{\pi w}2+A(w)/\zeta(w)$ is holomorphic on $\operatorname{Re}w>\frac12+\varepsilon$. The cotangent term is holomorphic on $0<\operatorname{Re}w<2$, so $A/\zeta$ is holomorphic on $\frac12+\varepsilon<\operatorname{Re}w<1$. A zero $\rho$ of $\zeta$ in that strip would be a pole of $A/\zeta$, because $A(\rho)\neq0$ by Theorem~\CGref{thm:Anonzero}. So there is none, and with $\zeta(s)\neq0$ on $\operatorname{Re}s\ge1$ \cite[\S3.1, p.~45]{TitchmarshHB1986} and $\varepsilon\to0$ this gives RH.

(d)$\Rightarrow$(a). The same, with $f(X)=W(X)+K_\varepsilon X^\beta$ and $g(s)=\widehat W(s)+K_\varepsilon/(s-\beta)$.
\end{proof}

\begin{theorem}[$\Omega_\pm$ bounds]\label{thm:omega}\VCGreviewed{1}
Let $\rho=\frac12+i\gamma$, $\gamma>0$, be a simple zero of $\zeta$, and
\begin{equation}\label{eq:rrho}
r_\rho=-\frac{2\pi}{1+\rho}\,\frac{A(\rho)}{\zeta'(\rho)}\neq0.
\end{equation}
Then $\limsup_{X\to\infty}W(X)X^{-3/2}\ge|r_\rho|$ and $\liminf_{X\to\infty}W(X)X^{-3/2}\le-|r_\rho|$. For the first zero, $\gamma_1=14.13472514\ldots$, which lies on the critical line and is simple (classical numerical facts, used here as inputs), numerically $|\zeta'(\rho_1)|=0.79316$ and $|r_{\rho_1}|=0.0821128$.
\end{theorem}

\begin{proof}
$r_\rho\neq0$ by Theorem~\CGref{thm:Anonzero}, and $r_\rho$ is the residue of $\widehat W$ at $s=1+\rho$ by \CGeqref{eq:WMellin}, since $\widehat F(\rho)=A(\rho)$ by \CGeqref{eq:Fhat}. Suppose there are $b<|r_\rho|$ and $X_0$ with $W(X)\le bX^{3/2}$ for $X\ge X_0$. Let $f(X)=bX^{3/2}-W(X)$ and $g(s)=b/(s-\frac32)-\widehat W(s)$. The function $f$ is continuous and $f\ge0$ on $[X_0,\infty)$; its Mellin integral $g$ converges absolutely for $\operatorname{Re}s>2$ because $|W(X)|\ll X^2$; and $g$ extends holomorphically to a neighbourhood of $[\sigma_2,2]$ for every $\sigma_2>\frac32$ by Proposition~\CGref{prop:Wbasic}(iii). Lemma~\CGref{lem:landau} gives absolute convergence of $\int_1^\infty fX^{-s-1}$ on $\operatorname{Re}s>\frac32$; hence $g$ is holomorphic there, and $g$ is meromorphic near $s_0=\frac32+i\gamma$ with $g(s)=-r_\rho/(s-s_0)+O(1)$. Since $f$ is real, $g(\bar s)=\overline{g(s)}$. For real $\sigma>\frac32$ and real $\phi_0$,
\begin{equation}\label{eq:fejer}
g(\sigma)+\tfrac12e^{i\phi_0}g(\sigma+i\gamma)+\tfrac12e^{-i\phi_0}g(\sigma-i\gamma)=\int_1^\infty f(X)X^{-\sigma-1}\bigl(1+\cos(\gamma\log X-\phi_0)\bigr)\,dX\ge-Q_0,
\end{equation}
with $Q_0=2\int_1^{X_0}|f|X^{-5/2}dX$, because the integrand is nonnegative on $[X_0,\infty)$. As $\sigma\to\frac32^+$, $\widehat W(\sigma)$ stays bounded by Proposition~\CGref{prop:Wbasic}(iii), so $g(\sigma)=b/(\sigma-\frac32)+O(1)$, while $g(\sigma\pm i\gamma)=-r_\rho^{\pm}/(\sigma-\frac32)+O(1)$ with $r_\rho^+=r_\rho$, $r_\rho^-=\overline{r_\rho}$. The left side of \CGeqref{eq:fejer} is therefore $(b-\operatorname{Re}(e^{i\phi_0}r_\rho))/(\sigma-\frac32)+O(1)$. Choosing $e^{i\phi_0}r_\rho=|r_\rho|$ it tends to $-\infty$, a contradiction. Hence $\limsup W/X^{3/2}\ge|r_\rho|$. For the $\liminf$ use $f=W+bX^{3/2}$, $g=\widehat W+b/(s-\frac32)$ and $e^{i\phi_0}r_\rho=-|r_\rho|$. The numerical values come from \texttt{route\_constants.py} (mpmath, $30$ digits); mpmath locates $\rho_1$ on the critical line by assumption, so the location of $\rho_1$ is an input and not a consequence of this computation; the computed value $\zeta'(\rho_1)\neq0$ agrees with the simplicity of $\rho_1$.
\end{proof}

\section{The \texorpdfstring{$d=2$}{d=2} transfer and the corner}\label{sec:corner}

\begin{lemma}\label{lem:vaughan}\VCGreviewed{1}
Let $U\ge1$ and $\mu_U(n)=\mu(n)$ for $n\le U$, $\mu_U(n)=0$ for $n>U$. For $1\le n\le U^2$,
\begin{equation}\label{eq:vaughan}
\mu(n)=2\mu_U(n)-(\mu_U*\mu_U*1)(n).
\end{equation}
\end{lemma}

\begin{proof}
Let $\mu^U=\mu-\mu_U$, and let $\mathbf 1_{\{1\}}$ be the unit of Dirichlet convolution, so that $\mu*1=\mathbf 1_{\{1\}}$. Then $\mu^U*1*\mu^U=\mu-2\mu_U+\mu_U*\mu_U*1$. A nonzero term of $(\mu^U*\mu^U)(n)$ needs $n=d_1d_2$ with $d_1,d_2>U$, so $n>U^2$; hence $(\mu^U*1*\mu^U)(n)=0$ for $n\le U^2$. This is the M\"obius form of the device behind Vaughan's identity for $\Lambda$ \cite[Ch.~24, pp.~138--139]{Davenport1980}.
\end{proof}

For $Y\ge1$ let
\begin{equation}\label{eq:Gdef}
G(Y)=\sum_{\ell\le Y}\ell\,\mathrm{Cl}_2\Bigl(\frac{2\pi Y}\ell\Bigr),\qquad E_G(Y)=G(Y)-aY^2,
\end{equation}
and for $\operatorname{Re}s>0$ let
\begin{equation}\label{eq:Thetadef}
\Theta(s)=\int_1^\infty\mathrm{Cl}_2(2\pi v)\,v^{-s-1}\,dv=-\frac{2\pi}s\,\widehat F(s-1).
\end{equation}
The second equality holds for $\operatorname{Re}s>1$ by parts, using $\mathrm{Cl}_2(2\pi)=0$ and \CGeqref{eq:Cl2}, and then for $\operatorname{Re}s>0$ by continuation; $\Theta(2)=-\pi\widehat F(1)=a$ by \CGeqref{eq:aconst}. $G$ is continuous for the same reason as $W$, $G(1)=0$, and $|G(Y)|\le\mathrm{Cl}_2(\pi/3)Y(Y+1)/2$.

\begin{lemma}\label{lem:EG}\VCGreviewed{1}
\textup{(i)} For $\operatorname{Re}s>2$, $\widehat G(s)=\int_1^\infty G(Y)Y^{-s-1}dY=\zeta(s-1)\Theta(s)$.
\textup{(ii)} For every $\delta\in(0,1)$ there is $c_\delta$ with $|E_G(Y)|\le c_\delta Y^{1+\delta}$ for $Y\ge1$.
\textup{(iii)} For $Y\notin\mathbb Z$, $E_G'(Y)=-2\pi R(Y)$, where $R(Y)=\sum_{\ell\le Y}F(Y/\ell)-(\log2-C)Y$.
\end{lemma}

\begin{proof}
(i) By Fubini's theorem, justified by Tonelli's theorem applied to the absolute values, and the substitution $Y=\ell v$, $\int_\ell^\infty\ell\,\mathrm{Cl}_2(2\pi Y/\ell)Y^{-s-1}dY=\ell^{1-s}\Theta(s)$, and $\sum_\ell\ell^{1-\operatorname{Re}s}<\infty$.

(ii) Put $w=s-1$, $s=\sigma+it$. By \CGeqref{eq:Thetadef} and \CGeqref{eq:Fhat}, for $0<\operatorname{Re}w<2$,
\begin{equation}\label{eq:Ghatform}
\widehat G(s)=-\frac{2\pi}s\,\zeta(w)\Bigl(\frac\pi w\cot\frac{\pi w}2\,\zeta(w)+A(w)\Bigr).
\end{equation}
For $|t|\ge2$ and $0<\operatorname{Re}w\le\frac32$ we have $|\frac\pi w\cot\frac{\pi w}2|\le\pi\coth(\pi)/|t|$, because $|w|\ge|t|$ and $|\cot z|\le\coth|\operatorname{Im}z|$; and $|A(w)|\ll\log|t|/|t|$, because $|2k-w|\ge|t|$ and $|2k-w|\ge2k-\frac32\ge k/2$ for $k\ge1$, so that the series in \CGeqref{eq:Adef} is at most $\sum_{k\le|t|}\zeta(2)/(k|t|)+\sum_{k>|t|}2\zeta(2)/k^2$. Hence
\begin{equation}\label{eq:Ghatbound}
|\widehat G(s)|\ll|t|^{-2}\bigl(|\zeta(w)|^2+|\zeta(w)|\log|t|\bigr)\qquad(1<\sigma\le\tfrac52,\ |t|\ge2).
\end{equation}
On $\sigma=\frac52$, $|\zeta(w)|\le\zeta(\frac32)$, so $\widehat G(\frac52+it)\in L^1(dt)$. The function $Y\mapsto G(Y)Y^{-5/2}$, extended by $0$ to $(0,1)$, is continuous, and in the variable $v=\log Y$ it is integrable with integrable Fourier transform; the inversion theorem \cite[Theorem~9.11, p.~185]{Rudin1987} gives
\begin{equation}\label{eq:Ginv}
G(Y)=\frac1{2\pi i}\int_{(5/2)}\widehat G(s)\,Y^s\,ds\qquad(Y\ge1).
\end{equation}
Move the line to $\sigma=1+\delta$. In $1+\delta\le\sigma\le\frac52$ the only pole of $\widehat G$ is $s=2$, from $\zeta(s-1)$, with residue $\Theta(2)Y^2=aY^2$; $\Theta$ is holomorphic on $\operatorname{Re}s>0$. On the horizontal segments $\sigma+iT$, $1+\delta\le\sigma\le\frac52$, Theorem~4.11 of \cite[p.~77]{TitchmarshHB1986}, applied with a partial sum of length $|T|$, gives $|\zeta(w)|\ll_\delta|T|^{1-\delta}$ uniformly for $\operatorname{Re}w\ge\delta$, so by \CGeqref{eq:Ghatbound} $|\widehat G(\sigma+iT)|\ll|T|^{-2\delta}\log|T|\to0$. On the line $\sigma=1+\delta$, $\operatorname{Re}w=\delta$ and the convexity bound $\zeta(\delta+it)\ll_\varepsilon|t|^{(1-\delta)/2+\varepsilon}$ \cite[(5.1.4), p.~96]{TitchmarshHB1986} with $\varepsilon=\delta/4$ gives $|\widehat G(1+\delta+it)|\ll|t|^{-1-\delta/2}$ for $|t|\ge2$; near $t=0$, $\widehat G$ is continuous on the line. Hence
\begin{equation}\label{eq:EGint}
E_G(Y)=\frac1{2\pi i}\int_{(1+\delta)}\widehat G(s)\,Y^s\,ds,\qquad|E_G(Y)|\le\frac{Y^{1+\delta}}{2\pi}\int_{-\infty}^\infty|\widehat G(1+\delta+it)|\,dt=c_\delta Y^{1+\delta}.
\end{equation}

(iii) On $(n,n+1)$, $n\ge1$ an integer, each term of \CGeqref{eq:Gdef} is differentiable and $\frac{d}{dY}\bigl(\ell\,\mathrm{Cl}_2(2\pi Y/\ell)\bigr)=-2\pi F(Y/\ell)$ by \CGeqref{eq:Cl2}, so $G'(Y)=-2\pi\sum_{\ell\le Y}F(Y/\ell)$; with $\frac{d}{dY}aY^2=2aY=2\pi(C-\log2)Y$ this gives $E_G'=-2\pi R$.
\end{proof}

For $X\ge1$ define
\begin{equation}\label{eq:Wflat}
W^\flat(X)=\sum_{n\le\sqrt X}\mu(n)\,n\,\mathrm{Cl}_2\Bigl(\frac{2\pi X}n\Bigr),
\end{equation}
\begin{equation}\label{eq:Rest}
\mathcal R(X)=\sum_{\substack{d,j\le\sqrt X\\ dj\le\sqrt X}}\mu(d)\mu(j)\,dj\,E_G\Bigl(\frac X{dj}\Bigr),\qquad\mathcal C(X)=\sum_{\substack{d,j\le\sqrt X\\ dj>\sqrt X}}\mu(d)\mu(j)\,dj\,E_G\Bigl(\frac X{dj}\Bigr).
\end{equation}
In $\mathcal C(X)$ the argument $Y=X/(dj)$ lies in $[1,\sqrt X)$; in $\mathcal R(X)$ it lies in $[\sqrt X,X]$.

\begin{proposition}[Exact decomposition]\label{prop:decomp}\VCGreviewed{1}
For $X\ge1$,
\begin{equation}\label{eq:decomp}
W(X)=2W^\flat(X)-aX^2m(\sqrt X)^2-\mathcal R(X)-\mathcal C(X).
\end{equation}
Moreover $|2W^\flat(X)|\le\mathrm{Cl}_2(\pi/3)(X+\sqrt X)$, and for every $\delta\in(0,1)$,
\begin{equation}\label{eq:Restbound}
|\mathcal R(X)|\le c_\delta X^{3/2+\delta/2}\bigl(1+\log\sqrt X\bigr)^2.
\end{equation}
\end{proposition}

\begin{proof}
By Lemma~\CGref{lem:vaughan} with $U=\sqrt X$, for $n\le X$, $\mu(n)=2\mu_U(n)-\sum_{dj\ell=n,\ d,j\le U}\mu(d)\mu(j)$. Hence
\begin{equation}\label{eq:decomp1}
\begin{split}
W(X)&=2W^\flat(X)-\sum_{d,j\le\sqrt X}\mu(d)\mu(j)\sum_{\ell\le X/(dj)}dj\ell\,\mathrm{Cl}_2\Bigl(\frac{2\pi X}{dj\ell}\Bigr)\\
&=2W^\flat(X)-\sum_{d,j\le\sqrt X}\mu(d)\mu(j)\,dj\,G\Bigl(\frac X{dj}\Bigr),
\end{split}
\end{equation}
where $dj\le X$, so $X/(dj)\ge1$. Inserting $G(Y)=aY^2+E_G(Y)$, the main terms give $aX^2$ times the double sum of $\mu(d)\mu(j)/(dj)$ over $d,j\le\sqrt X$, which is $aX^2m(\sqrt X)^2$, and the error terms split into $\mathcal R+\mathcal C$. The bound on $W^\flat$ follows from \CGeqref{eq:Clmax} and $2\sum_{n\le\sqrt X}n\le\sqrt X(\sqrt X+1)$. For \CGeqref{eq:Restbound}, the number of pairs with $dj=n$ is at most $\tau(n)$, so by Lemma~\CGref{lem:EG}(ii)
\begin{equation}\label{eq:Restpf}
|\mathcal R(X)|\le c_\delta\sum_{n\le\sqrt X}\tau(n)\,n\Bigl(\frac Xn\Bigr)^{1+\delta}\le c_\delta X^{1+\delta}(\sqrt X)^{1-\delta}\sum_{n\le\sqrt X}\frac{\tau(n)}n=c_\delta X^{3/2+\delta/2}\sum_{n\le\sqrt X}\frac{\tau(n)}n,
\end{equation}
and $\sum_{n\le y}\tau(n)/n\le(\sum_{n\le y}1/n)^2\le(1+\log y)^2$.
\end{proof}

\begin{theorem}[Reduction of RH to the corner]\label{thm:corner}\VCGreviewed{1}
The following are equivalent.
\textup{(a)} RH.
\textup{(b)} For every $\varepsilon>0$ there is $K_\varepsilon$ with $\mathcal C(X)\ge-K_\varepsilon X^{3/2+\varepsilon}$ for all $X\ge1$.
\textup{(c)} For every $\varepsilon>0$ there is $K_\varepsilon$ with $|\mathcal C(X)|\le K_\varepsilon X^{3/2+\varepsilon}$ for all $X\ge1$.
\end{theorem}

\begin{proof}
(b)$\Rightarrow$(a). By \CGeqref{eq:decomp}, $a>0$ and Proposition~\CGref{prop:decomp},
\begin{equation}\label{eq:cornerpf}
\begin{split}
W(X)&=2W^\flat(X)-aX^2m(\sqrt X)^2-\mathcal R(X)-\mathcal C(X)\\
&\le\mathrm{Cl}_2(\tfrac\pi3)(X+\sqrt X)+c_\delta X^{3/2+\delta/2}(1+\log\sqrt X)^2+K_\varepsilon X^{3/2+\varepsilon}.
\end{split}
\end{equation}
With $\delta=\min(\varepsilon,\frac12)$ the right side is at most $K'_\varepsilon X^{3/2+\varepsilon}$ for $X\ge1$, so condition (c) of Theorem~\CGref{thm:Wcrit} holds for every $\varepsilon>0$, and RH follows. No information on $m(\sqrt X)$ is used: the term $-aX^2m(\sqrt X)^2$ is $\le0$ and is dropped.

(a)$\Rightarrow$(c). Assume RH. By Theorem~\CGref{thm:Wcrit}, $W\ll X^{3/2+\varepsilon}$. By partial summation, for $\operatorname{Re}s>1$, $\sum_n\mu(n)n^{-s}=s\int_1^\infty M(v)v^{-s-1}dv$; under RH, $M(v)\ll v^{1/2+\varepsilon}$ \cite[Theorem~14.25(C), p.~370]{TitchmarshHB1986}, so the integral converges absolutely and uniformly for $s\in[1,2]$, and letting $s\to1^+$ gives $\int_1^\infty M(v)v^{-2}dv=\lim_{s\to1^+}1/(s\zeta(s))=0$. Then
\begin{equation}\label{eq:mbound}
m(y)=\frac{M(y)}y+\int_1^y\frac{M(v)}{v^2}\,dv=\frac{M(y)}y-\int_y^\infty\frac{M(v)}{v^2}\,dv\ll_\varepsilon y^{-1/2+\varepsilon}\qquad(y\ge1),
\end{equation}
so $aX^2m(\sqrt X)^2\ll X^{3/2+\varepsilon}$. Solving \CGeqref{eq:decomp} for $\mathcal C$ and using Proposition~\CGref{prop:decomp} gives (c).

(c)$\Rightarrow$(b) is trivial.
\end{proof}

\begin{remark}\label{rem:corner}\VCGreviewed{1}
$\mathcal C(X)$ depends on $\mu$ only through its values on $[1,\sqrt X]$, and on $E_G$ only on $[1,\sqrt X)$. The function $E_G$ involves no M\"obius values. By Lemma~\CGref{lem:EG}(iii), the derivative of $E_G$ is $-2\pi R$, and at each integer $n\ge2$ the function $R$ has logarithmic singularities, $(\tau(n)-1)\log|Y-n|$ from the left and $\tau(n)\log|Y-n|$ from the right, coming from the terms $\ell\mid n$ of $\sum_\ell F(Y/\ell)$. Every attempt recorded in Section~\CGref{sec:attempts} stopped at the size of $\sum|$terms$|$ of $\mathcal C$, which grows like $X^2$ in the measured range (Section~\CGref{sec:attempts}).
\end{remark}

\section{Other exact identities}\label{sec:other}

\begin{proposition}[Recursion over large primes]\label{prop:buchstab}\VCGreviewed{1}
For $X\ge1$ let $W^{\mathrm{sm}}(X)$ be the sum \CGeqref{eq:Wdef} restricted to $n$ whose prime factors are all $\le\sqrt X$. Then
\begin{equation}\label{eq:buchstab}
W(X)=W^{\mathrm{sm}}(X)-\sum_{\sqrt X<p\le X}p\,W\Bigl(\frac Xp\Bigr).
\end{equation}
\end{proposition}

\begin{proof}
A squarefree $n\le X$ has at most one prime factor $p>\sqrt X$, and then $n=pb$ with $b\le X/p<\sqrt X<p$, so $p\nmid b$ and $\mu(pb)=-\mu(b)$. Conversely each pair $(p,b)$ with $\sqrt X<p\le X$, $b\le X/p$, gives such an $n$, and $X/p\ge1$. Hence the non-smooth part of \CGeqref{eq:Wdef} is $-\sum_pp\sum_{b\le X/p}\mu(b)\,b\,\mathrm{Cl}_2(2\pi(X/p)/b)=-\sum_ppW(X/p)$.
\end{proof}

\begin{proposition}[lcm form]\label{prop:lcm}\VCGreviewed{1}
Let $X\ge1$, let $\mathcal D$ be the set of squarefree $d\le X$, let $\Xi_X(n)=n\,\mathrm{Cl}_2(2\pi X/n)$ for $n\le X$ and $\Xi_X(n)=0$ for $n>X$, and let $\mathbf K$ be the matrix $(\Xi_X([d,d']))_{d,d'\in\mathcal D}$, where $[d,d']$ is the least common multiple. Let $\boldsymbol\mu=(\mu(d'))_{d'\in\mathcal D}$ and let $\mathbf e_1$ be the coordinate vector of the index $1\in\mathcal D$. Then $\mathbf K\boldsymbol\mu=W(X)\,\mathbf e_1$. In particular $W(X)=\sum_{d,d'\in\mathcal D}\mu(d)\mu(d')\Xi_X([d,d'])$, and $W(X)=1/(\mathbf K^{-1})_{11}$ whenever $\mathbf K$ is invertible.
\end{proposition}

\begin{proof}
Fix $\bar d\in\mathcal D$. Each squarefree $d'$ factors uniquely as $d'=d_1d_2$ with $d_1\mid\bar d$, $(d_2,\bar d)=1$, and then $[\bar d,d']=\bar dd_2$ and $\mu(d')=\mu(d_1)\mu(d_2)$. Every squarefree $d'$ with $\Xi_X([\bar d,d'])\neq0$ has $d'\le[\bar d,d']\le X$, so it lies in $\mathcal D$. Thus $(\mathbf K\boldsymbol\mu)_{\bar d}=\sum_{d_1\mid\bar d}\mu(d_1)\sum_{(d_2,\bar d)=1}\mu(d_2)\Xi_X(\bar dd_2)$. The inner sum does not depend on $d_1$, and $\sum_{d_1\mid\bar d}\mu(d_1)=0$ for $\bar d>1$. For $\bar d=1$ the sum is $\sum_{d'\in\mathcal D}\mu(d')\Xi_X(d')=W(X)$. Taking the inner product with $\boldsymbol\mu$ gives the quadratic form, because $\mu(1)=1$. If $\mathbf K$ is invertible, $\boldsymbol\mu=W(X)\mathbf K^{-1}\mathbf e_1$, and the coordinate of the index $1$ gives $1=W(X)(\mathbf K^{-1})_{11}$.
\end{proof}

Invertibility of $\mathbf K$ is not proved. Numerically, for $X=300.5$, $1000.5$, $2500.5$, the identity $\mathbf K\boldsymbol\mu=W(X)\mathbf e_1$ holds to $10^{-11}$, $\det\mathbf K/\det\mathbf K^{(1)}$ reproduces $W(X)$, and the condition numbers are $4.9\cdot10^3$, $5.8\cdot10^{14}$, $8.9\cdot10^{14}$; the inertia of $\mathbf K^{(1)}$ is close to balanced ($84/98$, $281/326$, $695/827$ positive/negative eigenvalues), and $\boldsymbol\mu$ has almost no weight on the eigenvectors with eigenvalues of large modulus (the matrix $\mathbf K^{(1)}$ is $\mathbf K$ without the row and column of the index $1$).

\begin{proposition}[$d=3$]\label{prop:d3}\VCGreviewed{1}
Let $X\ge1$, $U=X^{1/3}$, $G_2(Y)=\sum_{\ell_1\ell_2\le Y}\ell_1\ell_2\,\mathrm{Cl}_2(2\pi Y/(\ell_1\ell_2))$, and
\begin{equation}\label{eq:d3}
\begin{split}
&T_1=\sum_{n\le U}\mu(n)\,n\,\mathrm{Cl}_2\Bigl(\frac{2\pi X}n\Bigr),\qquad T_2=\sum_{d,j\le U}\mu(d)\mu(j)\,dj\,G\Bigl(\frac X{dj}\Bigr),\\
&T_3=\sum_{d,j,h\le U}\mu(d)\mu(j)\mu(h)\,djh\,G_2\Bigl(\frac X{djh}\Bigr).
\end{split}
\end{equation}
Then $W(X)=3T_1-3T_2+T_3$.
\end{proposition}

\begin{proof}
Let $\mu^U=\mu-\mu_U$. A nonzero term of $(\mu^U)^{*3}(n)$ needs $n=d_1d_2d_3$ with $d_1,d_2,d_3>U$, so $n>U^3$; hence $(\mu^U)^{*3}*1^{*2}$ vanishes on $[1,U^3]$, and its expansion is $\mu-3\mu_U+3\mu_U^{*2}*1-\mu_U^{*3}*1^{*2}$, using $\mu*1=\mathbf 1_{\{1\}}$. Multiply by $n\,\mathrm{Cl}_2(2\pi X/n)$ and sum over $n\le X$.
\end{proof}

The route of Theorem~\CGref{thm:corner} does not use Proposition~\CGref{prop:d3}. In particular it needs no bound on $m(X^{1/3})$ or on any other tail of $m$; such quantities appear only if $T_2$ and $T_3$ are split into main terms and errors, and that split is not part of the route.

\section{Two facts about methods}\label{sec:methods}

\begin{proposition}\label{prop:musq}\VCGreviewed{1}
Let $W_{\mu^2}(X)=\sum_{n\le X}\mu(n)^2\,n\,\mathrm{Cl}_2(2\pi X/n)$. For every $\delta\in(0,\frac12)$,
\begin{equation}\label{eq:musq}
W_{\mu^2}(X)=\frac6{\pi^2}\,a\,X^2+O_\delta\bigl(X^{3/2+\delta}\bigr),\qquad\frac6{\pi^2}a=0.0967299459812\ldots
\end{equation}
\end{proposition}

\begin{proof}
By Fubini's theorem, justified as in Lemma~\CGref{lem:EG}(i) by $|\mathrm{Cl}_2|\le\mathrm{Cl}_2(\pi/3)$ and $\sum_nn^{1-\operatorname{Re}s}<\infty$, and the substitution $X=nv$, $\int_n^\infty n\,\mathrm{Cl}_2(2\pi X/n)X^{-s-1}dX=n^{1-s}\Theta(s)$, and $\sum_n\mu(n)^2n^{-z}=\zeta(z)/\zeta(2z)$ for $\operatorname{Re}z>1$, so $\int_1^\infty W_{\mu^2}X^{-s-1}dX=\Theta(s)\zeta(s-1)/\zeta(2s-2)$ for $\operatorname{Re}s>2$. On $\operatorname{Re}s\ge\frac32+\delta$, $|1/\zeta(2s-2)|\le\zeta(1+2\delta)$. With this extra bounded factor, the bound \CGeqref{eq:Ghatbound} holds for the new transform on $\frac32+\delta\le\sigma\le\frac52$, $|t|\ge2$; it is integrable on $\sigma=\frac52$, so the inversion formula \CGeqref{eq:Ginv} holds for $W_{\mu^2}$ by \cite[Theorem~9.11, p.~185]{Rudin1987}, since $W_{\mu^2}$ is continuous and $O(X^2)$. Move the line from $\frac52$ to $\frac32+\delta$: the only pole is $s=2$, with residue $\Theta(2)X^2/\zeta(2)=6aX^2/\pi^2$; on the horizontal segments $|\zeta(w)|\ll|T|^{1/2}$ uniformly for $\operatorname{Re}w\ge\frac12$ \cite[Theorem~4.11, p.~77]{TitchmarshHB1986}, so the integrand is $\ll|T|^{-1}\log|T|$; on the line, $|\zeta(\frac12+\delta+it)|\ll|t|^{1/4}$ by \cite[(5.1.4), p.~96]{TitchmarshHB1986}, so the integrand is $\ll|t|^{-3/2}$.
\end{proof}

\begin{remark}[What a bound for $\mathcal C$ must use]\label{rem:barrier}\VCGreviewed{1}
For squarefree $d,d'$ one has $\mu(d)\mu(d')=\mu(dd'/(d,d')^2)$, and the same relation holds with $\mu$ replaced by $\mu^2$, or by $n\mapsto\mu(n)^2\prod_{p\mid n}\chi(p)$ for any $\chi$ from the primes to $\{\pm1\}$. An argument that uses $\mu$ only through this relation, $|\mu|\le1$ and the support on squarefree integers applies equally to $\mu^2$. By Proposition~\CGref{prop:musq}, $W_{\mu^2}(X)\sim(6/\pi^2)aX^2$, so such an argument cannot give $W(X)\ll X^{3/2+\varepsilon}$. For the corner inequality \CGeqref{eq:target} the analogous conclusion is not proved: the route from the corner to RH uses $\mu*1=\mathbf 1_{\{1\}}$ through Lemma~\CGref{lem:vaughan}, which fails for $\mu^2$, so a bound for the $\mu^2$-corner $\sum\mu(d)^2\mu(j)^2dj\,E_G(X/(dj))$ would not bound $W_{\mu^2}$. The conclusion would follow if the $\mu^2$-corner violated the one-sided bound; the measured values in Section~\CGref{sec:attempts} ($-0.00090X^2$, $-0.00121X^2$, $-0.00148X^2$ for the part with $X/(dj)\le64$) are negative and grow like $X^2$, but this is not proved.
\end{remark}

\begin{remark}[Localization, not proved]\label{rem:local}\VCGreviewed{1}
An inequality of the form $W(X)\le c_1X+c_\varepsilon X^{3/2+\kappa+\varepsilon}(1+J(\sqrt X))$, with $J(y)=\int_1^yM(v)^2v^{-2}dv+M(y)^2/y$ and $\kappa$ a subconvexity exponent for $\zeta$ on the critical line, was derived in exploratory work through a mean-value argument in $\log y$. Its proof has not been written out or checked, and nothing in this document uses it.
\end{remark}

\section{A Hilbert-space form: the M\"untz transform of the log-sine kernel}\label{sec:muntz}

Let
\begin{equation}\label{eq:kdef}
\mathfrak k(t)=F(1/t)=\log\Bigl|2\sin\frac\pi t\Bigr|\quad(0<t\le1,\ 1/t\notin\mathbb Z),\qquad\mathfrak k(t)=0\quad(t>1).
\end{equation}
For $\sigma>0$ put
\begin{equation}\label{eq:Nsigma}
\mathcal N_\sigma=\int_1^\infty F(v)^2v^{-2\sigma-1}\,dv\le\frac{\pi^2}{12}\sum_{n\ge1}n^{-2\sigma-1}<\infty,
\end{equation}
using $\int_n^{n+1}F^2=\pi^2/12$ (proof of Theorem~\CGref{thm:mellin}). The substitution $t=1/v$ gives $\int_0^\infty|\mathfrak k(t)|t^{\sigma-1}dt=\int_1^\infty|F(v)|v^{-\sigma-1}dv<\infty$ and $\int_0^\infty\mathfrak k(t)^2t^{2\sigma-1}dt=\mathcal N_\sigma$, and
\begin{equation}\label{eq:kmellin}
\int_0^\infty\mathfrak k(t)\,t^{s-1}\,dt=\int_1^\infty F(v)\,v^{-s-1}\,dv=\widehat F(s)\qquad(\operatorname{Re}s>0),\qquad\int_0^\infty\mathfrak k=\widehat F(1)=\log2-C.
\end{equation}

\begin{proposition}[Dilation sums of $\mathfrak k$]\label{prop:muntz}\VCGreviewed{1}
Let $t>0$ with $1/t\notin\mathbb Z$ and put $x=1/t$. Then
\begin{equation}\label{eq:muntz}
\sum_{n\ge1}\mathfrak k(nt)-\frac1t\int_0^\infty\mathfrak k=R(x),\qquad\sum_{n\ge1}\mu(n)\,\mathfrak k(nt)=H(x)\quad(x\ge1),
\end{equation}
with $R$ as in Lemma~\CGref{lem:EG}\textup{(iii)}; both sums have the finitely many nonzero terms $n\le x$, and for $x<1$ the second sum is empty. For $\operatorname{Re}s>1$ the Mellin transforms are
\begin{equation}\label{eq:muntzmellin}
\int_0^\infty\Bigl(\sum_{n\ge1}\mathfrak k(nt)\Bigr)t^{s-1}\,dt=\zeta(s)\widehat F(s),\qquad\int_0^\infty\Bigl(\sum_{n\ge1}\mu(n)\mathfrak k(nt)\Bigr)t^{s-1}\,dt=\frac{\widehat F(s)}{\zeta(s)}.
\end{equation}
\end{proposition}

\begin{proof}
$\mathfrak k(nt)\neq0$ requires $nt\le1$, i.e.\ $n\le x$, and then $\mathfrak k(nt)=F(1/(nt))=F(x/n)$. So $\sum_n\mathfrak k(nt)=\sum_{n\le x}F(x/n)$, and subtracting $(\log2-C)x$ from \CGeqref{eq:kmellin} gives $R(x)$; the second sum is \CGeqref{eq:Hdef}. For \CGeqref{eq:muntzmellin}, $\int_0^\infty\mathfrak k(nt)t^{s-1}dt=n^{-s}\widehat F(s)$ by the substitution $u=nt$, and $\sum_nn^{-\operatorname{Re}s}\int_0^\infty|\mathfrak k(u)|u^{\operatorname{Re}s-1}du<\infty$ justifies the interchange by Fubini's theorem. The second transform is \CGeqref{eq:HMellin} in the variable $t=1/x$.
\end{proof}

The first operator in \CGeqref{eq:muntz}, $\mathfrak k\mapsto\sum_n\mathfrak k(n\,\cdot)-(1/t)\int\mathfrak k$, is the M\"untz (co-Poisson) transform of the strong Nyman--Beurling criterion of B\'aez-Duarte \cite{BaezDuarte2005}. By \CGeqref{eq:muntz} and Lemma~\CGref{lem:EG}\textup{(iii)}, the M\"untz transform of $\mathfrak k$ at $t$ is $R(1/t)=-E_G'(1/t)/(2\pi)$: the function whose integer singularities stopped the attempts of Section~\CGref{sec:attempts} is the M\"untz transform of the log-sine kernel. According to the abstract of \cite{BaezDuarte2005} (an unrefereed arXiv text), its necessity theorem (RH implies that $\mathfrak k$ lies in the $L^2$ closure of the span of the dilations of its M\"untz transform) requires $\mathfrak k$ continuously differentiable and vanishing at infinity with an integrability condition on $t|\mathfrak k'(t)|$, and its sufficiency theorem requires compact support in addition. The kernel \CGeqref{eq:kdef} is not continuously differentiable: it has logarithmic singularities at $t=1/n$, $n\ge1$. So that theorem is not applied here, in either direction, and nothing below depends on it.

For $N\ge1$ and $\sigma>0$ define the energy and the Gram matrix
\begin{equation}\label{eq:energy}
\begin{split}
\mathcal E_N(\sigma)&=\int_0^\infty\Bigl|\sum_{d\le N}\mu(d)\,\mathfrak k(dt)\Bigr|^2t^{2\sigma-1}\,dt,\\
\Pi_\sigma(d,e)&=\int_0^\infty\mathfrak k(dt)\mathfrak k(et)\,t^{2\sigma-1}\,dt=\int_{\max(d,e)}^\infty F\Bigl(\frac xd\Bigr)F\Bigl(\frac xe\Bigr)\frac{dx}{x^{2\sigma+1}},
\end{split}
\end{equation}
the last form by $x=1/t$, since $\mathfrak k(dt)\mathfrak k(et)\neq0$ requires $x\ge\max(d,e)$.

\begin{proposition}[Spectral form of the energy]\label{prop:energy}\VCGreviewed{1}
Let $\sigma>0$ and $N\ge1$. Then $\mathcal E_N(\sigma)<\infty$, and
\begin{equation}\label{eq:energyspec}
\mathcal E_N(\sigma)=\frac1{2\pi}\int_{-\infty}^\infty|\widehat F(\sigma+i\tau)|^2\Bigl|\sum_{d\le N}\frac{\mu(d)}{d^{\sigma+i\tau}}\Bigr|^2d\tau=\sum_{d,e\le N}\mu(d)\mu(e)\,\Pi_\sigma(d,e),
\end{equation}
\begin{equation}\label{eq:gram}
\Pi_\sigma(d,e)=\frac{\Delta_\sigma(\log(e/d))}{(de)^\sigma},\qquad\Delta_\sigma(u)=\frac1{2\pi}\int_{-\infty}^\infty|\widehat F(\sigma+i\tau)|^2e^{i\tau u}\,d\tau,\qquad\Pi_\sigma(d,d)=\frac{\mathcal N_\sigma}{d^{2\sigma}}.
\end{equation}
Here $\tau\mapsto\widehat F(\sigma+i\tau)$ lies in $L^2(\mathbb R)$ with $\frac1{2\pi}\int|\widehat F(\sigma+i\tau)|^2d\tau=\mathcal N_\sigma$, $\Delta_\sigma$ is real, even, continuous, $|\Delta_\sigma|\le\Delta_\sigma(0)=\mathcal N_\sigma$, and the matrix $(\Pi_\sigma(d,e))_{d,e\le N}$ is positive semidefinite. Consequently
\begin{equation}\label{eq:diag}
\sum_{d\le N}\mu(d)^2\,\Pi_\sigma(d,d)\le\mathcal N_\sigma\,\zeta(2\sigma)\qquad(\sigma>\tfrac12),
\end{equation}
uniformly in $N$.
\end{proposition}

\begin{proof}
For $d\ge1$ let $\mathfrak k_d(t)=\mathfrak k(dt)$ and $g_d(v)=\mathfrak k_d(e^v)e^{\sigma v}$, $v\in\mathbb R$. By the substitutions $t=e^v$ and $u=dt$, $\int_{\mathbb R}|g_d|\,dv=\int_0^\infty|\mathfrak k(dt)|t^{\sigma-1}dt=d^{-\sigma}\int_0^\infty|\mathfrak k(u)|u^{\sigma-1}du<\infty$ and $\int_{\mathbb R}|g_d|^2dv=\int_0^\infty\mathfrak k(dt)^2t^{2\sigma-1}dt=d^{-2\sigma}\mathcal N_\sigma$. So $g_d\in L^1\cap L^2(\mathbb R)$, its Fourier transform is $\int g_d(v)e^{-i\tau v}dv=\int_0^\infty\mathfrak k(dt)t^{\sigma+i\tau-1}dt=d^{-\sigma-i\tau}\widehat F(\sigma+i\tau)$ by \CGeqref{eq:kmellin}, and Plancherel's theorem \cite[Theorem~9.13, p.~186]{Rudin1987} applied to $g=\sum_{d\le N}\mu(d)g_d$ gives the first equality in \CGeqref{eq:energyspec}, with $\mathcal E_N(\sigma)=\int|g|^2dv<\infty$. With $d=1$ it gives $\frac1{2\pi}\int|\widehat F(\sigma+i\tau)|^2d\tau=\mathcal N_\sigma$. Expanding the square in \CGeqref{eq:energy} gives the second equality in \CGeqref{eq:energyspec}; the matrix is the Gram matrix of the functions $\mathfrak k_dt^{\sigma-1/2}$ in $L^2(0,\infty)$, hence positive semidefinite. Plancherel's theorem with polarization, applied to $g_d$ and $g_e$, gives
\begin{equation}\label{eq:gramproof}
\Pi_\sigma(d,e)=\int_{\mathbb R}g_dg_e\,dv=\frac1{2\pi}\int_{-\infty}^\infty d^{-\sigma-i\tau}\,\overline{e^{-\sigma-i\tau}}\,|\widehat F(\sigma+i\tau)|^2\,d\tau=\frac1{(de)^\sigma}\,\Delta_\sigma\Bigl(\log\frac ed\Bigr),
\end{equation}
where $\overline{\widehat F(\sigma+i\tau)}=\widehat F(\sigma-i\tau)$ because $F$ is real. Since $|\widehat F(\sigma+i\cdot)|^2\in L^1(\mathbb R)$ is even and nonnegative, $\Delta_\sigma$ is real, even, continuous and bounded by $\Delta_\sigma(0)=\mathcal N_\sigma$. The diagonal is $\Pi_\sigma(d,d)=\int\mathfrak k(dt)^2t^{2\sigma-1}dt=d^{-2\sigma}\mathcal N_\sigma$, and $\sum_{d\le N}\mu(d)^2d^{-2\sigma}\le\zeta(2\sigma)$ for $\sigma>\frac12$.
\end{proof}

\begin{theorem}[A sufficient condition for RH]\label{thm:energy}\VCGreviewed{1}
If for every $\sigma>\frac12$ one has $\liminf_{N\to\infty}\mathcal E_N(\sigma)<\infty$, then RH holds. In particular $\sup_{N}\mathcal E_N(\sigma)<\infty$ for every $\sigma>\frac12$ implies RH.
\end{theorem}

\begin{proof}
Let $1/N\le t\le1$ with $x=1/t\notin\mathbb Z$. Then $x\le N$, so every $d\le x$ satisfies $d\le N$, and the terms with $x<d\le N$ vanish because $dt>1$. Hence $\sum_{d\le N}\mu(d)\mathfrak k(dt)=H(x)$ by \CGeqref{eq:muntz}, and
\begin{equation}\label{eq:energylower}
\mathcal E_N(\sigma)\ge\int_{1/N}^1\Bigl|\sum_{d\le N}\mu(d)\mathfrak k(dt)\Bigr|^2t^{2\sigma-1}\,dt=\int_1^N|H(x)|^2x^{-2\sigma-1}\,dx.
\end{equation}
Letting $N\to\infty$ along a sequence realising the $\liminf$, $\int_1^\infty|H(x)|^2x^{-2\sigma-1}dx<\infty$ for every $\sigma>\frac12$, which is condition (b) of Theorem~\CGref{thm:L2}, and (b) implies RH there.
\end{proof}

The converse of Theorem~\CGref{thm:energy} is not proved: by \CGeqref{eq:energyspec}, it would need a bound for $\int|\widehat F(\sigma+i\tau)|^2|\sum_{d\le N}\mu(d)d^{-\sigma-i\tau}|^2d\tau$ uniform in $N$, and no such bound is established here. By \CGeqref{eq:diag}, the diagonal of the quadratic form \CGeqref{eq:energyspec} is bounded uniformly in $N$ for $\sigma>\frac12$ without any hypothesis, and the sign of $\mu$ enters only through the off-diagonal part. The statement
\begin{equation}\label{eq:target2}
\sup_{N\ge1}\sum_{\substack{d,e\le N\\ d\neq e}}\frac{\mu(d)\mu(e)}{(de)^\sigma}\,\Delta_\sigma\Bigl(\log\frac ed\Bigr)<\infty\qquad\text{for every }\sigma>\tfrac12
\end{equation}
is therefore sufficient for RH by Proposition~\CGref{prop:energy} and Theorem~\CGref{thm:energy}. It is not proved, and it is not known to be necessary.

\section{A weighted strong Nyman--Beurling criterion for the log-sine kernel}\label{sec:nb}

Throughout this section $\tfrac12<\sigma<1$, and $\mathcal H_\sigma=L^2\bigl((0,\infty),t^{2\sigma-1}dt\bigr)$ with norm $\|\cdot\|_\sigma$. For $\phi\in\mathcal H_\sigma$ with $\int_0^\infty|\phi(t)|t^{\sigma-1}dt<\infty$ the Mellin transform $\int_0^\infty\phi(t)t^{s-1}dt$ converges absolutely on $\operatorname{Re}s=\sigma$, and Plancherel's theorem \cite[Theorem~9.13, p.~186]{Rudin1987} in the variable $v=\log t$ gives $\|\phi\|_\sigma^2=\frac1{2\pi}\int_{-\infty}^\infty|\int_0^\infty\phi(t)t^{\sigma+i\tau-1}dt|^2d\tau$; this is the fact used in the proof of Proposition~\CGref{prop:energy}, and it is used below without further comment. The M\"untz transform of $\mathfrak k$ is written
\begin{equation}\label{eq:Pk}
(P\mathfrak k)(t)=\sum_{n\ge1}\mathfrak k(nt)-\frac1t\int_0^\infty\mathfrak k=R(1/t)\quad(0<t\le1),\qquad(P\mathfrak k)(t)=\frac{C-\log2}t\quad(t>1),
\end{equation}
by \CGeqref{eq:muntz} and \CGeqref{eq:kmellin}; for $t>1$ the sum is empty and $R(x)=-(\log2-C)x$ for $x<1$, so the two lines agree with $R(1/t)$ for all $t>0$ with $1/t\notin\mathbb Z$.

\begin{lemma}[Uniform convexity bound]\label{lem:convex}\VCGreviewed{1}
Let $\tfrac12<\sigma_0<1$ and $\varepsilon>0$. Then $|\zeta(\sigma'+i\tau)|\ll_{\sigma_0,\varepsilon}|\tau|^{(1-\sigma_0)/2+\varepsilon}$ uniformly for $\sigma_0\le\sigma'\le2$, $|\tau|\ge2$.
\end{lemma}

\begin{proof}
Put $\phi(s)=\zeta(s)(s-1)/(s+1)$, holomorphic on $\operatorname{Re}s\ge\sigma_0$. In the half-plane $\operatorname{Re}s\ge\sigma_0$ one has $\zeta(s)=O(|\tau|^{\kappa_0})$ for some $\kappa_0$ \cite[\S5.1, p.~95]{TitchmarshHB1986}, so $\phi(s)=O(e^{\varepsilon'|\tau|})$ for every $\varepsilon'>0$ in the strip $\sigma_0\le\operatorname{Re}s\le2$. On $\operatorname{Re}s=2$, $|\phi(s)|\le\zeta(2)$ because $|s-1|\le|s+1|$. On $\operatorname{Re}s=\sigma_0$, $|\phi(s)|\le|\zeta(\sigma_0+i\tau)|=O(|\tau|^{(1-\sigma_0)/2+\varepsilon})$ by \cite[(5.1.4), p.~96]{TitchmarshHB1986} and the definition of $\mu(\sigma_0)$ there. The Phragm\'en--Lindel\"of theorem for a strip \cite[\S5.65, pp.~180--181]{Titchmarsh1939} gives $\phi(\sigma'+i\tau)=O(|\tau|^{\kappa(\sigma')})$ uniformly for $\sigma_0\le\sigma'\le2$, with $\kappa$ linear, $\kappa(\sigma_0)=(1-\sigma_0)/2+\varepsilon$, $\kappa(2)=0$, hence $\kappa(\sigma')\le(1-\sigma_0)/2+\varepsilon$. Finally $|\zeta(s)|=|\phi(s)||s+1|/|s-1|\le2|\phi(s)|$ for $|\tau|\ge2$, since $|s+1|\le|s-1|+2\le2|s-1|$.
\end{proof}

\begin{lemma}[$R$ in the weighted spaces]\label{lem:Rweight}\VCGreviewed{1}
\textup{(i)} $\int_1^X R(x)^2dx\le\frac{\pi^2}{12}X^3+2(\log2-C)^2X^3$ for $X\ge1$, so that $\widehat R(s)=\int_1^\infty R(x)x^{-s-1}dx$ converges absolutely for $\operatorname{Re}s>1$, and there
\begin{equation}\label{eq:Rhat}
\widehat R(s)=\zeta(s)\widehat F(s)-\frac{\widehat F(1)}{s-1}.
\end{equation}
The right-hand side is holomorphic on $\operatorname{Re}s>0$.
\textup{(ii)} For every $\sigma_0>\tfrac12$, $\int_1^\infty R(x)^2x^{-2\sigma_0-1}dx<\infty$; the integral in \CGeqref{eq:Rhat} converges absolutely for $\operatorname{Re}s>\tfrac12$ and \CGeqref{eq:Rhat} holds there.
\textup{(iii)} For $\tfrac12<\sigma<1$: $P\mathfrak k\in\mathcal H_\sigma$, $\int_0^\infty|P\mathfrak k(t)|t^{\sigma-1}dt<\infty$, and $\int_0^\infty P\mathfrak k(t)\,t^{s-1}dt=\zeta(s)\widehat F(s)$ on $\operatorname{Re}s=\sigma$.
\end{lemma}

\begin{proof}
(i) By Cauchy--Schwarz over the at most $X$ terms, $(\sum_{\ell\le x}|F(x/\ell)|)^2\le X\sum_{\ell\le x}F(x/\ell)^2$ for $x\le X$; and $\int_\ell^XF(x/\ell)^2dx\le X\pi^2/12$ for each $\ell\le X$ (periodicity, $\int_0^1F^2=\pi^2/12$), so $\int_1^X(\sum_{\ell\le x}|F(x/\ell)|)^2dx\le X\cdot X\cdot X\pi^2/12$, and $(a+b)^2\le2a^2+2b^2$ with $b=(\log2-C)x$ gives the bound. Dyadic decomposition and Cauchy--Schwarz give $\int_1^\infty|R|x^{-\operatorname{Re}s-1}dx<\infty$ for $\operatorname{Re}s>1$. For such $s$, $\sum_\ell\int_\ell^\infty|F(x/\ell)|x^{-\operatorname{Re}s-1}dx=\zeta(\operatorname{Re}s)\int_1^\infty|F(v)|v^{-\operatorname{Re}s-1}dv<\infty$ by the substitution $x=\ell v$, so Fubini's theorem gives $\int_1^\infty\sum_{\ell\le x}F(x/\ell)x^{-s-1}dx=\zeta(s)\widehat F(s)$, and $\int_1^\infty\widehat F(1)x\cdot x^{-s-1}dx=\widehat F(1)/(s-1)$. The residue of $\zeta(s)\widehat F(s)$ at $s=1$ is $\widehat F(1)$, so the right-hand side of \CGeqref{eq:Rhat} is holomorphic on $\operatorname{Re}s>0$ ($\widehat F$ is holomorphic on $\operatorname{Re}s>-1$ by Theorem~\CGref{thm:mellin}).

(ii) Fix $\sigma_0\in(\tfrac12,1)$ and $\varepsilon\in(0,(2\sigma_0-1)/4)$. For $\sigma_0\le\sigma'\le\tfrac32$ and $|\tau|\ge2$, \CGeqref{eq:Fhat} with $|s|\ge|\tau|$, $|\cot(\pi s/2)|\le\coth(\pi|\tau|/2)\le\coth\pi$ and $|A(s)|\le c\log|\tau|/|\tau|$ (the bound on $A$ obtained in the proof of Theorem~\CGref{thm:L2} from $|2k-s|\ge\max(|\tau|,k/2)$) gives $|\widehat F(\sigma'+i\tau)|\le(\pi\coth\pi\,|\zeta(\sigma'+i\tau)|+c\log|\tau|)/|\tau|$, so by Lemma~\CGref{lem:convex} and \CGeqref{eq:Rhat}
\begin{equation}\label{eq:Rhatbound}
|\widehat R(\sigma'+i\tau)|\ll|\tau|^{-\sigma_0+2\varepsilon}+|\tau|^{-(1+\sigma_0)/2+\varepsilon}\log|\tau|+|\tau|^{-1}\ll|\tau|^{-\sigma_0+2\varepsilon},
\end{equation}
uniformly in $\sigma'$, and $2\sigma_0-4\varepsilon>1$; on $|\tau|\le2$ the function is bounded. Hence $\sup_{\sigma_0\le\sigma'\le3/2}\int|\widehat R(\sigma'+i\tau)|^2d\tau<\infty$. For $\sigma'\ge\tfrac32$ the function $v\mapsto R(e^v)e^{-\sigma'v}$ lies in $L^1\cap L^2(0,\infty)$ by (i), so Plancherel's theorem gives $\int|\widehat R(\sigma'+i\tau)|^2d\tau=2\pi\int_1^\infty R^2x^{-2\sigma'-1}dx\le2\pi\int_1^\infty R^2x^{-4}dx$. Thus $\sup_{\sigma'>\sigma_0}\int|\widehat R(\sigma'+i\tau)|^2d\tau<\infty$. Put $\mathcal P(z)=\widehat R(\sigma_0-iz)$ for $\operatorname{Im}z>0$. By the Paley--Wiener theorem \cite[Theorem~19.2, pp.~372--373]{Rudin1987} there is $\Psi\in L^2(0,\infty)$ with $\mathcal P(z)=\int_0^\infty\Psi(v)e^{ivz}dv$; for $\upsilon>\tfrac32-\sigma_0$ also $\mathcal P(\xi+i\upsilon)=\int_0^\infty R(e^v)e^{-(\sigma_0+\upsilon)v}e^{i\xi v}dv$, and both $\Psi(v)e^{-\upsilon v}$ and $R(e^v)e^{-(\sigma_0+\upsilon)v}$ lie in $L^1(0,\infty)$ with the same Fourier transform, so they agree almost everywhere \cite[Theorem~9.12, p.~185]{Rudin1987}. Hence $\int_1^\infty R^2x^{-2\sigma_0-1}dx=\int_0^\infty|\Psi|^2<\infty$. Absolute convergence of $\widehat R(s)$ for $\operatorname{Re}s>\tfrac12$ follows by Cauchy--Schwarz, $\int_1^\infty|R|x^{-\operatorname{Re}s-1}\le(\int R^2x^{-2\sigma_0-1})^{1/2}(\int x^{2\sigma_0-2\operatorname{Re}s-1})^{1/2}$ with $\tfrac12<\sigma_0<\operatorname{Re}s$, locally uniformly, so the integral is holomorphic on $\operatorname{Re}s>\tfrac12$ and agrees with the right-hand side of \CGeqref{eq:Rhat} there by the identity theorem.

(iii) By \CGeqref{eq:Pk} and $x=1/t$: $\int_0^1|P\mathfrak k|^2t^{2\sigma-1}dt=\int_1^\infty R^2x^{-2\sigma-1}dx<\infty$ by (ii), and $\int_1^\infty(\log2-C)^2t^{2\sigma-3}dt=(\log2-C)^2/(2-2\sigma)$ since $\sigma<1$. Likewise $\int_0^1|P\mathfrak k|t^{\sigma-1}dt=\int_1^\infty|R|x^{-\sigma-1}dx<\infty$ by (ii), and $\int_1^\infty t^{\sigma-2}dt<\infty$. On $\operatorname{Re}s=\sigma$, $\int_0^1P\mathfrak k(t)t^{s-1}dt=\widehat R(s)$ and $\int_1^\infty(C-\log2)t^{s-2}dt=(\log2-C)/(s-1)=\widehat F(1)/(s-1)$, and \CGeqref{eq:Rhat} gives the sum $\zeta(s)\widehat F(s)$.
\end{proof}

For $n\ge1$ let
\begin{equation}\label{eq:varrho}
\varrho_n(t)=(P\mathfrak k)(nt)-\frac1n(P\mathfrak k)(t)\qquad(t>0).
\end{equation}
For $t>1$ both terms equal $(C-\log2)/(nt)$, so $\varrho_n(t)=0$: $\varrho_n$ is supported in $(0,1]$. By Lemma~\CGref{lem:Rweight}(iii) and the substitution $u=nt$, $\varrho_n\in\mathcal H_\sigma$, $\int_0^\infty|\varrho_n|t^{\sigma-1}dt<\infty$, and its Mellin transform on $\operatorname{Re}s=\sigma$ is $(n^{-s}-n^{-1})\zeta(s)\widehat F(s)$. For a Dirichlet polynomial $D(s)=\sum_{n\le N}c_nn^{-s}$ with real $c_n$ put $\tilde D(s)=D(s)-D(1)$, so that $\tilde D(1)=0$; then
\begin{equation}\label{eq:approxmellin}
\int_0^\infty\Bigl(\mathfrak k(t)-\sum_{n\le N}c_n\varrho_n(t)\Bigr)t^{s-1}dt=\widehat F(s)\bigl(1-\zeta(s)\tilde D(s)\bigr)\qquad(\operatorname{Re}s=\sigma),
\end{equation}
and the function $\Phi_D(s)=\widehat F(s)(1-\zeta(s)\tilde D(s))$ is holomorphic on $\operatorname{Re}s>0$, the pole of $\zeta$ at $s=1$ being cancelled by $\tilde D(1)=0$. Conversely every Dirichlet polynomial $E$ with $E(1)=0$ is $\tilde D$ for $D=E$.

\begin{theorem}[Weighted strong criterion]\label{thm:nb}\VCGreviewed{1}
Let $\tfrac12<\sigma<1$.
\textup{(i)} If $\mathfrak k$ lies in the $\mathcal H_\sigma$-closure of the linear span of $\{\varrho_n:n\ge1\}$, then $\zeta(s)\neq0$ for $\operatorname{Re}s>\sigma$.
\textup{(ii)} Assume RH. Then for every $\varepsilon>0$,
\begin{equation}\label{eq:nbrate}
\Bigl\|\mathfrak k-\sum_{n\le N}\mu(n)\varrho_n\Bigr\|_\sigma\ll_{\sigma,\varepsilon}N^{-(\sigma-1/2)+\varepsilon}\qquad(N\ge2).
\end{equation}
Consequently RH holds if and only if $\mathfrak k$ lies in the $\mathcal H_\sigma$-closure of $\operatorname{span}\{\varrho_n\}$ for every $\sigma\in(\tfrac12,1)$.
\end{theorem}

\begin{proof}
(i) Suppose $\zeta(\rho)=0$ with $\operatorname{Re}\rho>\sigma$; then $\operatorname{Re}\rho<1$. Let $\eta>0$ and choose real $c_n$ with $\|\phi\|_\sigma<\eta$, where $\phi=\mathfrak k-\sum_{n\le N}c_n\varrho_n$. The function $\phi$ is supported in $(0,1]$, so for $\sigma'\ge\sigma$, $\|\phi\|_{\sigma'}^2=\int_0^1|\phi|^2t^{2\sigma'-1}dt\le\|\phi\|_\sigma^2<\eta^2$, and $\int_0^1|\phi|t^{\sigma'-1}dt\le\int_0^1|\phi|t^{\sigma-1}dt<\infty$. For $\operatorname{Re}s>\sigma$ the Mellin integral of $\phi$ therefore converges absolutely, is holomorphic, and equals $\Phi_D(s)$: the computation that gives \CGeqref{eq:approxmellin} holds on every line $\operatorname{Re}s=\sigma'\in[\sigma,1)$ (Lemma~\CGref{lem:Rweight}(iii) applies to every such $\sigma'$), so the two holomorphic functions agree on the open strip $\sigma<\operatorname{Re}s<1$, hence on $\operatorname{Re}s>\sigma$ by the identity theorem. Plancherel on each line gives $\frac1{2\pi}\int|\Phi_D(\sigma'+i\tau)|^2d\tau=\|\phi\|_{\sigma'}^2<\eta^2$ for all $\sigma'\ge\sigma$. Put $\mathcal P(z)=\Phi_D(\sigma-iz)$ for $\operatorname{Im}z>0$. By the Paley--Wiener theorem \cite[Theorem~19.2, pp.~372--373]{Rudin1987} there is $\Psi\in L^2(0,\infty)$ with $\mathcal P(z)=\int_0^\infty\Psi(v)e^{ivz}dv$ and $\int_0^\infty|\Psi|^2=\sup_{\upsilon>0}\frac1{2\pi}\int|\mathcal P(\xi+i\upsilon)|^2d\xi\le\eta^2$. For $z=\xi+i\upsilon$, Cauchy--Schwarz gives $|\mathcal P(z)|^2\le\int|\Psi|^2\int_0^\infty e^{-2\upsilon v}dv\le\eta^2/(2\upsilon)$. With $z$ such that $\sigma-iz=\rho$, i.e.\ $\upsilon=\operatorname{Re}\rho-\sigma>0$,
\begin{equation}\label{eq:nbpoint}
|\Phi_D(\rho)|^2\le\frac{\eta^2}{2(\operatorname{Re}\rho-\sigma)}.
\end{equation}
But $\Phi_D(\rho)=\widehat F(\rho)(1-\zeta(\rho)\tilde D(\rho))=\widehat F(\rho)=A(\rho)$ by \CGeqref{eq:Fhat}, and $A(\rho)\neq0$ by Theorem~\CGref{thm:Anonzero}. Since $\eta>0$ is arbitrary, this is a contradiction.

(ii) Assume RH. Fix $\delta$ with $0<\delta<\sigma-\tfrac12$, put $\eta_0=\sigma-\tfrac12-\delta>0$, and let $\varepsilon\in(0,\eta_0/8)$. Under RH, for $\tfrac12+\delta\le\operatorname{Re}u\le2$ and $|\operatorname{Im}u|\ge2$,
\begin{equation}\label{eq:RHbounds}
|\zeta(u)|\ll_{\delta,\varepsilon}|\operatorname{Im}u|^\varepsilon,\qquad|1/\zeta(u)|\ll_{\delta,\varepsilon}|\operatorname{Im}u|^\varepsilon:
\end{equation}
for $\tfrac12+\delta\le\operatorname{Re}u\le1$ this is \cite[Theorem~14.2, (14.2.1), p.~336]{TitchmarshHB1986}, which is uniform in that range and gives $|\log\zeta(u)|\ll(\log|\operatorname{Im}u|)^{1-2\delta+\varepsilon}$; for $1\le\operatorname{Re}u\le2$, $|\zeta(u)|\ll\log|\operatorname{Im}u|$ by \cite[Theorem~4.11, p.~77]{TitchmarshHB1986} with a partial sum of length $|\operatorname{Im}u|$, and $|1/\zeta(u)|\ll\log^7|\operatorname{Im}u|$ by \cite[(3.6.5), p.~51]{TitchmarshHB1986}. On the compact set $\tfrac12+\delta\le\operatorname{Re}u\le2$, $|\operatorname{Im}u|\le2$, $1/\zeta$ is continuous under RH (no zeros, and $1/\zeta(1)=0$), hence bounded. So $|1/\zeta(u)|\ll(2+|\operatorname{Im}u|)^\varepsilon$ on the whole strip.

Let $N\ge2$, $x=N+\tfrac12$, $T=N^2$, $c=2-\sigma$, and $s=\sigma+i\tau$. The lemma of \cite[\S3.12, Lemma~3.12, pp.~60--61]{TitchmarshHB1986} applied to $f=1/\zeta$ (coefficients $\mu(n)$, so $\psi\equiv1$ and $\sum|\mu(n)|n^{-\sigma'}\le\zeta(\sigma')\ll(\sigma'-1)^{-1}$, i.e.\ $\alpha=1$), with $\sigma+c=2>1$ and $|x-N|=\tfrac12$, gives
\begin{equation}\label{eq:perron}
D_N(s):=\sum_{n\le N}\frac{\mu(n)}{n^s}=\frac1{2\pi i}\int_{c-iT}^{c+iT}\frac{x^w}{\zeta(s+w)}\frac{dw}w+O\Bigl(\frac{x^c}{T}\Bigr)+O\Bigl(\frac{x^{1-\sigma}\log x}T\Bigr),
\end{equation}
and the error terms are $\ll N^{-\sigma}$. Move the segment to $\operatorname{Re}w=-\eta_0$, so that $\operatorname{Re}(s+w)=\tfrac12+\delta$. In the rectangle $-\eta_0\le\operatorname{Re}w\le c$, $|\operatorname{Im}w|\le T$, the integrand has the single pole $w=0$, with residue $1/\zeta(s)$: $1/\zeta(s+w)$ is holomorphic there under RH, and $w=1-s$ is a zero of it. On the horizontal segments $w=v\pm iT$, $-\eta_0\le v\le c$, by \CGeqref{eq:RHbounds} $|1/\zeta(s+w)|\ll(T+|\tau|+2)^\varepsilon$ and $|x^w/w|\le x^c/T$, so their contribution is $\ll(N^2+|\tau|)^\varepsilon N^{-\sigma}$. On the vertical segment $w=-\eta_0+iv$, $|v|\le T$,
\begin{equation}\label{eq:perronvert}
\Bigl|\int\frac{x^w}{\zeta(s+w)}\frac{dw}w\Bigr|\ll x^{-\eta_0}\int_{-T}^T\frac{(2+|\tau|+|v|)^\varepsilon}{\eta_0+|v|}\,dv\ll_{\eta_0}N^{-\eta_0}(N^2+|\tau|)^\varepsilon\log N.
\end{equation}
Since $(N^2+|\tau|)^\varepsilon\le N^{2\varepsilon}(2+|\tau|)^\varepsilon$ and $N^{-\sigma}\le N^{-\eta_0}$,
\begin{equation}\label{eq:DNbound}
\Bigl|D_N(s)-\frac1{\zeta(s)}\Bigr|\ll_{\sigma,\delta,\varepsilon}N^{-\eta_0+3\varepsilon}(2+|\tau|)^\varepsilon\qquad(\operatorname{Re}s=\sigma).
\end{equation}
Multiplying by $\zeta(s)$ and using $|\zeta(\sigma+i\tau)|\ll(2+|\tau|)^\varepsilon$ (by \CGeqref{eq:RHbounds} for $|\tau|\ge2$, and boundedness for $|\tau|\le2$, $\sigma<1$),
\begin{equation}\label{eq:oneminus}
|1-\zeta(s)D_N(s)|\ll N^{-\eta_0+3\varepsilon}(2+|\tau|)^{2\varepsilon}\qquad(\operatorname{Re}s=\sigma).
\end{equation}
Now $\tilde D_N=D_N-m(N)$, and by \CGeqref{eq:approxmellin}, Plancherel and \CGeqref{eq:oneminus},
\begin{equation}\label{eq:nbnorm}
\Bigl\|\mathfrak k-\sum_{n\le N}\mu(n)\varrho_n\Bigr\|_\sigma\le\Bigl(\frac1{2\pi}\int|\widehat F(\sigma+i\tau)|^2|1-\zeta D_N|^2d\tau\Bigr)^{1/2}+|m(N)|\,\|P\mathfrak k\|_\sigma,
\end{equation}
where the second term comes from $\widehat F(s)\zeta(s)m(N)$, the Mellin transform of $m(N)P\mathfrak k$ by Lemma~\CGref{lem:Rweight}(iii). The first term is at most $N^{-\eta_0+3\varepsilon}$ times a constant multiple of the square root of $\int|\widehat F(\sigma+i\tau)|^2(2+|\tau|)^{4\varepsilon}d\tau$, and this integral is finite because $|\widehat F(\sigma+i\tau)|\ll(|\zeta(\sigma+i\tau)|+\log|\tau|)/|\tau|\ll|\tau|^{-1+\varepsilon}$ for $|\tau|\ge2$ under RH, with $6\varepsilon<1$. The second term is $\ll N^{-1/2+\varepsilon}$ by \CGeqref{eq:mbound} and Lemma~\CGref{lem:Rweight}(iii). Altogether the norm is $\ll N^{-\eta_0+3\varepsilon}+N^{-1/2+\varepsilon}\ll N^{-(\sigma-1/2)+\delta+3\varepsilon}$, since $\eta_0=\sigma-\frac12-\delta$ and $\frac12\ge\sigma-\frac12$. Given $\varepsilon'>0$, choosing $\delta=\varepsilon'/2$ and $\varepsilon=\varepsilon'/6$ gives the exponent $-(\sigma-\frac12)+\varepsilon'$ of \CGeqref{eq:nbrate}.

The equivalence: (ii) gives one direction; for the other, (i) for every $\sigma\in(\tfrac12,1)$ gives $\zeta(s)\neq0$ for $\operatorname{Re}s>\tfrac12$, which is RH.
\end{proof}

The classical strong Nyman--Beurling criterion is the case of the kernel $\mathbf 1_{(0,1]}$, whose M\"untz transform is $-\{1/t\}$ and whose Mellin transform is $1/s$; Theorem~\CGref{thm:nb} replaces $1/s$ by $\widehat F$, the weight $1/|s|^2$ by $|\widehat F(\sigma+i\tau)|^2$, and the space $L^2(0,\infty)$ by $\mathcal H_\sigma$. The two ingredients that make the kernel $\mathfrak k$ admissible, although it is not continuously differentiable, are the support of $\mathfrak k-\sum c_n\varrho_n$ in $(0,1]$, which makes the weighted norms monotone in $\sigma'$, and $A\neq0$ at the zeros, which is what \CGeqref{eq:nbpoint} contradicts.

\begin{proposition}[Finite prime cut]\label{prop:primecut}\VCGreviewed{1}
For $Y\ge2$ let $\mathcal S_Y$ be the set of squarefree integers all of whose prime factors are $\le Y$, and
\begin{equation}\label{eq:zetaY}
\Pi_Y(s)=\prod_{p\le Y}(1-p^{-s})=\sum_{n\in\mathcal S_Y}\mu(n)n^{-s},\qquad\zeta_{>Y}(s)=\zeta(s)\Pi_Y(s),
\end{equation}
so that $\zeta_{>Y}(s)=\prod_{p>Y}(1-p^{-s})^{-1}$ for $\operatorname{Re}s>1$, and $\zeta_{>Y}$ has the same zeros as $\zeta$ in $0<\operatorname{Re}s<1$. Let $\tfrac12<\sigma<1$. If
\begin{equation}\label{eq:primecut}
\liminf_{Y\to\infty}\int_{-\infty}^\infty|\widehat F(\sigma+i\tau)|^2\,|1-\zeta_{>Y}(\sigma+i\tau)|^2\,d\tau=0,
\end{equation}
then $\zeta(s)\neq0$ for $\operatorname{Re}s>\sigma$. In particular, \CGeqref{eq:primecut} for every $\sigma\in(\tfrac12,1)$ implies RH.
\end{proposition}

\begin{proof}
$\mathcal S_Y$ is finite, so $\Pi_Y$ is a Dirichlet polynomial, and $\Pi_Y(1)=\prod_{p\le Y}(1-1/p)\le(\sum_{n\le Y}1/n)^{-1}\le1/\log Y$, because $\prod_{p\le Y}(1-1/p)^{-1}=\prod_{p\le Y}\sum_{k\ge0}p^{-k}\ge\sum_{n\le Y}1/n$. By \CGeqref{eq:approxmellin} with $D=\Pi_Y$ and Plancherel,
\begin{equation}\label{eq:primecutnorm}
\Bigl\|\mathfrak k-\sum_{n\in\mathcal S_Y}\mu(n)\varrho_n\Bigr\|_\sigma\le\Bigl(\frac1{2\pi}\int|\widehat F(\sigma+i\tau)|^2|1-\zeta_{>Y}(\sigma+i\tau)|^2d\tau\Bigr)^{1/2}+\frac{\|P\mathfrak k\|_\sigma}{\log Y},
\end{equation}
which tends to $0$ along a sequence of $Y$ by \CGeqref{eq:primecut} and Lemma~\CGref{lem:Rweight}(iii). Theorem~\CGref{thm:nb}(i) applies.
\end{proof}

The converse of Proposition~\CGref{prop:primecut}, that RH implies \CGeqref{eq:primecut}, is not proved here. The proof of Theorem~\CGref{thm:nb}(ii) does not transfer: it rests on Perron's formula for the sharp cutoff $n\le N$, whose error terms are controlled by the length $N$, whereas $\Pi_Y$ has length $\prod_{p\le Y}p$, and no bound for $|\zeta_{>Y}(\sigma+i\tau)|$ uniform in $Y$ on the line $\operatorname{Re}s=\sigma<1$ is available to us. The factor $\Pi_Y(s)$ is the inverse of the Euler product of $\zeta$ truncated at $Y$; as $Y\to\infty$ it converges to $1/\zeta(s)$ for $\operatorname{Re}s>1$ only, and on the line $\operatorname{Re}s=\sigma<1$ the statement \CGeqref{eq:primecut} is a statement about $\zeta$ itself, not about a function with finitely many parameters.

\section{The corner as a weighted Liouville sum, and the Clausen kernel}\label{sec:liouville}

Write $U=\sqrt X$ and group the corner by $n=dj$:
\begin{equation}\label{eq:aU}
\mathcal C(X)=\sum_{U<n\le X}a_U(n)\,n\,E_G\Bigl(\frac Xn\Bigr),\qquad a_U(n)=\sum_{\substack{dj=n\\ d\le U,\ j\le U}}\mu(d)\mu(j).
\end{equation}
Let $\lambda(n)=(-1)^{\Omega(n)}$ be Liouville's function, $\Omega(n)$ the number of prime factors of $n$ counted with multiplicity.

\begin{proposition}[Liouville form of the corner]\label{prop:liouville}\VCGreviewed{1}
For $n\ge1$ and $U\ge1$, $a_U(n)=\lambda(n)\,b_U(n)$ with $b_U(n)\ge0$ the number of factorizations $n=dj$ with $d,j\le U$ both squarefree. Explicitly, $b_U(n)=0$ unless $n=ab^2$ with $a,b$ squarefree and $(a,b)=1$, and then
\begin{equation}\label{eq:bU}
b_U(n)=\#\Bigl\{t\mid a:\ \frac{ab}U\le t\le\frac Ub\Bigr\}.
\end{equation}
Hence
\begin{equation}\label{eq:Cliouville}
\mathcal C(X)=\sum_{\sqrt X<n\le X}\lambda(n)\,b_{\sqrt X}(n)\,n\,E_G\Bigl(\frac Xn\Bigr),
\end{equation}
and by Theorem~\CGref{thm:corner}, RH is equivalent to $\sum_{\sqrt X<n\le X}\lambda(n)b_{\sqrt X}(n)\,nE_G(X/n)\ge-K_\varepsilon X^{3/2+\varepsilon}$ for all $X\ge1$, for every $\varepsilon>0$.
\end{proposition}

\begin{proof}
If $\mu(d)\mu(j)\neq0$ then $d$ and $j$ are squarefree, so $\Omega(d)=\omega(d)$, $\Omega(j)=\omega(j)$, and $\mu(d)\mu(j)=(-1)^{\omega(d)+\omega(j)}=(-1)^{\Omega(dj)}=\lambda(n)$. Thus every nonzero term of $a_U(n)$ equals $\lambda(n)$, which gives $a_U=\lambda b_U$ with $b_U(n)$ the number of such pairs. If $n=dj$ with $d,j$ squarefree, put $b=(d,j)$ and $a=n/b^2=(d/b)(j/b)$; the factors $d/b$ and $j/b$ are coprime and squarefree, so $a$ is squarefree, and a prime dividing $b$ divides neither $d/b$ nor $j/b$ (as $d$, $j$ are squarefree), so $(a,b)=1$. Conversely, given $n=ab^2$ with $a,b$ squarefree and coprime, the pairs are exactly $d=bt$, $j=ab/t$ with $t\mid a$ (both squarefree because $(t,b)=1$), and the condition $d,j\le U$ is $ab/U\le t\le U/b$. This is \CGeqref{eq:bU}; if $n$ has a cube factor it is not of the form $ab^2$ with $a,b$ squarefree, so $b_U(n)=0$. The identity \CGeqref{eq:bU} was also checked against the definition for every $n\le U^2$ with $U=7,20,45$ (\texttt{liouville\_corner.py}).
\end{proof}

In \CGeqref{eq:Cliouville} the sign of each term is $\lambda(n)\operatorname{sgn}E_G(X/n)$ and the weight $b_{\sqrt X}(n)$ is nonnegative: there is no cancellation inside $a_U(n)$, only between different $n$. The weight counts the divisors $t$ of $a$ in the multiplicative window $[ab/U,U/b]$, whose ratio is $U^2/(ab^2)=X/n$; it vanishes unless the largest prime factor of $n$ is at most $\sqrt X$.

\begin{proposition}[The Clausen kernel and its M\"untz transform]\label{prop:clausenmuntz}\VCGreviewed{1}
Let $\mathfrak c(t)=t\,\mathrm{Cl}_2(2\pi/t)$ for $0<t\le1$ and $\mathfrak c(t)=0$ for $t>1$. Then $\int_0^\infty\mathfrak c(t)t^{s-1}dt=\Theta(s)$ for $\operatorname{Re}s>0$, $\int_0^\infty\mathfrak c=a$, and for $Y\ge1$, $Y\notin\mathbb Z$,
\begin{equation}\label{eq:EGmuntz}
\frac{E_G(Y)}Y=\sum_{\ell\ge1}\mathfrak c\Bigl(\frac\ell Y\Bigr)-Y\int_0^\infty\mathfrak c=(P\mathfrak c)(1/Y),
\end{equation}
the M\"untz transform of $\mathfrak c$ at $t=1/Y$. Consequently
\begin{equation}\label{eq:Cmuntz}
\mathcal C(X)=X\sum_{\sqrt X<n\le X}\lambda(n)\,b_{\sqrt X}(n)\,(P\mathfrak c)\Bigl(\frac nX\Bigr),
\end{equation}
and, with $P_U(s)=\sum_{n\le U}\mu(n)n^{-s}$ and
\begin{equation}\label{eq:LU}
L_U(s)=\sum_{U<n\le X}\lambda(n)b_U(n)n^{-s}=P_U(s)^2-\sum_{n\le U}(\mu*\mu)(n)n^{-s},
\end{equation}
\begin{equation}\label{eq:Cspectral}
\mathcal C(X)=\frac1{2\pi i}\int_{(5/2)}\zeta(s-1)\Theta(s)\,L_U(s-1)\,X^s\,ds-aX^2\sum_{U<n\le X}\frac{\lambda(n)b_U(n)}n.
\end{equation}
\end{proposition}

\begin{proof}
With $v=1/t$, $\int_0^1t\,\mathrm{Cl}_2(2\pi/t)t^{s-1}dt=\int_1^\infty\mathrm{Cl}_2(2\pi v)v^{-s-1}dv=\Theta(s)$ by \CGeqref{eq:Thetadef}, and $\Theta(2)=a$. Next $\mathfrak c(\ell/Y)=(\ell/Y)\mathrm{Cl}_2(2\pi Y/\ell)$ for $\ell\le Y$ and $0$ for $\ell>Y$, so $\sum_\ell\mathfrak c(\ell/Y)=G(Y)/Y$, and subtracting $aY=Y\int\mathfrak c$ gives $E_G(Y)/Y$; this is the M\"untz transform $\sum_\ell\mathfrak c(\ell t)-t^{-1}\int\mathfrak c$ at $t=1/Y$. Then \CGeqref{eq:Cmuntz} is \CGeqref{eq:Cliouville} with $nE_G(X/n)=X\cdot E_G(X/n)/(X/n)$. For \CGeqref{eq:Cspectral}, \CGeqref{eq:Ginv} gives $G(X/n)=\frac1{2\pi i}\int_{(5/2)}\zeta(s-1)\Theta(s)(X/n)^sds$ for each of the finitely many $n$, and $E_G=G-aY^2$; the series $\sum_{d,j\le U}\mu(d)\mu(j)(dj)^{-s}=P_U(s)^2$ splits into $dj\le U$ and $dj>U$, which gives the formula for $L_U$.
\end{proof}

Thus the corner is $X$ times a weighted Liouville sum of dilations of the M\"untz transform of the Clausen kernel $\mathfrak c$, the kernel one integration above $\mathfrak k$ ($\mathfrak k(t)=F(1/t)$, $\mathfrak c(t)=t\,\mathrm{Cl}_2(2\pi/t)$, and $\mathrm{Cl}_2'=-F(\cdot/2\pi)$ by \CGeqref{eq:Cl2}). The Mellin factor $\zeta(s-1)\Theta(s)$ in \CGeqref{eq:Cspectral} is the transform of $G$; the arithmetic is entirely in $L_U(s-1)=P_U(s-1)^2-\sum_{n\le U}(\mu*\mu)(n)n^{1-s}$, a Dirichlet polynomial of length $U^2=X$ built from $\mu$ on $[1,\sqrt X]$ only.

\begin{proposition}[One-sided Liouville--Clausen criterion]\label{prop:lambdacrit}\VCGreviewed{1}
Let $\Lambda_G(X)=\sum_{n\le X}\lambda(n)\,n\,E_G(X/n)$ for $X\ge1$. Suppose that for every $\varepsilon>0$ there is $K_\varepsilon$ such that either $\Lambda_G(X)\le K_\varepsilon X^{3/2+\varepsilon}$ for all $X\ge1$, or $\Lambda_G(X)\ge-K_\varepsilon X^{3/2+\varepsilon}$ for all $X\ge1$. Then RH holds.
\end{proposition}

\begin{proof}
By Lemma~\CGref{lem:EG}(ii) with $\delta=\frac12$,
\begin{equation}\label{eq:lambdatrivial}
|\Lambda_G(X)|\le\sum_{n\le X}n\,c_{1/2}\Bigl(\frac Xn\Bigr)^{3/2}\le c_{1/2}X^{3/2}\sum_{n\le X}n^{-1/2}\le2c_{1/2}X^2,
\end{equation}
so $\widehat\Lambda_G(s)=\int_1^\infty\Lambda_G(X)X^{-s-1}dX$ converges absolutely for $\operatorname{Re}s>2$. There, by Fubini (the same bound gives absolute convergence of the double integral), $\widehat\Lambda_G(s)=\sum_n\lambda(n)n^{1-s}\int_n^\infty E_G(X/n)X^{-s-1}dX=\widehat E_G(s)\sum_n\lambda(n)n^{1-s}$, where $\widehat E_G(s)=\int_1^\infty E_GY^{-s-1}dY=\zeta(s-1)\Theta(s)-a/(s-2)$ by Lemma~\CGref{lem:EG}(i), and $\sum_n\lambda(n)n^{-w}=\zeta(2w)/\zeta(w)$ for $\operatorname{Re}w>1$. Hence
\begin{equation}\label{eq:lambdamellin}
\widehat\Lambda_G(s)=\Bigl(\zeta(s-1)\Theta(s)-\frac a{s-2}\Bigr)\frac{\zeta(2s-2)}{\zeta(s-1)}\qquad(\operatorname{Re}s>2).
\end{equation}
The right-hand side is meromorphic on $\operatorname{Re}s>1$; on the real half-line $s>\frac32$ it is holomorphic, since $\zeta(s-1)\neq0$ for real $s-1>0$ (proof of Proposition~\CGref{prop:Wbasic}(iii)), the first factor is holomorphic on $\operatorname{Re}s>0$, and the only real pole of $\zeta(2s-2)$ is $s=\frac32$. Fix $\varepsilon>0$, $\beta=\frac32+\varepsilon$, and suppose $\Lambda_G(X)\le K_\varepsilon X^\beta$. Then $f=K_\varepsilon X^\beta-\Lambda_G$ is nonnegative, locally bounded and measurable, its Mellin integral $g(s)=K_\varepsilon/(s-\beta)-\widehat\Lambda_G(s)$ converges absolutely for $\operatorname{Re}s>2$, and $g$ extends holomorphically to a neighbourhood of every real segment $[\sigma_2,2]$ with $\sigma_2>\beta$; Lemma~\CGref{lem:landau} gives absolute convergence of the Mellin integral of $f$, hence of $\widehat\Lambda_G$, on $\operatorname{Re}s>\beta$, so the right-hand side of \CGeqref{eq:lambdamellin} is holomorphic there. Let $\rho$ be a zero of $\zeta$ with $\frac12+\varepsilon<\operatorname{Re}\rho<1$. At $s=1+\rho$ the denominator vanishes, $\zeta(2s-2)=\zeta(2\rho)\neq0$ because $\operatorname{Re}2\rho>1$, and the first factor equals $-a/(\rho-1)\neq0$ because $\zeta(\rho)\Theta(1+\rho)=0$. So $s=1+\rho$ would be a pole in $\operatorname{Re}s>\beta$, a contradiction. The other one-sided bound is handled with $f=\Lambda_G+K_\varepsilon X^\beta$. Letting $\varepsilon\to0$ and using $\zeta\neq0$ on $\operatorname{Re}s\ge1$ gives RH.
\end{proof}

The converse, RH $\Rightarrow\Lambda_G(X)\ll X^{3/2+\varepsilon}$, is not written out here. Proposition~\CGref{prop:lambdacrit} uses $\lambda$ on all of $[1,X]$ and no weight; it is recorded because the measurements below show that the weight $b_{\sqrt X}$ in \CGeqref{eq:Cliouville} does not change the scale of the sum.

\begin{remark}[Measurements on the Liouville form]\label{rem:liouvillenum}\VCGreviewed{1}
At $X=30000.5$, $120000.5$, $380000.5$ (\texttt{liouville\_weight.py}): $b_{\sqrt X}(n)>0$ for $14.2\%$, $13.9\%$, $13.8\%$ of the $n\in(\sqrt X,X]$, with mean $2.48$, $2.61$, $2.66$ on its support and maximum $16$, $22$, $26$. The contributions of the $n$ with a fixed value $b_{\sqrt X}(n)=k$ are each of size $0.01$--$0.2\,X^{3/2}$ with varying signs ($k=2$: $-0.089$, $-0.085$, $-0.047$; $k=4$: $+0.220$, $-0.077$, $+0.189$). The indicator sum $\sum_{b>0}\lambda(n)nE_G(X/n)$ is $+0.026$, $-0.013$, $+0.033\,X^{3/2}$, and the unweighted $\sum_{\sqrt X<n\le X}\lambda(n)nE_G(X/n)$ is $-0.059$, $-0.095$, $+0.011\,X^{3/2}$, against $\sum_n n|E_G(X/n)|=81$, $163$, $290\,X^{3/2}$.

One-sided Abel form (\texttt{corner\_onesided.py}). With $g(n)=nE_G(X/n)$, $b_X(n)=g(n+1)-g(n)$ and $A_U(t)=\sum_{U<n\le t}a_U(n)$, $\mathcal C(X)=A_U(\lfloor X\rfloor)g(\lfloor X\rfloor)-\sum_{U<n<\lfloor X\rfloor}A_U(n)b_X(n)$ exactly (residual $<10^{-6}$), and $\mathcal C\ge-K_\varepsilon X^{3/2+\varepsilon}$ follows from $\sum_n[A_U(n)b_X(n)]_+\ll X^{3/2+\varepsilon}$ plus the boundary term. Measured, in units of $X^{3/2}$: $\sum[A_Ub_X]_+=3.77$, $7.01$, $5.92$; $\sum|A_Ub_X|=7.69$, $14.10$, $11.95$; boundary term $-0.007$, $+0.009$, $+0.002$; $\max_{n}|A_U(n)|/\sqrt X=0.36$, $0.45$, $0.41$. The Abel regrouping brings the sum of absolute values from $26$--$95\,X^{3/2}$ (termwise, Table~\CGref{tab:corner}) to $8$--$14\,X^{3/2}$, no longer growing like $X^{1/2}$ in this range; the positive part is one half of it, so the one-sided form gains nothing over the absolute one here.
\end{remark}

\begin{remark}[The $\mu^2$-corner]\label{rem:musqcorner}\VCGreviewed{1}
Let $\mathcal C_{\mu^2}(X)$ be \CGeqref{eq:Rest} with $\mu(d)\mu(j)$ replaced by $\mu(d)^2\mu(j)^2$. Replacing the sum over squarefree $d=Us$, $j=Ut$ by the integral with density $6/\pi^2$ in each variable suggests $\mathcal C_{\mu^2}(X)\sim\kappa X^2$ with
\begin{equation}\label{eq:kappa}
\begin{split}
\kappa&=\Bigl(\frac6{\pi^2}\Bigr)^2\int_0^1\!\!\int_0^1st\,E_G\Bigl(\frac1{st}\Bigr)ds\,dt=\Bigl(\frac6{\pi^2}\Bigr)^2\int_1^\infty\frac{E_G(y)\log y}{y^3}\,dy,\\
\int_1^\infty\frac{E_G(y)\log y}{y^3}\,dy&=-\widehat E_G{}'(2)=-\frac{\Theta''(2)}2-\gamma_E\Theta'(2)+\gamma_1a,
\end{split}
\end{equation}
where $\gamma_E$ is Euler's constant and $\gamma_1$ the first Stieltjes constant; the second equality uses $\int_0^1\!\int_0^1\phi(st)\,ds\,dt=\int_0^1\phi(p)\log(1/p)\,dp$ and $p=1/y$, the third $\widehat E_G(s)=\zeta(s-1)\Theta(s)-a/(s-2)$ expanded at $s=2$. Numerically (\texttt{musq\_corner\_const.py}) $\Theta'(2)=-0.0146808536$, $\Theta''(2)=-0.000882717$, the integral is $-0.00267067$ (confirmed by direct integration of $E_G$ up to $y=3000$: $-0.002671$), and $\kappa=-0.000987014$. The passage from the sum to the integral is not proved here, and the measured values of $\mathcal C_{\mu^2}(X)/X^2$ at $X=30000.5$, $120000.5$, $380000.5$ are $-0.000762$, $-0.000988$, $-0.001327$: they cross $\kappa$ and drift by about $-0.00022$ per unit of $\log X$, so the asymptotic $\kappa X^2$ is not confirmed in this range. What is used in Remark~\CGref{rem:barrier} is only that $\mathcal C_{\mu^2}$ is measured to be negative and of order $X^2$. The near-diagonal part $|d-j|\le H$ of the corner has at most $(H+1)\sqrt X$ pairs, each $\ll_\delta X^{1+\delta}(dj)^{-\delta}$ by Lemma~\CGref{lem:EG}(ii), hence is $\ll_\delta(H+1)X^{3/2}$ for $H\le X^{1/4}/2$; this is a count of pairs, not a cancellation, and it removes nothing from the problem.
\end{remark}

\section{The value of \texorpdfstring{$\widehat E_G$}{the transform of E\_G} at the zeros, and the centred primitive of the corner}\label{sec:centred}

\begin{proposition}[$\widehat E_G$ at $1+\rho$]\label{prop:EGrho}\VCGreviewed{1}
The integral $\widehat E_G(s)=\int_1^\infty E_G(Y)Y^{-s-1}dY$ converges absolutely for $\operatorname{Re}s>1$, and there
\begin{equation}\label{eq:EGhat}
\widehat E_G(s)=\zeta(s-1)\Theta(s)-\frac a{s-2},
\end{equation}
where the right-hand side is holomorphic on $\operatorname{Re}s>1$ (the singularity at $s=2$ is removable). In particular, for every zero $\rho$ of $\zeta$ with $\operatorname{Re}\rho>0$,
\begin{equation}\label{eq:EGrho}
\widehat E_G(1+\rho)=-\frac a{\rho-1}.
\end{equation}
\end{proposition}

\begin{proof}
By Lemma~\CGref{lem:EG}(ii), $|E_G(Y)|\le c_\delta Y^{1+\delta}$ for every $\delta\in(0,1)$, so the integral converges absolutely and locally uniformly on $\operatorname{Re}s>1+\delta$ for every $\delta>0$, and is holomorphic on $\operatorname{Re}s>1$. For $\operatorname{Re}s>2$, Lemma~\CGref{lem:EG}(i) and $\int_1^\infty aY^2Y^{-s-1}dY=a/(s-2)$ give \CGeqref{eq:EGhat}. The function $\Theta$ is holomorphic on $\operatorname{Re}s>0$ by \CGeqref{eq:Thetadef}, $\zeta(s-1)$ is meromorphic with a simple pole of residue $1$ at $s=2$, and $\Theta(2)=a$, so the right-hand side is holomorphic on $\operatorname{Re}s>1$; by the identity theorem \CGeqref{eq:EGhat} holds there. At $s=1+\rho$, $\operatorname{Re}s>1$ and $\zeta(\rho)=0$, which gives \CGeqref{eq:EGrho}.
\end{proof}

\begin{remark}[Bands and zeros]\label{rem:pairs}\VCGreviewed{1}
Proposition~\CGref{prop:EGrho} is the exact form of the following computation. Replacing $\mu(d)$ formally by $\sum_\rho d^{\rho-1}/\zeta'(\rho)$ in the corner, the substitution $d=\sqrt X\,x$, $j=\sqrt X\,y$, $p=xy$, $q=x/y$ shows that the pair $(\rho,\bar\rho)$, $\rho=\beta+i\gamma$, produces $X^{1+\beta}J(\rho)/|\zeta'(\rho)|^2$ with
\begin{equation}\label{eq:Jbands}
\begin{split}
J_k(\rho)&=\frac1\gamma\int_{2^k}^{2^{k+1}}E_G(Y)\,Y^{-\beta-2}\sin(\gamma\log Y)\,dY,\\
J(\rho)&=\sum_{k\ge0}J_k(\rho)=-\frac{\operatorname{Im}\widehat E_G(1+\rho)}\gamma=-\frac a{|\rho-1|^2}.
\end{split}
\end{equation}
The last two equalities are exact (the second is \CGeqref{eq:EGrho}); the formal replacement of $\mu$ is not a theorem and is used only in this remark. The band $k=0$ is the top block $Y\in[1,2)$ with kernel $\mathrm{Cl}_2(2\pi Y)/Y-aY$. Measured (\texttt{bands\_pairs.py}, \texttt{block\_top.py}): for $\rho_1,\rho_2,\rho_3$ the bands $k=0,1,2,3$ are $(-5.6,-66.4,+72.4,-11.3)$, $(+8.1,-29.0,+14.2,-1.1)$, $(+1.0,-9.0,-6.2,+10.3)$ in units of $10^{-4}$, and their partial sums up to $Y<4096$ are $-7.9508$, $-3.5964$, $-2.5438$ against $-a/|\rho-1|^2=-7.9541$, $-3.5984$, $-2.5426$ ($\times10^{-4}$). In the top block alone the Clausen part of the kernel has the opposite sign to the part $-aY$ at the first ten zeros and the total has both signs. The same residue $-a/((\rho-1)\zeta'(\rho))$ squared is the contribution of $(\rho,\bar\rho)$ to $-aX^2m(\sqrt X)^2$, in agreement with the exact identity $\mathcal C=2W^\flat-\mathcal R-W-aX^2m(\sqrt X)^2$ of Proposition~\CGref{prop:decomp}.
\end{remark}

For $U\ge1$ and $1\le n\le U^2$ let
\begin{equation}\label{eq:DUdef}
D_U(n)=\sum_{\substack{d,j\le U\\ dj>n}}\frac{\mu(d)\mu(j)}{dj}=m(U)^2-\sum_{\substack{d,j\le U\\ dj\le n}}\frac{\mu(d)\mu(j)}{dj}=\sum_{n<k\le U^2}\frac{\lambda(k)\,b_U(k)}k,
\end{equation}
the deficit of the partial sums of $a_U(k)/k$ with respect to their final value $m(U)^2$; the last form is Proposition~\CGref{prop:liouville}. Note $D_U(U^2)=0$.

\begin{lemma}\label{lem:Pbound}\VCGreviewed{1}
Assume RH and let $\varepsilon>0$. Then, uniformly in $U\ge1$ and $t\in\mathbb R$,
\begin{equation}\label{eq:Pbound}
\Bigl|P_U(\tfrac12+it)\Bigr|=\Bigl|\sum_{d\le U}\frac{\mu(d)}{d^{1/2+it}}\Bigr|\ll_\varepsilon U^\varepsilon(2+|t|)^\varepsilon .
\end{equation}
\end{lemma}

\begin{proof}
Let $M(y;t)=\sum_{k\le y}\mu(k)k^{-it}$. For $x\ge1$ half an odd integer, Lemma~3.12 of \cite[pp.~60--61]{TitchmarshHB1986}, applied to $f(w)=1/\zeta(w+it)=\sum_k\mu(k)k^{-it}k^{-w}$ at $w=0$ (coefficients bounded by $1$, $\sum_k k^{-\sigma}\le\sigma/(\sigma-1)$), with $c=1+1/\log x$ and $T=x^2$, gives
\begin{equation}\label{eq:Mt}
M(x;t)=\frac1{2\pi i}\int_{c-iT}^{c+iT}\frac{x^w}{\zeta(w+it)}\frac{dw}w+O\Bigl(\frac{\log x}x\Bigr).
\end{equation}
Under RH, $1/\zeta(u)\ll_\varepsilon(2+|\operatorname{Im}u|)^\varepsilon$ uniformly for $\operatorname{Re}u\ge\frac12+\varepsilon$ (\cite[Theorem~14.2, (14.2.6), p.~337]{TitchmarshHB1986} for $|\operatorname{Im}u|\ge2$, and continuity and non-vanishing on the compact remainder, $1/\zeta$ being holomorphic at $u=1$). Move the segment to $\operatorname{Re}w=\frac12+\varepsilon$; the integrand is holomorphic in between, since $w=0$ lies to the left. The new vertical segment contributes $\ll x^{1/2+\varepsilon}(2+T+|t|)^\varepsilon\log T$, the horizontal ones $\ll x^c(2+T+|t|)^\varepsilon/T$. Hence $M(x;t)\ll_\varepsilon x^{1/2+3\varepsilon}(2+|t|)^\varepsilon$, and the same for integer $y$ since $M(y;t)=M(y+\frac12;t)$. Partial summation gives $P_U(\frac12+it)=M(U;t)U^{-1/2}+\frac12\int_1^UM(y;t)y^{-3/2}dy\ll_\varepsilon U^{3\varepsilon}(2+|t|)^\varepsilon$, and $\varepsilon$ is arbitrary.
\end{proof}

\begin{theorem}[Centred primitive criterion]\label{thm:DU}\VCGreviewed{1}
The following are equivalent.
\textup{(a)} RH.
\textup{(b)} For every $\varepsilon>0$ there is $K_\varepsilon$ such that, for all $U\ge1$ and all integers $U\le n<U^2$,
\begin{equation}\label{eq:DUcrit}
|D_U(n)|\le K_\varepsilon\,\frac{U^\varepsilon}{\sqrt n}.
\end{equation}
\end{theorem}

\begin{proof}
(a)$\Rightarrow$(b). Let $U\le n<U^2$, $x=n+\frac12$, $T=n^2$, $c=1/\log x$. The Dirichlet series $\sum_ka_U(k)k^{-1-w}$ is the polynomial $P_U(1+w)^2$. Lemma~3.12 of \cite[pp.~60--61]{TitchmarshHB1986}, applied to the Dirichlet series $P_U(w)^2=\sum_ka_U(k)k^{-w}$ at the point $w=1$, with coefficients $a_U(k)$ bounded by $\tau(k)\ll_\varepsilon k^\varepsilon$ and $\sum_k|a_U(k)|k^{-\sigma}\le\zeta(\sigma)^2\ll(\sigma-1)^{-2}$ uniformly in $U$, gives
\begin{equation}\label{eq:DUperron}
\sum_{k\le n}\frac{a_U(k)}k=\frac1{2\pi i}\int_{c-iT}^{c+iT}P_U(1+w)^2\frac{x^w}w\,dw+O_\varepsilon\Bigl(\frac{\log^2x}{T}+\frac{x^\varepsilon\log x}T\Bigr).
\end{equation}
The integrand is holomorphic except at $w=0$, where the residue is $P_U(1)^2=m(U)^2$. Moving the segment to $\operatorname{Re}w=-\frac12$,
\begin{equation}\label{eq:DUshift}
D_U(n)=-\frac1{2\pi i}\int_{-\frac12-iT}^{-\frac12+iT}P_U(\tfrac12+iv)^2\frac{x^w}w\,dw-\frac1{2\pi i}\Bigl(\int_{-\frac12+iT}^{c+iT}-\int_{-\frac12-iT}^{c-iT}\Bigr)P_U(1+w)^2\frac{x^w}w\,dw+O_\varepsilon\bigl(n^{-2+\varepsilon}\bigr).
\end{equation}
By Lemma~\CGref{lem:Pbound} and $|w|\ge\max(\frac12,|v|)$, the vertical integral is $\ll_\varepsilon x^{-1/2}U^{2\varepsilon}(2+T)^{2\varepsilon}\log T\ll_\varepsilon U^{11\varepsilon}n^{-1/2}$, since $T=n^2<U^4$. On the horizontal segments $|P_U(1+w)|\le\sum_{d\le U}d^{-1/2}\le2\sqrt U$, $|x^w|\le x^c\le e$ and $|w|\ge T$, so they contribute $\ll U/n^2\le n^{-1}$, using $n\ge U$. This gives \CGeqref{eq:DUcrit} with $11\varepsilon$ in place of $\varepsilon$.

(b)$\Rightarrow$(a). Let $X\ge4$, $U=\lfloor\sqrt X\rfloor$, $N=\lfloor X\rfloor$; the corner \CGeqref{eq:Rest} is $\mathcal C(X)=\sum_{U<k\le N}c(k)\mathcal V(k)$ with $c(k)=a_U(k)/k$ and $\mathcal V(k)=k^2E_G(X/k)$, because for integers $d,j\le\sqrt X$ is $d,j\le U$ and $dj>\sqrt X$ is $dj\ge U+1$. Since $c(k)=D_U(k-1)-D_U(k)$ and $D_U(N)=0$ (as $N\ge U^2$), summation by parts gives exactly
\begin{equation}\label{eq:DUabel}
\mathcal C(X)=D_U(U)\,\mathcal V(U+1)+\sum_{U<k<U^2}D_U(k)\bigl(\mathcal V(k+1)-\mathcal V(k)\bigr).
\end{equation}
The function $E_G$ is continuous on $[1,\infty)$ and, on each interval between consecutive integers, continuously differentiable with $E_G'=-2\pi R$, which is integrable there because $F$ is locally integrable; so $\mathcal V$ is absolutely continuous on $[U+1,N]$ with
\begin{equation}\label{eq:hprime}
\mathcal V'(\xi)=2\xi E_G(X/\xi)+2\pi X\,R(X/\xi)\qquad\text{for almost every }\xi,
\end{equation}
by Lemma~\CGref{lem:EG}(iii), and $|\mathcal V(k+1)-\mathcal V(k)|\le\int_k^{k+1}|\mathcal V'|$. By \CGeqref{eq:DUcrit} and $\xi\le k+1\le2k$, $|D_U(k)|\le\sqrt2K_\varepsilon U^\varepsilon\xi^{-1/2}$ on $[k,k+1]$. Fix $\delta\in(0,\frac12)$ and let $c_\delta$ be the constant of Lemma~\CGref{lem:EG}(ii). Then
\begin{equation}\label{eq:DUestimate}
|\mathcal C(X)|\le K_\varepsilon U^\varepsilon\Bigl(U^{-1/2}|\mathcal V(U+1)|+\sqrt2\int_U^X\xi^{-1/2}\bigl(2\xi|E_G(X/\xi)|+2\pi X|R(X/\xi)|\bigr)d\xi\Bigr).
\end{equation}
First, $|\mathcal V(U+1)|\le4U^2c_\delta(X/(U+1))^{1+\delta}\le4c_\delta X^{3/2+\delta/2}$, so the boundary term is $\ll X^{5/4+\delta/2}$. Second, $\int_U^X\xi^{1/2}c_\delta(X/\xi)^{1+\delta}d\xi\le c_\delta X^{1+\delta}X^{1/2-\delta}/(\frac12-\delta)\ll X^{3/2}$. Third, with $Y=X/\xi$, $\int_U^X\xi^{-1/2}X|R(X/\xi)|d\xi=X^{3/2}\int_1^{X/U}|R(Y)|Y^{-3/2}dY$ and $X/U\le2\sqrt X$. Fix $\sigma_0=\frac12+\varepsilon$; by Lemma~\CGref{lem:Rweight}(ii), $\int_{2^q}^{2^{q+1}}R^2\le(2^{q+1})^{2\sigma_0+1}\int_1^\infty R^2Y^{-2\sigma_0-1}dY\ll_\varepsilon2^{q(2+2\varepsilon)}$, so by Cauchy--Schwarz $\int_{2^q}^{2^{q+1}}|R|Y^{-3/2}\ll_\varepsilon2^{-3q/2}2^{q/2}2^{q(1+\varepsilon)}=2^{q\varepsilon}$, and summing over $0\le q\le\log_2(2\sqrt X)$ gives $\ll_\varepsilon X^{\varepsilon}\log X$. Altogether $|\mathcal C(X)|\ll_\varepsilon X^{3/2+3\varepsilon}$ for $X\ge4$, and trivially for $1\le X<4$. Since $\varepsilon$ is arbitrary, Theorem~\CGref{thm:corner}(c)$\Rightarrow$(a) gives RH.
\end{proof}

\begin{remark}\label{rem:DUnum}\VCGreviewed{1}
Two comments on \CGeqref{eq:DUcrit}. First, bounding the terms of $D_U(n)$ one by one does not give it even under RH: for $U\le n<U^2$, $D_U(n)=m(U)^2-\sum_{k\le n}(\mu*\mu)(k)/k+2\sum_{j<n/U}\mu(j)j^{-1}(m(n/j)-m(U))$, and inserting $m(y)\ll y^{-1/2+\varepsilon}$ termwise gives only $U^{-1/2+\varepsilon}$; the proof above goes through the polynomial $P_U$ on the critical line instead. Second, \CGeqref{eq:DUcrit} is a statement about a weighted mean of $\lambda$ (last form of \CGeqref{eq:DUdef}). Writing $j=d+h$ turns $m(U)^2$ into a sum of two-point correlations over all shifts $|h|<U$,
\begin{equation}\label{eq:shifts}
m(U)^2=\sum_{|h|<U}\ \sum_{\substack{d\le U\\1\le d+h\le U}}\frac{\mu(d)\mu(d+h)}{d(d+h)},
\end{equation}
so bounds for two-point correlations at fixed shifts, without uniformity up to $|h|\sim U$ and without a power saving, do not give \CGeqref{eq:DUcrit}. Measured (\texttt{DU\_test.py}): $\max_{1\le n<U^2}\sqrt n\,|D_U(n)|=2.03$, $2.47$, $3.11$, $3.39$ for $U=100$, $300$, $1000$, $2000$, attained near $n=U$ ($n/U=0.96$, $0.66$, $0.59$, $1.68$), with medians $0.12$, $0.11$, $0.05$, $0.13$.
\end{remark}

\section{Attempts on the corner and where they stopped}\label{sec:attempts}

This section records the methods tried on \CGeqref{eq:target}, with the measured quantity at which each stopped. None of them is claimed to be impossible; they were not carried through.

Termwise bounds. Lemma~\CGref{lem:EG}(ii) gives $|\mathcal C(X)|\le\sum_{d,j}dj|E_G(X/(dj))|$, which loses all cancellation from $\mu(d)\mu(j)$. With the measured size $|E_G(Y)|\approx Y^{3/4}$ the sum of absolute values of the terms of $\mathcal C$ is $26.0$, $53.5$, $95.5$ in units of $X^{3/2}$ at $X=3\cdot10^4$, $1.2\cdot10^5$, $3.8\cdot10^5$, while $|\mathcal C|$ itself is $0.13$, $0.08$, $0.11$ in the same units (Table~\CGref{tab:corner}). In units of $X^{3/2}$ the sum of absolute values grows like $X^{1/2}$: from $X=3\cdot10^4$ to $1.2\cdot10^5$ it doubles, and from $1.2\cdot10^5$ to $3.8\cdot10^5$ it grows by $1.78\approx(3.8/1.2)^{1/2}$.

Induction through $E_G'=-2\pi R$. Summing the singular parts of $R$ at the integers (Remark~\CGref{rem:corner}) over the pairs $(d,j)$ was attempted. In the attempts made, the sum over the corner did not show cancellation, and no input on the signs $\mu(d)\mu(j)$ was found that controls it.

Quadratic certificates. By Remark~\CGref{rem:barrier}, a positive semidefinite certificate that uses only the relation $\mu(d)\mu(d')=\mu(dd'/(d,d')^2)$ applies equally to $\mu^2$; for $W_{\mu^2}$ the sum is of order $X^2$ (proved, Proposition~\CGref{prop:musq}), and for the $\mu^2$-corner it is measured to be of order $X^2$. Measured: the part of the corner with $X/(dj)\le64$, with $|\mu(d)\mu(j)|$ in place of $\mu(d)\mu(j)$, is $-0.00090X^2$, $-0.00121X^2$, $-0.00148X^2$ at $X=30000.5$, $120000.5$, $480000.5$; the same part of $\mathcal C(X)$ itself is $+0.051X^{3/2}$, $+0.108X^{3/2}$, $-0.054X^{3/2}$ at these $X$. The first two agree with the sum of the bands $k\le5$ of Table~\CGref{tab:bands} (which lists $-\mathcal C$); $X=480000.5$ is not in the tables. Writing $W_\chi$ for \CGeqref{eq:Wdef} with $\mu(n)$ replaced by $\mu(n)^2\prod_{p\mid n}\chi(p)$, and fixing $\chi(p)=-1$ for $p\le\sqrt X$ and choosing the signs of the larger primes freely still allows $W_\chi(X)$ up to $0.055X^2$, $0.044X^2$, $0.042X^2$ at $X=10^4$, $10^5$, $2\cdot10^5$.

Parity. At $X=1.2\cdot10^5$ let $S_{\mathrm{even}}$ and $S_{\mathrm{odd}}$ be the sums of $dj\,E_G(X/(dj))$, without the sign $\mu(d)\mu(j)$, over the squarefree pairs $d,j\le\sqrt X$ with $X/(dj)\le64$ and $\omega(d)+\omega(j)$ even, respectively odd, where $\omega(n)$ is the number of prime factors of $n$. Then $S_{\mathrm{even}}=-6.459\cdot10^6$ and $S_{\mathrm{odd}}=-1.0958\cdot10^7$, with absolute sums $1.082\cdot10^9$ and $1.066\cdot10^9$, and the corresponding part of the corner is $S_{\mathrm{even}}-S_{\mathrm{odd}}=4.499\cdot10^6=0.108X^{3/2}$.

Dyadic bands. Splitting $\mathcal C$ by $X/(dj)\in[2^k,2^{k+1})$, single bands reach $0.73X^{3/2}$ in absolute value, larger than the total, and neighbouring bands $k=1,2$ have opposite signs in all three cases (Table~\CGref{tab:bands}).

\section{Numerical observations}\label{sec:num}

All computations are in double precision unless stated, with exact integer M\"obius values; scripts and outputs are listed in Appendix~\CGref{app:code}.

\begin{table}[ht]
\centering
\small
\caption{The exact decomposition \CGeqref{eq:decomp}, in units of $X^{3/2}$; in parentheses the sum of the absolute values of the terms. The last column is the residual of \CGeqref{eq:decomp}.}\label{tab:corner}
\begin{tabular}{lrrrrrr}
\toprule
$X$ & $W$ & $2W^\flat$ & $-aX^2m^2$ & $-\mathcal R$ & $-\mathcal C$ & residual \\
\midrule
$30000.5$ & $-0.1579$ & $-0.0004$ & $-0.0006$ & $-0.0245\ (0.442)$ & $-0.1323\ (26.0)$ & $3\cdot10^{-8}$ \\
$120000.5$ & $-0.0948$ & $-0.0000$ & $-0.0112$ & $-0.0054\ (0.421)$ & $-0.0783\ (53.5)$ & $3\cdot10^{-7}$ \\
$380000.5$ & $-0.1052$ & $+0.0001$ & $-0.0000$ & $+0.0091\ (0.359)$ & $-0.1142\ (95.5)$ & $1\cdot10^{-6}$ \\
\bottomrule
\end{tabular}
\end{table}

\begin{table}[ht]
\centering
\small
\caption{$-\mathcal C$ split into bands $X/(dj)\in[2^k,2^{k+1})$, in units of $X^{3/2}$; in parentheses the sum of absolute values.}\label{tab:bands}
\begin{tabular}{lrrrrrr}
\toprule
$X$ & $k=0$ & $k=1$ & $k=2$ & $k=3$ & $k=4$ & $k=5$ \\
\midrule
$30000.5$ & $+0.037\ (4.8)$ & $+0.273\ (7.0)$ & $-0.434\ (6.8)$ & $-0.065\ (3.3)$ & $-0.005\ (2.1)$ & $+0.143\ (1.1)$ \\
$120000.5$ & $+0.206\ (10.3)$ & $+0.309\ (14.3)$ & $-0.725\ (13.8)$ & $+0.250\ (6.7)$ & $-0.116\ (4.3)$ & $-0.032\ (2.2)$ \\
$380000.5$ & $-0.095\ (18.7)$ & $-0.252\ (25.6)$ & $+0.547\ (24.5)$ & $-0.249\ (11.8)$ & $-0.103\ (7.7)$ & $+0.080\ (3.8)$ \\
\bottomrule
\end{tabular}
\end{table}

\begin{table}[ht]
\centering
\small
\caption{Residues $r_k=r_{\rho_k}$ of $\widehat W$ at $s=1+\rho_k$ for the first five zeros (mpmath, $30$ digits).}\label{tab:res}
\begin{tabular}{rrrrr}
\toprule
$k$ & $\gamma_k$ & $|A(\rho_k)|$ & $|\zeta'(\rho_k)|$ & $|r_k|$ \\
\midrule
$1$ & $14.1347251417$ & $0.147337$ & $0.79316$ & $0.0821128$ \\
$2$ & $21.0220396388$ & $0.112785$ & $1.13684$ & $0.0295771$ \\
$3$ & $25.0108575801$ & $0.100191$ & $1.37172$ & $0.0183161$ \\
$4$ & $30.4248761259$ & $0.0875278$ & $1.30394$ & $0.0138456$ \\
$5$ & $32.9350615877$ & $0.0828281$ & $1.38212$ & $0.011421$ \\
\bottomrule
\end{tabular}
\end{table}

Size of $E_G$. Over $20$ random $Y$ in $[10^q,2\cdot10^q]$, the root mean square of $E_G(Y)/Y^{3/4}$ is $0.948$, $1.119$, $0.676$ for $q=3,4,5$. Over $4000$ values of $Y$ drawn log-uniformly in $[10,2\cdot10^5]$, $\max|E_G(Y)|/Y^{3/4}=2.830$, and over the first $400$ of them $\max|E_G(Y)|/Y=1.025$. The identity of Lemma~\CGref{lem:EG}(iii) was checked by numerical differentiation at $Y=37.3$, $512.77$, $4096.41$ (agreement to $10^{-5}$).

Size of $W$ and the zeros. Over $60$ random $X$ in $[10^q,2\cdot10^q]$, $q=3,4,5$: $\max|W|/X^{3/2}=0.2335$, $0.2827$, $0.2346$; root mean square of $W/X^{3/2}$: $0.133$, $0.127$, $0.141$. The correlation of $W(X)$ with $2\operatorname{Re}\sum_{k\le100}r_kX^{1+\rho_k}$ is $0.99985$, $0.99993$, $0.99986$, with regression slopes $0.9975$, $0.9953$, $0.9976$. For $H$, over $150$ random $x$ in $[10^q,2\cdot10^q]$, $q=3,\dots,6$, the root mean square of $H/\sqrt x$ is $0.399$, $0.417$, $0.413$, $0.475$, and the correlation with $2\operatorname{Re}\sum_{k\le150}A(\rho_k)x^{\rho_k}/\zeta'(\rho_k)$ is $0.937$, $0.935$, $0.930$, $0.951$.

The no-go constant. $W_{\mu^2}(X)/X^2=0.0977$, $0.0975$, $0.0968$ at $X=1000.5$, $10000.5$, $100000.5$, against $(6/\pi^2)a=0.09673$ from \CGeqref{eq:musq}.

Identities. \CGeqref{eq:Hsinc} was checked at $x=2.37$, $17.61$ (agreement to $10^{-8}$; at $x=1234.57$ and $98765.43$ the truncation $d\le4\cdot10^6$ of the series dominates the difference); \CGeqref{eq:Hcomplex} to $2\cdot10^{-11}$; \CGeqref{eq:Hplus} for $2\le n<400$, evaluating $H$ at distance $10^{-9}$ to the right of $n$ (difference at most $1.2\cdot10^{-4}$); Proposition~\CGref{prop:Wbasic}(i) at $X=57.3$, $300.7$ (agreement to $10^{-6}$) and (ii) at one value of $X$ to $20$ digits; Proposition~\CGref{prop:d3} at $X=20000.5$, $100000.5$ (residual $<10^{-5}X^{3/2}$).

\section{What the route reduces the Riemann hypothesis to}\label{sec:remaining}

By Theorem~\CGref{thm:corner}, RH is equivalent to the following inequality, which is not proved: for every $\varepsilon>0$ there is $K_\varepsilon$ such that
\begin{equation}\label{eq:target}
\sum_{\substack{d,j\le\sqrt X\\ dj>\sqrt X}}\mu(d)\mu(j)\,dj\,E_G\Bigl(\frac X{dj}\Bigr)\ge-K_\varepsilon X^{3/2+\varepsilon}\qquad(X\ge1),
\end{equation}
where $E_G(Y)=\sum_{\ell\le Y}\ell\,\mathrm{Cl}_2(2\pi Y/\ell)-\pi(C-\log2)Y^2$. Only $\mu$ on $[1,\sqrt X]$ and $E_G$ on $[1,\sqrt X)$ enter. Inequality \CGeqref{eq:target} is one-sided, and the term $-aX^2m(\sqrt X)^2$ of \CGeqref{eq:decomp} plays no role in it. This document does not prove it. By Remark~\CGref{rem:barrier}, a proof of it must use more about $\mu$ than the multiplicativity relation on squarefree integers.

By Proposition~\CGref{prop:liouville}, \CGeqref{eq:target} reads $\sum_{\sqrt X<n\le X}\lambda(n)b_{\sqrt X}(n)\,nE_G(X/n)\ge-K_\varepsilon X^{3/2+\varepsilon}$, a one-sided bound for a Liouville sum with the explicit nonnegative weight $b_{\sqrt X}$, and by Proposition~\CGref{prop:clausenmuntz} it is $X$ times a weighted Liouville sum of dilations of the M\"untz transform of the Clausen kernel. By Theorem~\CGref{thm:DU}, RH is also equivalent to the bound \CGeqref{eq:DUcrit} for the centred primitive $D_U(n)=\sum_{n<k\le U^2}\lambda(k)b_U(k)/k$, which is not proved either. A further unproved statement, \CGeqref{eq:target2}, is sufficient for RH but not known to be equivalent to it. It is a bound for a positive semidefinite quadratic form in $(\mu(d))_{d\le N}$ whose matrix \CGeqref{eq:gram} is explicit, depends on $d,e$ only through $(de)^{-\sigma}$ and $\log(e/d)$, and has a diagonal bounded uniformly in $N$.

\begin{thebibliography}{9}

\bibitem{BaezDuarte2005} L. B\'aez-Duarte, A general strong Nyman--Beurling criterion for the Riemann hypothesis, arXiv:math/0505453, version 1 (2005), \url{https://arxiv.org/abs/math/0505453} (only its abstract is used here).

\bibitem{Davenport1980} H. Davenport, \emph{Multiplicative Number Theory}, second edition, revised by H.~L. Montgomery, Graduate Texts in Mathematics 74, Springer, New York, 1980, \href{https://doi.org/10.1007/978-1-4757-5927-3}{doi:10.1007/978-1-4757-5927-3}.

\bibitem{Rudin1987} W. Rudin, \emph{Real and Complex Analysis}, third edition, McGraw-Hill, New York, 1987, ISBN 0-07-054234-1 (no DOI).

\bibitem{Titchmarsh1939} E.~C. Titchmarsh, \emph{The Theory of Functions}, second edition, Oxford University Press, Oxford, 1939 (no DOI).

\bibitem{TitchmarshHB1986} E.~C. Titchmarsh, \emph{The Theory of the Riemann Zeta-Function}, second edition, revised by D.~R. Heath-Brown, Clarendon Press, Oxford, 1986 (no DOI).

\end{thebibliography}

\clearpage
\appendix
\clearpage
\section{Code and data}\label{app:code}

The zip file that accompanies this document contains every script with its output (\texttt{X.py} and \texttt{X\_output.txt}, and a copy \texttt{X.py.txt}). The constants of Sections~\CGref{sec:mellin}--\CGref{sec:W} and Table~\CGref{tab:res} come from \texttt{route\_constants.py}. Table~\CGref{tab:corner} and the size of $E_G$ come from \texttt{corner\_exact.py} and \texttt{bilinear\_W.py}; Table~\CGref{tab:bands} from \texttt{corner\_bands.py}; parity, $E_G'=-2\pi R$ and the restricted corner from \texttt{corner\_parity.py} and \texttt{corner\_scale.py}; $W$ against the zeros from \texttt{mu\_clausen.py}; the no-go constant and the $\chi$ experiments from \texttt{nogo\_chi.py}; Proposition~\CGref{prop:lcm} from \texttt{lcm\_form.py} and \texttt{redheffer\_W.py}; Proposition~\CGref{prop:d3} from \texttt{huxley\_watt3.py}; \CGeqref{eq:WMK} from \texttt{W\_kernel2.py}; Proposition~\CGref{prop:Hid} and \CGeqref{eq:Hplus} from \texttt{telescoping.py}, \texttt{idcheck.py} and \texttt{left\_sing.py}; the explicit formula for $H$ from \texttt{H\_explicit.py}; Section~\CGref{sec:liouville} from \texttt{liouville\_corner.py}, \texttt{liouville\_weight.py}, \texttt{corner\_onesided.py} and \texttt{musq\_corner\_const.py}; Section~\CGref{sec:centred} from \texttt{bands\_pairs.py}, \texttt{block\_top.py} and \texttt{DU\_test.py}, and the profile of the Abel form from \texttt{corner\_cancel.py} and \texttt{corner\_cancel2.py}. The remaining scripts of the folder \texttt{pairs} belong to earlier exploration of the same functions and are included unchanged.

\end{document}

# BEGIN CORRECTVCG  (CorrectVCG checker, ignored by LaTeX; extract with the sed line at the top)
#!/usr/bin/env python3
r"""CorrectVCG checker (version 10) for Doktor's manuscripts and for Theorems and Proofs (TP).

Every .tex delivered to him carries a copy of this script after \end{document} (LaTeX never reads it).
Extract it with:   sed -n '/^# BEGIN CORRECTVCG/,/^# END CORRECTVCG/p' FILE.tex > correctvcg_check.py
Read every rule first:   python3 correctvcg_check.py --rules
Main uses:
    python3 correctvcg_check.py --paper FILE.tex                      format + mathematical checks of a manuscript
    python3 correctvcg_check.py CHAPTER.tex --chapter N --previous OLD  TP chapter
    python3 correctvcg_check.py --paper NEW.tex --diff OLD.tex --allowed labels.txt   no unrequested change
    python3 correctvcg_check.py FILE.tex --mark OUT.tex               review marks \VCGreviewed{n}
    python3 correctvcg_check.py FILE.tex --mark OUT.tex --reviewer codex     Codex marks \VCGreviewedCodex{n}
    python3 correctvcg_check.py FILE.tex --mark OUT.tex --reviewer chatgpt   ChatGPT marks \VCGreviewedChatGPT{n}
    python3 correctvcg_check.py FILE.tex --mark OUT.tex --reviewer copilot   Copilot marks \VCGreviewedCopilot{n}
    python3 correctvcg_check.py FILE.tex --embed OUT.tex              put this checker into a .tex
Exit code 0 only if there are no ERROR lines.  WARN lines must be read and fixed or justified one by one.
"""
import re
import sys
import argparse
import os
from collections import Counter

RULES = r'''
CorrectVCG RULES (v5). Every correction Doktor (Victor Casana-Giner) has had to make, mathematical and formal.
They apply to ALL his papers, notes, reviews and to every chapter of the book Theorems and Proofs (TP), not to one paper.
Read all of them before touching any of his .tex files. The past cases behind each rule are listed (in Spanish) in the
project docs claude/CorrectVCG-errores-matematicos-v4.md, claude/CorrectVCG-formato-v2.md, claude/CorrectVCG-listado-correcciones-v1.md.
This block and the checker below travel inside every .tex after \end{document}. NEVER remove them automatically:
they are removed only when Doktor himself asks for it (final versions).

RULE ZERO. Every error in what I deliver is mine, wherever it came from (source, other author, other AI, one of his drafts).
"It came from the source" is never an excuse.

THE NINE PATTERNS I REPEAT (check each one, in writing, before delivering anything):
 1. Marking a point as corrected without looking at the final output.
 2. Not running through the boundary cases: 0, negative values, 1, the first value of every index, empty ranges, degenerate parameters, the sign when dividing.
 3. Leaving symbols undefined, without domain, used before being defined, with two meanings, or renamed in only part of the text.
 4. Writing "equivalent", "suffices", "necessary", "unconditional" without both directions or the quantitative form.
 5. Leaving hard steps as "the same argument", "applies verbatim", "similarly", "as in [ref]".
 6. Marking a statement as verified without redoing the computations.
 7. Ignoring corrections he already gave: before delivering, run through all the accumulated rules.
 8. Fixing only the first occurrence of an error (first paragraph or section) and not the others.
 9. Importing results and errors from his papers and drafts without re-proving them, or citing him in the bibliography.
 And the one that annoys him most: in a revision, changing things that were right and breaking them. Only what he asks, or a proven error, changes; everything else stays letter by letter (checker: --diff OLD --allowed labels.txt).

ONE COMPLETE REVIEW SHOULD SUFFICE. Every review is complete, as if it were the only one: every statement, the whole checklist
below on each, every occurrence of every error, checker with 0 ERROR. In practice there will be several; none relaxes anything.
Every reviewed statement carries the reviewer's own mark in the LaTeX (n = number of reviews of that statement by that reviewer;
checker --mark --reviewer R): Claude \VCGreviewed{n}, Codex \VCGreviewedCodex{n}, ChatGPT \VCGreviewedChatGPT{n},
Copilot \VCGreviewedCopilot{n}.
Each reviewer touches only its own marks; the marks, counters and preamble block of the other reviewers stay letter by letter
(checker --diff gives ERROR otherwise). Review marks are not text changes and are ignored by every content check and by the
unrequested-change check. The embedded checker is never replaced by an older version (--embed refuses), and never removed unless
Doktor asks.

A. WHAT IS ASKED, AND APPLYING A CORRECTION REPORT
 A1. Deliver exactly the object asked (a chapter of TP is a chapter of TP, not a new version of its source paper; "only that part" means nothing else). A new version of a paper or of TP only when he asks.
 A2. If object or scope is unclear, ask ONCE before generating. Never regenerate "just in case".
 A3. For every report: point -> cause (mathematics / generator / text) -> where it is fixed -> LITERAL QUOTE of the corrected text in the FINAL OUTPUT. A point without that quote is not corrected.
 A4. Every correction in all its occurrences: grep the error and its variants (same symbol, formula, phrase, copies in settings, tooltips, abstract, introduction, other sections, other chapters, the checker, the skill); count before, 0 after. A general correction becomes a global convention, and every place that depended on it is revisited.
 A5. Never mark corrections as applied in block; each one is checked in the output.
 A6. Regression: in every new version re-check ALL points of ALL previous rounds.
 A7. If a report flags some cases of a pattern, search the pattern everywhere: there are usually more.
 A8. Mathematical errors are serious, never "minor". Referees' style recommendations are not applied. References cite the version actually checked (arXiv stays arXiv).

B. MATHEMATICAL CHECKLIST (each statement and proof, in writing; verified only after redoing every equality and boundary case)
 B1. Domains: every parameter, variable and index has its domain (N, Z, R, >=1, coprime, prime...) declared before use, in the statement or its setting.
 B2. First value of every index: substitute it into every sum, product, recurrence and induction; denominators, arithmetic functions and logarithms must be defined there (arguments >=1, nothing divided by 0, log of positive numbers). Counting from 0 or from 1 changes counts by one.
 B3. Division by zero: no denominator that can vanish for an allowed value of a parameter; the excluded value is treated separately.
 B4. Sign when dividing or multiplying an inequality by an expression whose sign depends on a parameter: each sign case separately. Excluded terms and the O(1) boundary values stated exactly.
 B5. Exact versus asymptotic: every "=" is exact; asymptotic steps carry their +o()/+O(); every intermediate equality redone.
 B6. Numerator and denominator: 1/min = max of reciprocals and conversely; a ratio defined as a/b never reappears as b/a; min/max read literally.
 B7. Exponents: braces (x^{-1}, never x^-1); x^{1/2+eps} not x^{1/2}+eps; the exponent of the statement equals the one at the end of the proof; the line with the total exponent written.
 B8. Floors and ceilings, ranges of blocks and windows; constants with all their dependencies; no loose bound passed off as sharp.
 B9. Symbols: every symbol and named object (norms, operators, transforms and their normalisation, coefficients, hypotheses with names) defined, with domain, before its first use; every new term defined or removed; quantifiers explicit. A letter that occurs once, in a formula of a proof, and nowhere else in the result, its proof or its Setting is undefined (usually a leftover of a renaming: o(log K) after K was renamed Y).
 B10. Collisions: one symbol, one meaning in the whole text (a fixed parameter is never also a summation index; a letter is never two different constants).
 B11. Renaming: a renamed symbol is renamed everywhere; never mix epsilon/varepsilon, phi/varphi, rho/varrho, theta/vartheta for one object, nor the same named hypothesis with different parameter letters.
 B12. Scope: "equivalent", "iff", "necessary", "suffices", "unconditional", "the theorem holds": both directions, the exact quantitative form used, the exact parameter range proved; a property of a sum is not a property of each summand; write the precise conclusion, not "the theorem holds".
 B13. Logical status: conditional on a named conjecture / classical unconditional / unconditional only through an external reference (say if it is an unrefereed preprint) / conditional on named hypotheses; every result in the right class; say what depends on which reference. A title that says "unconditional" may not hide a dependence on a cited result (proved only in a reference, e.g. an unrefereed preprint): the title says "unconditional through Theorem N" or the result goes to the conditional class.
 B14. An example satisfying a numerical inequality does not satisfy the structural hypotheses (multiplicativity, independence, ...) of the theorem it is compared with.
 B15. Nothing delegated: no "the same argument", "applies verbatim", "similarly", "as in [ref]", "standard" without statement and page at a load-bearing step. Write the intermediate lemma with every substitution, factor, congruence, Fourier inversion, the check that the hypotheses still hold, and every combinatorial fact used.
 B16. His own work: never cited as bibliography (no bibitem or \cite of Casana-Giner, Casanya, LawCasanya, Rotkiv Aynasac); his name appears only as author of the document itself. Nothing rests on his papers, drafts or compendium: whatever is taken from them is re-proved from scratch (drafts may be wrong; a draft is no authority).
 B17. Literature checked in LIBRARY TOTAL (GitHub repository vcasana1/LIBRARY-TOTAL-28oct2026, always accessible: git clone --filter=blob:none --no-checkout, then check out only INDEX.md and the PDFs needed; never say it is not reachable) (statement, hypotheses, page) ONCE: recorded at once in the registry claude/<paper>-fuentes-verificadas.md and never reopened (checker --verified). After compiling, look at the bibliography pages as PNG (checker --bib-pdf): DOI or arXiv in each entry, "(Cited on p. ...)", no [?] or ??.
 B18. When a statement passed B1-B17 with nothing found, go through B1-B4 and B12 again: the errors that escaped were there.

C. FORMAT OF EVERY .tex (canonical preamble of excursion_area_section_v39, in memory topics/latex-pdf-config)
 C1. Every displayed equation numbered on the right with \label (equation/align/gather/multline); never \[ \], $$ or starred environments for mathematics; \numberwithin{equation}{section}; amsart with reqno.
 C2. Links: v39 configuration (colorlinks=false, boxes: internal red, citations blue, URL/DOI teal; never hidelinks). All references with \CGref/\CGeqref, never \ref/\eqref; each \label with its CGTip tooltip inline in the preamble (no separate tips file). TOC. Every section with \label.
 C3. Bibliography: manual thebibliography, AuthorYear keys, DOI in every entry unless none exists (arXiv when appropriate), back-references "(Cited on p./pp. ...)".
 C4. Author: \author{Rotkiv Ayn\"asac} (never "Rotkin Aynasac"), no affiliation, e-mail rotkivanyasac@gmail.com below; or Victor Casana-Giner. Both are valid.
 C5. Prose: no itemize/enumerate in a paper; more full stops than semicolons; long paragraphs split. No "Questions"/"Open problems" sections; no heuristics (prove it or leave it out); no paper history (versions, retractions, referees); no "Provenance" paragraphs; no non-standard tags ([N], [C], [LIT]); explicit AMS status.
 C6. Words and symbols: clock, ordinate(s), envelope forbidden (also in labels, keys, captions); mu only Moebius; pi the prime-counting function and pi_RvM his; tau^c_k; no symbol with two meanings.
 C7. Appendices: \clearpage before \appendix and before each appendix.
 C8. Figures: PNG (never PDF images), never lower resolution; zip with figures at the root, Python code and its output.
 C9. Compile with pdflatex several passes: 0 errors, 0 undefined references, 0 multiply defined; a missing package is installed, never removed; the preamble is never changed to compile.
 C10. His papers: change only what he asks (words, commas, format untouched); propose exact lines and wait for his order; work on a copy and deliver the diff. A new version only when he asks.
 C11. Deliveries: .tex + .txt (literal copy) + .pdf as loose files and a zip that compiles; new integer version number in every file name; LaTeX and code in English; code in Python with its .txt.

D. THE BOOK THEOREMS AND PROOFS (TP)
 D1. Single edition: no edition, date, version or "supplement" marks.
 D2. Hyperlinks to everything. Literature in blue with [Literature]; background chapter of literature theorems.
 D3. Every entry has SettingBlock or SettingAsIn with the paragraphs defining its symbols and the displays it references (every \ref, \eqref, \CGref, \CGeqref, \autoref, \cref variant); copies identical to the original; settings never truncated.
 D4. Equation references never resolved to the id of the entry containing them; no self references; equation numbers = real chapter number, no gaps, every quoted number displayed.
 D5. No generator placeholder ("PROOF NOT LOCATED..."); a cited result carries "cited result" status with its reference. No corrupted text; umlauts as TeX accents.
 D6. Permanent ids (bookbuild/id_registry.json): never changed or reused; new = max+1.
 D7. Release with bookbuild/release.sh: VERSION, README version table, AI_GUIDE (where the database is; automatic publication of every version); old atlas/INDEX/SQLite deleted; id_registry.json and reviews/ kept. Say which version and commit were published. "sube git" means everything pending. Process improvements (checker, rules, docs) always go to correctpool at once, with no order needed; papers and drafts are never put on GitHub without his order; TP is released when he asks for a new version.

E. HONESTY IN THE ANSWER
 E1. Classical results are called classical from the first sentence, with their author; never disguised as new.
 E2. Work inside his programme and his own theorems; literature from LIBRARY TOTAL, not the web.
 E3. Never "corrected" without having seen it in the output. Report in each answer: the table of A3, the number of changed blocks (--diff) and that all were authorised, the number of reviewed statements and their review number, and what is still pending.
 E4. Chat in Spanish, formulas in LaTeX display, no "Provenance" or meta paragraphs at the end.
'''

KINDS = r'(?:Teorema|Lema|Proposici\\?\'?[oó]n|Corolario|Conjetura|Observaci\\?\'?[oó]n|Definici\\?\'?[oó]n|Theorem|Lemma|Proposition|Corollary|Conjecture|Remark|Definition|Hypothesis|Example|Condition|Assumption|Step|Part|Claim)'
FORBIDDEN_WORDS = [r'\bclock', r'\bordinates?\b', r'\benvelope']
OWN_NAMES = [r'Casa(?:\\~n|ñ|n|ny)a', r'Ayn(?:\\"a|ä|a)sac', r'Rotkiv']
LEFTOVERS = [
    (r'PROOF NOT LOCATED', 'generator placeholder left in the book'),
    (r'applies verbatim|verbatim the same|carries over verbatim|the same expansion leads to', 'proof step delegated instead of written out'),
    (r'H\s{2,}lder|H\?lder|H\ufffdlder', 'corrupted Hölder'),
    (r'\ufffd', 'mojibake character U+FFFD'),
    (r'[äöüÄÖÜ]', 'raw umlaut: use TeX accents (\\"o) in TP sources'),
]
EDITION_MARKS = [r'\bSupplement\b', r'\bsupplement(?:ary)? (?:to|of) (?:the )?(?:book|edition)',
                 r'\b(?:new|previous|this|second|first) edition\b', r'\bin version v?\d', r'(?<![./\w])v\d+\b(?![.](?:tex|n\d))']


def _braced(depth):
    """Regex for text with balanced braces nested up to depth levels (titles may contain \\ref{...}, \\textnormal{[\\ResolvedLink{..}{..}]})."""
    p = r'[^{}]*'
    for _ in range(depth):
        p = r'(?:[^{}]|\{' + p + r'\})*'
    return p


BRACED = _braced(6)


def entries(text):
    """Split into (uid, title, body) blocks by \\ResultHeading."""
    parts = re.split(r'(?=\\ResultHeading\{)', text)
    out = []
    for p in parts:
        m = re.match(r'\\ResultHeading\{(U\d{5})\}\{(' + BRACED + r')\}', p)
        if m:
            out.append((m.group(1), m.group(2), p))
    return out


def setting_spans(text):
    """Character spans of the arguments of \\SettingBlock{...} (balanced braces)."""
    spans = []
    for m in re.finditer(r'\\SettingBlock\{', text):
        depth, i = 1, m.end()
        while i < len(text) and depth:
            c = text[i]
            if c == '\\':
                i += 2
                continue
            depth += (c == '{') - (c == '}')
            i += 1
        spans.append((m.end(), i))
    return spans


PROOF_MARK = r'\\ProofHeading|\\begin\{proof\}'


def tp_statement(b, sp):
    """Statement of a TP entry: after its Setting (or after its heading notes) and before its proof.
    The note "PROOF INCLUDED" sits in the heading, so it never marks where the proof starts."""
    if sp:
        st = b[sp[-1][1]:]
    else:
        m = re.search(r'\\SourceNote\{[^{}]*(?:\{[^{}]*\}[^{}]*)*\}', b)
        st = b[m.end():] if m else b
    return re.split(PROOF_MARK, st, maxsplit=1)[0]


def setting_context(uid, byid, seen=None):
    """Everything a TP entry may rely on: its own Setting and, through SettingAsIn, the whole chain of earlier entries."""
    seen = set() if seen is None else seen
    if uid in seen or uid not in byid:
        return ''
    seen.add(uid)
    b = byid[uid]
    out = ''.join(b[x:y] for x, y in setting_spans(b))
    for r in re.findall(r'\\SettingAsIn\{(?:result:)?(U\d{5})\}', b):
        out += byid.get(r, '') + setting_context(r, byid, seen)
    return out


def balance(s):
    """Return a list of problems with braces and environments in s."""
    probs = []
    s2 = re.sub(r'\\[{}]', '', s)
    s2 = re.sub(r'(?<!\\)%.*', '', s2)
    depth = 0
    for ch in s2:
        if ch == '{':
            depth += 1
        elif ch == '}':
            depth -= 1
            if depth < 0:
                probs.append('closing brace without opening brace')
                depth = 0
    if depth:
        probs.append('%d unclosed brace(s)' % depth)
    stack = []
    for m in re.finditer(r'\\(begin|end)\{([^}]+)\}', s2):
        if m.group(1) == 'begin':
            stack.append(m.group(2))
        elif not stack or stack[-1] != m.group(2):
            probs.append('\\end{%s} does not match %s' % (m.group(2), stack[-1] if stack else 'nothing'))
        else:
            stack.pop()
    if stack:
        probs.append('unclosed environment(s): ' + ', '.join(stack))
    return probs


def check_paper(text, errors, warns):
    """Format rules for Doktor's manuscripts (LatexRules + excursion_area_section_v39 configuration)."""
    body = text.split('\\begin{document}', 1)[-1]
    pre = text.split('\\begin{document}', 1)[0]
    ctx = lambda m, t=body: t[max(0, m.start() - 50):m.end() + 50].replace('\n', ' ')
    # Every display numbered on the right and labelled.
    for m in re.finditer(r'(?<!\\)\\\[|\$\$', body):
        errors.append('unnumbered display (\\[ or $$): every display must be a numbered, labelled environment: %s' % ctx(m))
    for m in re.finditer(r'\\begin\{(equation|align|gather|multline|eqnarray)\*\}', body):
        errors.append('starred (unnumbered) display %s*: %s' % (m.group(1), ctx(m)))
    for m in re.finditer(r'\\begin\{(equation|align|gather|multline)\}(.*?)\\end\{\1\}', body, re.S):
        if '\\label' not in m.group(2):
            warns.append('numbered display without \\label: %s' % ctx(m))
    if 'reqno' not in pre and 'amsart' in pre:
        warns.append('amsart without reqno: equation numbers must be on the right')
    if not re.search(r'\\numberwithin\{equation\}\{section\}', pre):
        warns.append('missing \\numberwithin{equation}{section}')
    # Links with border + hover previews + back-references.
    if re.search(r'hidelinks|colorlinks\s*=\s*true', pre):
        errors.append('hidelinks/colorlinks=true: use the v39 configuration (colorlinks=false, boxed links)')
    for need, why in [('pagebackref', 'back-references (Cited on p. ...)'), ('pdfcomment', 'hover tooltips'),
                      (r'\\CGref', '\\CGref/\\CGeqref macros'), ('pdfborder', 'boxed links'), (r'\\backrefalt', 'Cited on p./pp. format')]:
        if not re.search(need, pre):
            errors.append('preamble lacks %s (%s)' % (need.replace('\\\\', '\\'), why))
    for m in re.finditer(r'(?<![A-Za-z@])\\(eq)?ref\{', body):
        errors.append('raw \\%sref: use \\CG%sref (hover preview): %s' % (m.group(1) or '', m.group(1) or '', ctx(m)))
    tips = set(re.findall(r'\\csname CGTip@([^\\]+)\\endcsname', pre))
    for lb in re.findall(r'\\label\{([^}]+)\}', body):
        if lb not in tips:
            warns.append('label %s has no CGTip tooltip in the preamble' % lb)
    for m in re.finditer(r'\\(section|subsection)\*?\{[^}]*\}(?!\s*\\label)', body):
        warns.append('%s without \\label: %s' % (m.group(1), ctx(m)))
    if '\\tableofcontents' not in body:
        warns.append('no \\tableofcontents')
    # Author.
    au = re.search(r'\\author(?:\[[^\]]*\])?\{([^}]*)\}', text)
    rotkiv = au and re.search(r'Rotkiv Ayn(?:\\"a|\\"\{a\}|ä)sac', au.group(1))
    civil = au and re.search(r'Casa(?:\\~n|\\~\{n\}|ñ)a-Giner', au.group(1))
    if not (rotkiv or civil):
        errors.append('author must be Rotkiv Ayn\\"asac or V\\\'ictor Casa\\~na-Giner (found: %s)' % (au.group(1) if au else 'no \\author'))
    if rotkiv and 'rotkivanyasac@gmail.com' not in text:
        warns.append('author e-mail rotkivanyasac@gmail.com missing')
    # Bibliography: thebibliography, DOI or arXiv in every item.
    for m in re.finditer(r'\\bibitem(?:\[[^\]]*\])?\{([^}]*)\}(.*?)(?=\\bibitem|\\end\{thebibliography\})', body, re.S):
        if not re.search(r'doi|arxiv', m.group(2), re.I):
            warns.append('reference %s without DOI or arXiv link' % m.group(1))
    # Prose, not AI-style lists; no Questions/Open problems; no heuristics; no version history.
    for m in re.finditer(r'\\begin\{(itemize|enumerate|description)\}', body):
        errors.append('%s list in the paper body: write prose: %s' % (m.group(1), ctx(m)))
    for m in re.finditer(r'\\section\*?\{[^}]*(Question|Open problem|Open Problem|Problems)[^}]*\}', body):
        errors.append('section of questions/open problems: %s' % ctx(m))
    for m in re.finditer(r'\\begin\{heuristic\}|[Hh]euristic(?:ally)?\b', body):
        warns.append('heuristic material: either prove it or remove it: %s' % ctx(m))
    for m in re.finditer(r'previous version|earlier version|this version|in version|(?<![A-Za-z\\])v\d+\b|retract|withdrawn claim|\breferee', body, re.I):
        errors.append('version-history / referee mention inside the paper: %s' % ctx(m))
    for m in re.finditer(r'[Pp]rocedencia|[Pp]rovenance', body):
        errors.append('provenance paragraph: %s' % ctx(m))
    for m in re.finditer(r'Platt|3\s*\\cdot\s*10\^\{?12', body):
        warns.append('result depending on numerical zero verification (Platt / 3e12): %s' % ctx(m))
    # Appendices start on a new page.
    for m in re.finditer(r'\\appendix', body):
        if not re.search(r'\\clearpage\s*$', body[:m.start()]):
            errors.append('\\appendix not preceded by \\clearpage')
        tail = body[m.end():]
        for s in re.finditer(r'\\section\{', tail):
            if not re.search(r'\\clearpage\s*$', tail[:s.start()]):
                errors.append('appendix section without \\clearpage before it: %s' % tail[s.start():s.start() + 60].replace('\n', ' '))


STANDARD_OPS = {'Re', 'Im', 'mod', 'div', 'li', 'Li', 'd', 'e', 'i', 'tr', 'rank', 'supp', 'sgn', 'lcm', 'gcd', 'card', 'meas', 'dist', 'diam', 'ord', 'deg', 'id', 'T'}


def lambda_convention(text):
    """A global convention fixing the domain of the arithmetic functions (terms with argument <= 0 dropped) covers affine arguments."""
    return re.search(r'defined only for (?:positive )?(?:integers )?\$?[a-z]\s*(?:\\ge|\\geq)\s*1|every sum[^.]{0,60}restricted to arguments|terms? (?:with|whose) (?:an )?argument[^.]{0,60}(?:omitted|excluded|dropped)', text)


MATH_ENVS = r'(?:theorem|lemma|proposition|corollary|conjecture|definition|remark|hypothesis|example)'


def math_units(text, paper):
    """Units for the mathematical pass: TP entries, or statement+proof blocks of a manuscript."""
    if not paper:
        return [(u, t, b) for u, t, b in entries(text)]
    out = []
    for m in re.finditer(r'\\begin\{(' + MATH_ENVS + r')\}(.*?)\\end\{\1\}(\s*\\begin\{proof\}.*?\\end\{proof\})?', text, re.S):
        lab = re.search(r'\\label\{([^}]+)\}', m.group(2))
        out.append((lab.group(1) if lab else 'line %d' % (text[:m.start()].count('\n') + 1), m.group(1), m.group(0)))
    return out


def _math_spans(t):
    """Text of every inline or displayed formula of t."""
    t = re.sub(r'(?<!\\)%.*', '', t)
    out = re.findall(r'\$([^$]+)\$', t)
    out += re.findall(r'\\begin\{(?:equation|align|gather|multline)\*?\}(.*?)\\end\{(?:equation|align|gather|multline)\*?\}', t, re.S)
    return out


def check_orphan_letters(units, errors):
    # TP: the unit is the entry with its Setting chain (ERROR); papers: the result, its proof and the preceding text (WARN)
    """B9: a single letter that occurs in a formula of a proof and nowhere else in the unit (statement, Setting, proof)
    is undefined; usually it is a leftover of a renaming (e.g. o(log K) after K was renamed Y)."""
    for uid, setting, stmt, proof in units:
        if not proof:
            continue
        everything = setting + ' ' + stmt + ' ' + proof
        seen = set()
        for f in _math_spans(proof):
            # a symbol introduced by a reference in the same sentence ("By Lemma 8.4 ... over $(A_X,B_X]$") is defined there
            i = proof.find(f)
            sentence = re.split(r'\.\s', proof[max(0, i - 300):i])[-1] if i >= 0 else ''
            if re.search(r'\\(?:ref|CGref|CGeqref|eqref|cite)\{|\\ResolvedLink\{', sentence):
                continue
            g = re.sub(r'\\(?:label|tag|ref|eqref|CGref|CGeqref|cite|ResolvedLink|UniqueTarget)\s*(?:\[[^\]]*\])?\{[^{}]*\}(?:\{[^{}]*\})?', ' ', f)
            g = re.sub(r'\\(?:mathrm|text|operatorname|mathbb|mathbf|mathcal|mathscr|mathfrak|textnormal)\s*(?:\{[^{}]*\}|\w)', ' ', g)
            g = re.sub(r'(?<![A-Za-z\\])[Oo]\s*(?:_\s*(?:\{[^{}]*\}|\w))?\s*(?:\\[Bb]igg?[lr]?|\\left)?\s*[(\[]', ' (', g)
            for v in sorted(set(re.findall(r'(?<![\\A-Za-z])([A-Za-z])(?![A-Za-z])', g)) - set('deio')):
                n = len(re.findall(r'(?<![\\A-Za-z])%s(?![A-Za-z])' % v, ' '.join(_math_spans(everything))))
                if n == 1 and v not in seen:
                    seen.add(v)
                    errors.append('MATH %s: the letter %s occurs once, in a formula of the proof, and is defined nowhere in the unit (leftover of a renaming?): %s' % (uid, v, re.sub(r'\s+', ' ', f)[:90]))


QUALIFIED = r'through|relative to|via|assuming|given|(?<!un)conditional|modulo|on Theorem|on Conjecture|of the cited'


def check_unconditional_titles(text, paper, errors, warns):
    """B13: a result whose title says "unconditional" may not depend on a cited result (proved only in a reference,
    e.g. an unrefereed preprint) unless the title says so ("unconditional through Theorem 7.1")."""
    if not paper:
        ents = entries(text)
        cited = {u for u, t, b in ents if re.search(r'\\StateNote\{CITED RESULT', b)}
        items = [(u, t, set(re.findall(r'result:(U\d{5})', b)) & cited - {u}) for u, t, b in ents]
    else:
        body = text.split('\\begin{document}', 1)[-1]
        envs = list(re.finditer(r'\\begin\{(' + REVIEW_ENVS + r')\}(?:\[((?:[^\[\]]|\[[^\]]*\])*)\])?(.*?)\\end\{\1\}', body, re.S))
        lab = lambda m: (re.search(r'\\label\{([^}]+)\}', m.group(3)) or [None, None])[1]
        cited = {lab(m) for m in envs if lab(m) and not re.search(r'\\begin\{proof\}', body[m.end():m.end() + 300])
                 and (re.search(r'\\cite', m.group(2) or '') or re.search(r'unrefereed|preprint', m.group(3)))}
        items = []
        for m in envs:
            # the environment and its proof, up to the next environment
            pm = re.match(r'\s*\\begin\{proof\}.*?\\end\{proof\}', body[m.end():], re.S)
            seg = m.group(3) + (pm.group(0) if pm else '')
            refs = set(re.findall(r'\\CG(?:eq)?ref\{([^}]+)\}', seg)) & cited - {lab(m)}
            items.append((lab(m) or m.group(1), m.group(2) or '', refs))
    for u, title, refs in items:
        if not refs or not re.search(r'unconditional', title, re.I) or re.search(QUALIFIED, title, re.I):
            continue
        strong = re.search(r'(?:^|[:,(]\s*)unconditional\b|unconditional\)?\s*$', re.sub(r'\\\w+\{[^{}]*\}|[{}]', '', title), re.I)
        msg = 'MATH %s: the title says "unconditional" but the result uses the cited result(s) %s, proved only in a reference: say "through ..." in the title or move it to the conditional class' % (u, ', '.join(sorted(refs)))
        (errors if strong else warns).append(msg + ('' if strong else ' (check the wording)'))


def check_math(text, paper, warns, errors):
    """Patterns of the mathematical errors that Doktor had to correct (see RULES, section B).
    These are candidates: every MATH line is checked by hand against RULES section B."""
    defs = {}
    # a global convention fixing the domain of lambda/mu (and dropping terms with argument <= 0) covers affine arguments
    convention = lambda_convention(text)
    # a letter used as a parameter of a named hypothesis, e.g. LU_k(A,K), and also assigned a value elsewhere
    params = set(re.findall(r'\\mathrm\{[A-Z]+\}_\{?\w*\}?\(\s*([A-Z])\s*,', text))
    for v in sorted(params):
        for m in re.finditer(r'(?<![A-Za-z\\_{])%s\s*(?::=|=)\s*(e\^|\\exp|[0-9])' % v, text):
            warns.append('MATH symbol %s is a parameter of a named hypothesis and is also assigned a value: %s' % (v, text[max(0, m.start() - 60):m.end() + 30].replace('\n', ' ')))
            break
    for uid, kind, b in math_units(text, paper):
        flat = re.sub(r'\s+', ' ', b)
        c = lambda m: flat[max(0, m.start() - 60):m.end() + 40]
        # 3a. division by a parameter whose stated range contains 0
        for m in re.finditer(r'(?:/\(|/|\\frac\{[^{}]*\}\{)\s*([a-z]?)\s*([a-zA-Z])\s*(?:\)|\})', flat):
            v = m.group(2)
            if v in 'nkpqmdNXYTHxzw':
                continue
            near = flat[max(0, m.start() - 120):m.end() + 120]
            if re.search(r'\\mathbf\s*1_\{?%s\s*>\s*0|%s\s*>\s*0|%s\s*\\ne(?:q)?\s*0|%s\s*(?:\\ge|\\geq|>)\s*1' % (v, v, v, v), near):
                continue
            if re.search(r'(?<![A-Za-z\\_{])%s\s*(?:\\ge|\\geq|\\geqslant)\s*0|(?<![A-Za-z\\_{])%s\s*\\in\s*\\mathbb\s*\{?[RZ]' % (v, v), flat):
                warns.append('MATH %s: division by %s while %s=0 is allowed -- treat %s=0 separately: %s' % (uid, v, v, v, c(m)))
        # 3a. an arithmetic function at an affine argument with no positivity condition in the unit
        for m in re.finditer(r'\\(lambda|mu|Lambda|Omega|omega|tau|sigma|varphi)\(\s*([^()]*?[a-zA-Z]\s*[a-zA-Z]?\s*[+-]\s*[^()]*?)\)', flat):
            if not convention and not re.search(r'\\ge\s*1|\\geq\s*1|\\ge 1|>\s*0|\\geq 1|\\mathbb\s*\{?N|positive', flat):
                warns.append('MATH %s: \\%s(%s) at an affine argument, no condition making it >=1 in this unit' % (uid, m.group(1), m.group(2)[:30]))
                break
        # 3a. inequality divided by a factor whose sign depends on a parameter allowed to be negative
        negs = set(re.findall(r'(?<![A-Za-z\\_{])([a-zA-Z])\s*<\s*0|(?<![A-Za-z\\_{])([a-zA-Z])\s*\\in\s*\\mathbb\s*\{?[RZ]', flat))
        negs = {x or y for x, y in negs}
        if negs and re.search(r'[Dd]ivid|[Mm]ultipl\w* (?:both sides|the inequality|by)', flat):
            warns.append('MATH %s: an inequality is divided or multiplied while %s may be negative -- treat each sign separately' % (uid, ', '.join(sorted(negs))))
        # 3d. scope claims that need both directions / the quantitative form
        for m in re.finditer(r'equivalent to|if and only if|\\iff|\\Longleftrightarrow|necessary and sufficient|(?<!in)\bsufficient\b(?!ly)|it suffices|is necessary|is now (?:a theorem|unconditional|\\cite)|holds unconditionally|unconditionally', flat):
            warns.append('MATH %s: scope claim "%s" -- prove both directions, quantitative form, exact parameter range: %s' % (uid, m.group(0), c(m)))
        # 3e. unrefereed input used as unconditional
        if re.search(r'unconditional', flat) and re.search(r'\\cite', flat) and not re.search(r'preprint|unrefereed', flat):
            warns.append('MATH %s: "unconditional" next to a citation without saying whether it is a refereed result' % uid)
        # 3f. delegated steps
        for m in re.finditer(r'[Ss]imilarly|[Aa]nalogously|the same argument|the same expansion|the same proof|as in the proof of|as in the two-point case|mutatis mutandis|one checks|it is easy to see|easily seen|standard argument|by a standard|adaptation of|adapting|we omit|left to the reader|details are omitted', flat, re.I):   # also at the start of a sentence
            if re.match(r'[^.]{0,20}(?:written out|proved|carried out) in (?:Lemma|Proposition|Theorem|\\S|Section|\\CGref|~)', flat[m.end():m.end() + 60]):
                continue
            warns.append('MATH %s: delegated step "%s" -- write the step out: %s' % (uid, m.group(0), c(m)))
        # 3c. symbol given two different definitions
        for m in re.finditer(r'(?<![A-Za-z\\_^{])([A-Z])\s*(?::=|=)\s*(e\^\{[^}]*\}|\\[a-z]+\{?[^,;$]{0,15}|[0-9][^,;$]{0,10})', flat):
            defs.setdefault(m.group(1), {}).setdefault(re.sub(r'\s', '', m.group(2))[:20], uid)
    for v, d in defs.items():
        if len(d) > 1:
            warns.append('MATH symbol %s defined with %d different values (%s) -- one symbol, one meaning' % (v, len(d), '; '.join('%s in %s' % (k, u) for k, u in list(d.items())[:4])))
    # 3c. TP: \mathrm/\mathcal names in a statement that its Setting never defines
    if not paper:
        byid = {u: b for u, _, b in entries(text)}
        for uid, title, b in entries(text):
            sp = setting_spans(b)
            setting = setting_context(uid, byid)
            stmt = tp_statement(b, sp)
            for name in sorted(set(re.findall(r'\\mathrm\{([A-Z][A-Za-z]{1,4})\}', stmt))):
                if name not in setting and not re.search(r'(?:say|define|denote|let|write)[^.]{0,80}\\mathrm\{%s\}' % name, stmt, re.I):
                    warns.append('MATH %s: \\mathrm{%s} used in the statement but not defined in its Setting' % (uid, name))


# ---------------------------------------------------------------------------
# Deeper mathematical checks (first index values, reciprocals, domains,
# definitions before use) and bibliography (registry of verified sources,
# rendering of the bibliography pages of the PDF).

import random

def _to_py(expr):
    """Very small LaTeX -> Python arithmetic converter (no sympy needed)."""
    e = expr
    e = re.sub(r'\\left|\\right|\\,|\;|\\!|\\big|\\Big|\\bigl|\\bigr|\\Bigl|\\Bigr', '', e)
    e = re.sub(r'\\frac\s*\{([^{}]*)\}\s*\{([^{}]*)\}', r'((\1)/(\2))', e)
    e = re.sub(r'\\frac\s*(\w)\s*(\w)', r'((\1)/(\2))', e)
    e = re.sub(r'\\cdot|\\times', '*', e)
    e = re.sub(r'\\(log|exp|sqrt)\b', r'\1', e)
    e = e.replace('{', '(').replace('}', ')').replace('^', '**')
    e = re.sub(r'_\(?\w+\)?', '', e)                      # drop subscripts: a_1 -> a
    if re.search(r'\\[A-Za-z]', e):
        return None
    e = re.sub(r'(?<=[0-9A-Za-z)])\s*(?=[A-Za-z(])', '*', e)   # implicit products
    e = re.sub(r'\b(log|exp|sqrt)\*\(', r'\1(', e)
    return e


def _is_zero(expr, var, value):
    """True if expr vanishes identically (for random values of the other letters) when var=value."""
    py = _to_py(expr)
    if not py:
        return False
    import math
    vals = []
    for seed in (1, 2):
        rnd = random.Random(seed)
        env = {k: rnd.uniform(1.3, 3.7) for k in sorted(set(re.findall(r'[A-Za-z]', py))) if k not in ('log', 'exp', 'sqrt')}
        env.update(log=math.log, exp=math.exp, sqrt=math.sqrt)
        env[var] = value
        try:
            vals.append(eval(re.sub(r'\b%s\b' % var, '(%r)' % value, py), {'__builtins__': {}}, env))
        except ZeroDivisionError:
            return True
        except Exception:
            return False
    return all(abs(v) < 1e-12 for v in vals)


def _lower_limit(sub):
    """(variable, first value) of a summation/product subscript, if it can be read."""
    for pat, shift in [(r'^\s*([a-zA-Z])\s*=\s*(-?\d+)', 0), (r'(-?\d+)\s*\\le(?:q|qslant)?\s*([a-zA-Z])(?![A-Za-z])', 0),
                       (r'([a-zA-Z])\s*\\ge(?:q|qslant)?\s*(-?\d+)', 0), (r'(-?\d+)\s*<\s*([a-zA-Z])(?![A-Za-z])', 1),
                       (r'([a-zA-Z])\s*>\s*(-?\d+)', 1)]:
        m = re.search(pat, sub)
        if m:
            a, b = m.group(1), m.group(2)
            var, start = (b, int(a)) if a.lstrip('-').isdigit() else (a, int(b))
            return var, start + shift
    return None


def check_deep(text, paper, warns, errors):
    for uid, kind, b in math_units(text, paper):
        flat = re.sub(r'\s+', ' ', b)
        # First value of every summation/product index: denominators and arithmetic arguments.
        for m in re.finditer(r'\\(sum|prod)\s*_\s*\{((?:[^{}]|\{[^{}]*\})*)\}(\s*\^\s*(?:\{[^{}]*\}|\S))?', flat):
            sub = re.sub(r'\\substack\{|\\\\', ' ', m.group(2))
            ll = _lower_limit(sub)
            if not ll:
                continue
            var, start = ll
            summand = flat[m.end():m.end() + 160]
            summand = re.split(r'\\end\{|\\label|\\quad|\\qquad|\.\s|,\s+(?:where|and|so)|\\tag', summand)[0]
            dens = re.findall(r'\\frac\s*\{[^{}]*\}\s*\{([^{}]*)\}', summand) + re.findall(r'\\frac\s*\w\s*(\w)', summand) \
                + re.findall(r'/\s*\(([^()]*)\)', summand) + re.findall(r'/\s*([A-Za-z0-9_]+)', summand)
            for d in dens:
                if re.search(r'(?<![A-Za-z\\])%s(?![A-Za-z])' % var, d) and _is_zero(d, var, start):
                    errors.append('MATH %s: first value %s=%d makes the denominator %s zero in \\%s_{%s}' % (uid, var, start, d.strip(), m.group(1), m.group(2)[:40]))
            for f, arg in re.findall(r'\\(lambda|mu|Lambda|Omega|omega|tau|sigma|varphi|log)\s*\(([^()]*)\)', summand):
                if re.search(r'(?<![A-Za-z\\])%s(?![A-Za-z])' % var, arg):
                    py = _to_py(arg)
                    if py and not re.search(r'[A-Za-z]', re.sub(r'\b%s\b' % var, '', py)):
                        try:
                            val = eval(re.sub(r'\b%s\b' % var, '(%d)' % start, py), {'__builtins__': {}}, {})
                        except Exception:
                            continue
                        if (f != 'log' and val < 1) or (f == 'log' and val <= 0):
                            errors.append('MATH %s: first value %s=%d gives \\%s(%s) with argument %s' % (uid, var, start, f, arg.strip(), val))
                    elif f != 'log' and not lambda_convention(text):
                        warns.append('MATH %s: first value %s=%d in \\%s(%s): check the argument is >=1' % (uid, var, start, f, arg.strip()[:30]))
        # Numerator/denominator: reciprocal of a min/max, and a ratio defined one way and used the other.
        for m in re.finditer(r'\\(min|max)[^$]{0,60}?\\\}\s*\^\s*\{?-1|\\(min|max)[^$]{0,60}?\)\s*\^\s*\{?-1|(?:\\frac\s*\{?1\}?\s*\{|1/)\s*\\(min|max)', flat):
            warns.append('MATH %s: reciprocal of a min/max (1/min = max of reciprocals) -- check numerator and denominator: %s' % (uid, flat[max(0, m.start() - 40):m.end() + 40]))
    ratios = {}
    for m in re.finditer(r'(\\?[A-Za-z]+(?:_\{?\w+\}?)?)\s*:=\s*\\frac\s*\{([^{}]+)\}\s*\{([^{}]+)\}', text):
        ratios[m.group(1)] = (re.sub(r'\s', '', m.group(2)), re.sub(r'\s', '', m.group(3)))
    for name, (num, den) in ratios.items():
        for m in re.finditer(re.escape(name) + r'\s*=\s*\\frac\s*\{([^{}]+)\}\s*\{([^{}]+)\}', text):
            if (re.sub(r'\s', '', m.group(1)), re.sub(r'\s', '', m.group(2))) == (den, num):
                errors.append('MATH %s is defined as %s/%s but later written as %s/%s (numerator and denominator swapped)' % (name, num, den, den, num))

    # Domains declared and symbols defined before use.
    DOMAIN = r'(?:\\in|\\ge|\\geq|\\le|\\leq|>|<|\\ne|\\neq|:=|\\to|\\mid|\\subseteq|\\subset)'
    if not paper:
        byid = {u: b for u, _, b in entries(text)}
        for uid, title, b in entries(text):
            sp = setting_spans(b)
            ctxt = setting_context(uid, byid)
            stmt = tp_statement(b, sp)
            known = ctxt + stmt
            math = ' '.join(re.findall(r'\$([^$]+)\$', stmt))
            math = re.sub(r'\\(?:mathrm|text|operatorname|label|ref|tag|mathbb|mathbf|mathcal|mathscr|mathfrak)\s*(?:\{[^{}]*\}|\w)', ' ', math)
            math = re.sub(r'(?<![A-Za-z\\])[Oo]\s*(?:_\s*(?:\{[^{}]*\}|\w))?\s*\(', ' (', math)   # O(...), O_{p,q}(...)
            math = re.sub(r'U\^\s*(?:\{[^{}]*\}|\w)', ' ', math)                               # Gowers norms U^k
            bound = set(re.findall(r'\\(?:sum|prod|max|min|sup|inf|lim|int)\s*_\s*\{?\\?(?:substack\{)?\s*([a-zA-Z])', math))
            letters = set(re.findall(r'(?<![\\A-Za-z])([A-Za-z])(?![A-Za-z])', math)) - bound - set('edio')
            missing = []
            for v in sorted(letters):
                # a letter may be declared with a subscript (a_1,a_2\in\mathbb N, g_{\mathcal C}) and arguments (C_X(u):=)
                SUB = r'(?:\s*[_^]\s*(?:\{(?:[^{}]|\{[^{}]*\})*\}|\\?\w+|\*)){0,3}(?:\s*\([^()]*\))?'
                lst = r'(?<![A-Za-z\\])%s(?![A-Za-z])' % v + SUB + r'\s*(?:,\s*[A-Za-z]' + SUB + r'\s*)*'
                v_ = r'(?<![A-Za-z\\])%s(?![A-Za-z])' % v + SUB
                defined = (
                    # symbol set equal to something after write/put/let/set/define/denote/the value/there is, or at the start of a display line
                    # ($q=1$ after "when" or "if" is a condition, not a definition)
                    re.search(r'(?:\\q?quad|&|\\\\|\\begin\{[a-z*]+\}|(?:[Ww]rit(?:e|ing)|[Pp]ut|[Ll]et|[Ss]et|[Dd]efine|[Dd]enote|value|[Tt]here (?:is|are))\s+\$)\s*' + v_ + r'\s*=', known)
                    # inside a leading tuple set equal to a tuple: $(A_1,B_1,Q)=(1,0,1)$ (not $(p,q)=1$, a gcd condition)
                    or re.search(r'\$\s*\((?:[^()$]*[,(]\s*)?' + v_ + r'\s*[,)][^()$]*\)?\s*=\s*\(', known)
                    # let/define ... be, maps g:A\to B
                    or re.search(r'(?:[Ll]et|[Dd]efine|[Ww]rite|[Dd]enote)\s+\$\s*' + v_ + r'\s*(?::|\$\s*(?:be|denote|for))', known)
                    or re.search(v_ + r'\s*:\s*(?:\\mathbb|\\\{|[A-Z])', known)
                    # defined by reference: $G_y$ of Definition 9.3, the twists $b_i$ of [2]
                    or re.search(v_ + r'\s*\$\s*(?:of|as in|from|in)\s+(?:Definition|Lemma|Theorem|Proposition|Hypothesis|Subsection|Section|\\CG|\\ref|\\cite|\{?\\textnormal)', known)
                    or re.search(r'(?:sets?|intervals?|functions?|maps?|twists?|cutoffs?|coefficients?|sequences?|vectors?|famil(?:y|ies)|class(?:es)?|residues?|norms?)\s+\$\s*' + v_, known))
                if not defined and not re.search(lst + DOMAIN, known) and not re.search(r'(?:[A-Za-z]' + SUB + r'\s*,\s*)+%s' % v + SUB + r'\s*' + DOMAIN, known) \
                        and not re.search(r'(?:integers?|reals?|real numbers?|primes?|numbers?|parameters?|fixed|every|each|all|any|constants?)\s*\$?\s*(?:[A-Za-z](?:_\{?\w+\}?)?\s*,\s*)*%s(?![A-Za-z])' % v, known):
                    missing.append(v)
            if missing:
                warns.append('MATH %s: no domain/definition for %s in the statement or its Setting' % (uid, ', '.join(missing)))
    else:
        body = text.split('\\begin{document}', 1)[-1]
        names = set(re.findall(r'\\math(?:rm|cal|scr|frak)\s*\{([A-Za-z]{1,6})\}', body)) - STANDARD_OPS
        for nm in sorted(names):
            pat = r'\\math(?:rm|cal|scr|frak)\s*\{%s\}' % nm
            uses = [m.start() for m in re.finditer(pat, body)]
            dm = re.search(r'(?:[Dd]efine|[Dd]enote|[Ll]et|[Ww]rite|[Ss]ay that|[Cc]all)[^.]{0,120}' + pat + r'|' + pat + r'[^$.]{0,40}?(?::=|\\coloneqq)|\\begin\{(?:hypothesis|definition|conjecture)\}\[[^\]]*' + pat, body)
            if dm and uses and uses[0] < dm.start() - 200:
                warns.append('MATH \\mathrm{%s} is used (offset %d) before it is defined (offset %d): %s' % (nm, uses[0], dm.start(), body[uses[0]:uses[0] + 80].replace('\n', ' ')))
            if not dm and len(uses) > 1:
                warns.append('MATH \\mathrm{%s} is used %d times and never defined' % (nm, len(uses)))


def check_bib_registry(text, registry, warns):
    """References already verified in the original are listed in the registry (e.g. claude/<paper>-fuentes-verificadas.md);
    they are not re-checked.  The others must be checked once and added to the registry at once."""
    reg = open(registry, encoding='utf-8', errors='replace').read()
    keys = re.findall(r'\\bibitem(?:\[[^\]]*\])?\{([^}]*)\}', text)
    todo = [k for k in keys if k not in reg]
    print('BIB registry %s: %d of %d references already verified (do not re-check them)' % (registry, len(keys) - len(todo), len(keys)))
    for k in todo:
        warns.append('BIB %s not in the registry of verified sources: check it once in LIBRARY TOTAL and record it' % k)


def render_bib(pdf, outdir, errors, warns):
    """Render the bibliography pages of the compiled PDF as PNG (to be looked at) and check their text."""
    import subprocess, os
    os.makedirs(outdir, exist_ok=True)
    n = int(re.search(r'Pages:\s*(\d+)', subprocess.run(['pdfinfo', pdf], capture_output=True, text=True).stdout).group(1))
    full = subprocess.run(['pdftotext', '-layout', pdf, '-'], capture_output=True, text=True).stdout
    pages = full.split('\f')
    for m in re.finditer(r'\?\?|\[\?\]', full):
        errors.append('PDF: undefined reference or citation "%s": %s' % (m.group(0), full[max(0, m.start() - 60):m.end() + 20].replace('\n', ' ')))
    start = None
    for i in range(len(pages) - 1, -1, -1):
        if re.search(r'^\s*(References|Bibliography|REFERENCES|BIBLIOGRAPHY)\s*$', pages[i], re.M):
            start = i + 1
            break
    if not start:
        warns.append('PDF: no References/Bibliography page found')
        return
    subprocess.run(['pdftoppm', '-png', '-r', '110', '-f', str(start), '-l', str(n), pdf, os.path.join(outdir, 'bib')], check=True)
    bibtext = '\n'.join(pages[start - 1:])
    bibtext = bibtext.split(re.search(r'(References|Bibliography|REFERENCES|BIBLIOGRAPHY)', bibtext).group(0), 1)[-1]
    items = [x for x in re.split(r'\n\s*(?=\[\w[^\]]{0,30}\]\s)', bibtext) if x.strip()]
    for it in items:
        head = ' '.join(it.split())[:70]
        if not re.search(r'doi|arXiv|arxiv', it):
            warns.append('PDF bibliography: entry without DOI/arXiv as printed: %s' % head)
        if not re.search(r'Cited on', it):
            warns.append('PDF bibliography: entry without "(Cited on p. ...)": %s' % head)
    for f in sorted(os.listdir(outdir)):
        if f.startswith('bib') and f.endswith('.png'):
            print('BIB PNG to look at:', os.path.join(outdir, f))



def check_names_exponents(text, paper, warns, errors):
    """Variable collisions, renamings between parts of the text, and exponent errors."""
    body = text.split('\\begin{document}', 1)[-1] if paper else text
    # Renaming: variant letters for the same role used side by side.
    for a, b in [(r'\\epsilon\b', r'\\varepsilon\b'), (r'\\phi\b', r'\\varphi\b'), (r'\\rho\b', r'\\varrho\b'),
                 (r'\\theta\b', r'\\vartheta\b'), (r'\\kappa\b', r'\\varkappa\b')]:
        na, nb = len(re.findall(a, body)), len(re.findall(b, body))
        if na and nb:
            warns.append('MATH both %s (%d times) and %s (%d times) are used: same object renamed, or two objects too alike' % (a.strip('\\b').replace('\\\\', '\\'), na, b.strip('\\b').replace('\\\\', '\\'), nb))
    # Renaming: a named hypothesis or object written with different parameter letters in different places.
    sig = {}
    for m in re.finditer(r'\\mathrm\{([A-Z][A-Za-z]{0,5})\}\s*_\s*\{?([A-Za-z0-9\']{1,4})\}?\s*\(([^()]{1,30})\)', body):
        args = tuple(x.strip() for x in m.group(3).split(','))
        if all(re.fullmatch(r'\\?[A-Za-z]+', x) for x in args):
            sig.setdefault(m.group(1), {}).setdefault(args, body[:m.start()].count('\n') + 1)
    for nm, d in sig.items():
        if len(d) > 1:
            warns.append('MATH \\mathrm{%s} appears with different parameter letters: %s -- renamed in part of the text?' % (nm, '; '.join('(%s) line %d' % (', '.join(k), v) for k, v in d.items())))
    for uid, kind, b in math_units(text, paper):
        if not paper:      # in TP the Setting is context, not part of the statement or proof of the entry
            sp = setting_spans(b)
            b = ''.join(b[i] for i in range(len(b)) if not any(x <= i < y for x, y in sp)) if sp else b
        flat = re.sub(r'\s+', ' ', b)
        # Collision: a letter fixed as a parameter and also used as a summation/product index in the same unit.
        # "for every k" is a universal quantifier, not a fixed parameter: only Let / Fix / fixed fix a letter
        fixed = set(re.findall(r'(?:[Ll]et|[Ff]ix|fixed)\s*\$\s*([A-Za-z])\s*(?:\\ge|\\geq|\\in|>|,|\$)', flat))
        bound = set(re.findall(r'\\(?:sum|prod)\s*_\s*\{\s*(?:\\substack\{)?\s*(?:-?\d+\s*\\le(?:q)?\s*)?([A-Za-z])(?![A-Za-z])', flat))
        for v in sorted(fixed & bound):
            warns.append('MATH %s: %s is fixed as a parameter and also used as a summation index -- variable collision' % (uid, v))
        # Collision: one letter declared in two different sets in the same unit.
        doms = {}
        for m in re.finditer(r'(?<![A-Za-z\\_])([A-Za-z])\s*\\in\s*(\\mathbb\s*\{?[NZQRC]\}?)', flat):   # a subscript (q_i) is not the letter declared
            doms.setdefault(m.group(1), set()).add(re.sub(r'[\s{}]', '', m.group(2)))
        for v, d in doms.items():
            if len(d) > 1:
                warns.append('MATH %s: %s is declared in different sets (%s) in the same unit' % (uid, v, ', '.join(sorted(d))))
        # Exponents: braces, ambiguous 1/2+epsilon, statement vs end of proof.
        for m in re.finditer(r'\^\s*-\s*[0-9A-Za-z]', flat):
            errors.append('MATH %s: exponent without braces (renders only the sign): %s' % (uid, flat[max(0, m.start() - 20):m.end() + 10]))
        for m in re.finditer(r'\^\s*\{?\s*1/2\s*\}\s*\+\s*\\varepsilon|\^\s*\\frac12\s*\+\s*\\varepsilon|\^\s*\\tfrac12\s*\+\s*\\varepsilon', flat):
            errors.append('MATH %s: "+\\varepsilon" outside the exponent (x^{1/2}+eps instead of x^{1/2+eps}): %s' % (uid, flat[max(0, m.start() - 20):m.end() + 10]))
        for m in re.finditer(r'\}\s*\^\s*\{[^{}]*\}\s*\^\s*[{\w]', flat):
            warns.append('MATH %s: double superscript / nested exponent, check grouping: %s' % (uid, flat[max(0, m.start() - 20):m.end() + 10]))
        if re.search(PROOF_MARK, flat):
            stmt, proof = re.split(PROOF_MARK, flat, maxsplit=1)
            def exps(t):
                return set(re.sub(r'\s', '', x) for x in re.findall(r'(?:O\s*\(|\\ll|\\le|o\s*\()\s*[^$]{0,25}?(?<![A-Za-z])([XxNTBHY])\s*\^\s*\{([^{}]{1,25})\}', t) for x in [x[0] + '^' + x[1]])
            es, ep = exps(stmt), exps(proof[-1500:])
            for e in es:
                base = e.split('^')[0]
                ends = [x for x in ep if x.startswith(base + '^')]
                if ends and e not in ends:
                    warns.append('MATH %s: statement has %s but the end of the proof has %s -- check the exponent' % (uid, e, ', '.join(sorted(ends))))



# ---------------------------------------------------------------------------
# Unrequested changes, and review marks in LaTeX.

REVIEWERS = {
    # name: (macro, counter, colour name, rgb, label printed in the margin)
    'claude': ('VCGreviewed', 'VCGreviewcount', 'VCGgreen', '0,0.45,0.15', 'Claude'),
    'codex': ('VCGreviewedCodex', 'VCGreviewcountCodex', 'VCGblue', '0,0.25,0.7', 'Codex'),
    'chatgpt': ('VCGreviewedChatGPT', 'VCGreviewcountChatGPT', 'VCGorange', '0.75,0.35,0', 'ChatGPT'),
    'copilot': ('VCGreviewedCopilot', 'VCGreviewcountCopilot', 'VCGpurple', '0.45,0.1,0.55', 'Copilot'),
}


def review_preamble(reviewer='claude'):
    """Preamble block defining one reviewer's mark; each reviewer has its own macro, counter, colour and final total.
    The Claude block is exactly the one of checker v7, so files marked before stay recognised."""
    mac, cnt, col, rgb, who = REVIEWERS[reviewer]
    return (PREAMBLE_TEMPLATE.replace('@WHO@', who).replace('@MAC@', mac).replace('@CNT@', cnt)
            .replace('@COL@', col).replace('@RGB@', rgb))


PREAMBLE_TEMPLATE = r"""
% CorrectVCG review marks: every statement reviewed by @WHO@ carries \@MAC@{n},
% n = how many times that statement has been reviewed.  \VCGshowreviewsfalse hides the marks.
\makeatletter
\@ifundefined{@MAC@}{%
\newcounter{@CNT@}
\newif\ifVCGshowreviews \VCGshowreviewstrue
\definecolor{@COL@}{rgb}{@RGB@}
\newcommand{\@MAC@}[1]{\stepcounter{@CNT@}\ifVCGshowreviews\marginpar{\raggedright\tiny\textcolor{@COL@}{Reviewed by @WHO@\\ (review~#1)}}\fi}
\AtEndDocument{\ifVCGshowreviews\par\medskip\noindent{\small\textcolor{@COL@}{Statements reviewed by @WHO@: \arabic{@CNT@}.}}\fi}
}{}
\makeatother
"""


REVIEW_PREAMBLE = review_preamble('claude')
MARK_ANY = r'\\VCGreviewed(?:Codex|ChatGPT|Copilot)?\{\d+\}'


def mark_re(reviewer):
    return r'\\' + REVIEWERS[reviewer][0] + r'\{(\d+)\}'


def strip_review_marks(text):
    """Review marks and their preamble blocks are not text changes."""
    for r in REVIEWERS:
        text = text.replace(review_preamble(r) + '\n', '').replace(review_preamble(r), '')
    return re.sub(MARK_ANY, '', text)


def marks_by_unit(text, reviewer):
    """(unit, n) for every mark of one reviewer."""
    return sorted((unit_of(text, m.start()), int(m.group(1))) for m in re.finditer(mark_re(reviewer), text))

REVIEW_ENVS = r'(?:theorem|lemma|proposition|corollary|conjecture|definition|remark|hypothesis|example)'


def unit_of(text, pos):
    """Label of the statement (or TP entry) that contains position pos."""
    tp = [m for m in re.finditer(r'\\ResultHeading\{(U\d{5})\}', text[:pos])]
    if tp:
        return tp[-1].group(1)
    lab = [m for m in re.finditer(r'\\label\{([^}]+)\}', text[:pos])]
    return lab[-1].group(1) if lab else 'preamble'


def check_diff(old_path, text, allowed_path, errors, warns, reviewer='claude'):
    """Every change between the previous version and this one must be a requested correction.
    allowed_path lists, one per line, the labels / TP ids whose text may change."""
    import difflib
    old_text = strip_embedded(open(old_path, encoding='utf-8', errors='replace').read())
    for r in sorted(REVIEWERS):
        if r != reviewer and marks_by_unit(old_text, r) != marks_by_unit(text, r):
            errors.append('REVIEW marks of %s changed by %s: %s -> %s (only %s may touch them)' % (
                r, reviewer, marks_by_unit(old_text, r), marks_by_unit(text, r), r))
        if r != reviewer and (review_preamble(r) in old_text) != (review_preamble(r) in text):
            errors.append('REVIEW preamble block of %s added or removed by %s' % (r, reviewer))
    old = strip_review_marks(old_text).rstrip('\n').splitlines()
    text = strip_review_marks(text)
    new = text.rstrip('\n').splitlines()
    allowed = set()
    if allowed_path:
        allowed = {x.strip() for x in open(allowed_path, encoding='utf-8').read().split() if x.strip()}
    offsets, acc = [], 0
    for ln in new:
        offsets.append(acc)
        acc += len(ln) + 1
    n = 0
    for tag, i1, i2, j1, j2 in difflib.SequenceMatcher(None, old, new, autojunk=False).get_opcodes():
        if tag == 'equal':
            continue
        n += 1
        where = unit_of(text, offsets[min(j1, len(new) - 1)] if new else 0)
        snippet = (' '.join(old[i1:i2])[:90] + '  -->  ' + ' '.join(new[j1:j2])[:90]).replace('\t', ' ')
        if allowed and where not in allowed:
            errors.append('UNREQUESTED change in %s (lines %d-%d), not in the list of requested corrections: %s' % (where, j1 + 1, j2, snippet))
        elif not allowed:
            warns.append('CHANGE %d in %s (lines %d-%d): justify it with a requested correction or undo it: %s' % (n, where, j1 + 1, j2, snippet))
    print('DIFF against %s: %d changed blocks' % (old_path, n))


def mark_reviewed(path, out, which, reviewer='claude'):
    """Write a copy of path in which every reviewed statement carries the reviewer's mark (n incremented on each review).
    Only that reviewer's marks are added or incremented; every other reviewer's mark stays as it is."""
    mac = REVIEWERS[reviewer][0]
    full = open(path, encoding='utf-8').read()
    t = strip_embedded(full)
    tail = full[len(t):] if full.startswith(t) else ''
    if '\\newcommand{\\%s}' % mac not in t:
        t = t.replace('\\begin{document}', review_preamble(reviewer) + '\n\\begin{document}', 1)
    sel = None if which == 'all' else {x.strip() for x in open(which, encoding='utf-8').read().split() if x.strip()}
    count = [0]

    def bump(block, key):
        if sel is not None and key not in sel:
            return block
        count[0] += 1
        m = re.search(mark_re(reviewer), block)
        if m:
            return block[:m.start()] + '\\%s{%d}' % (mac, int(m.group(1)) + 1) + block[m.end():]
        return None

    def env(m):
        lab = re.search(r'\\label\{([^}]+)\}', m.group(0))
        key = lab.group(1) if lab else None
        b = bump(m.group(0), key)
        if b is not None:
            return b
        head = re.match(r'\\begin\{' + REVIEW_ENVS + r'\}(?:\[(?:[^\[\]]|\[[^\]]*\])*\])?(?:\s*\\label\{[^}]+\})?', m.group(0))
        rest = m.group(0)[head.end():]
        own = re.match(r'(?:' + MARK_ANY + r')*', rest).end()   # after the marks already there
        return m.group(0)[:head.end()] + rest[:own] + '\\%s{1}' % mac + rest[own:]
    if '\\ResultHeading{' in t:
        def entry(m):
            b = bump(m.group(0), m.group(1))
            if b is not None:
                return b
            i = m.start(2) - m.start(0) + len(m.group(2))
            own = re.match(r'(?:' + MARK_ANY + r')*', m.group(0)[i:]).end()
            return m.group(0)[:i + own] + '\\%s{1}' % mac + m.group(0)[i + own:]
        t = re.sub(r'\\ResultHeading\{(U\d{5})\}(\{' + BRACED + r'\})(?:(?!\\ResultHeading\{).)*', entry, t, flags=re.S)
    else:
        t = re.sub(r'\\begin\{(' + REVIEW_ENVS + r')\}.*?\\end\{\1\}', env, t, flags=re.S)
    head = EMBED_HEADER if full.startswith(EMBED_HEADER) else ''
    i = full.find('\n' + EMBED_BEGIN)
    open(out, 'w', encoding='utf-8').write(head + t + (full[i + 1:] if i >= 0 else ''))
    print('MARKED %d statements reviewed by %s in %s' % (count[0], REVIEWERS[reviewer][4], out))


def check_review_marks(text, paper, errors, reviewer='claude'):
    """After a review, every statement must carry this reviewer's mark."""
    own = '\\%s{' % REVIEWERS[reviewer][0]
    if '\\ResultHeading{' in text:
        for u, title, b in entries(text):
            if own not in b:
                errors.append('REVIEW mark of %s missing in %s (%s)' % (REVIEWERS[reviewer][4], u, title[:50]))
    else:
        for m in re.finditer(r'\\begin\{(' + REVIEW_ENVS + r')\}(.*?)\\end\{\1\}', text, re.S):
            if own not in m.group(2):
                lab = re.search(r'\\label\{([^}]+)\}', m.group(2))
                errors.append('REVIEW mark of %s missing in %s %s' % (REVIEWERS[reviewer][4], m.group(1), lab.group(1) if lab else 'line %d' % (text[:m.start()].count('\n') + 1)))


EMBED_BEGIN = '# BEGIN CORRECTVCG'
EMBED_END = '# END CORRECTVCG'
EMBED_HEADER = ('% CorrectVCG: after \\end{document} this file carries every correction rule and the checker (LaTeX ignores them).\n'
                '% Never remove that block unless Doktor asks. Before any change run:\n'
                "% sed -n '/^# BEGIN CORRECTVCG/,/^# END CORRECTVCG/p' THIS.tex > correctvcg_check.py && python3 correctvcg_check.py --rules && python3 correctvcg_check.py --paper THIS.tex\n")


def strip_embedded(text):
    """The checker never checks its own embedded copy (nor the header comment that announces it)."""
    i = text.find('\n' + EMBED_BEGIN)
    if i >= 0:
        text = text[:i + 1]
    return text.replace(EMBED_HEADER, '')


def embed(path, out):
    """Write a copy of path that carries this checker after \\end{document} and a header comment saying how to use it."""
    full = open(path, encoding='utf-8').read()
    i = full.find('\n' + EMBED_BEGIN)
    old = re.search(r'CorrectVCG checker \(version (\d+)\)', full[i:]) if i >= 0 else None
    mine = int(re.search(r'CorrectVCG checker \(version (\d+)\)', __doc__).group(1))
    if old and int(old.group(1)) > mine:
        raise SystemExit('%s already carries checker version %s, newer than this one (%d): extract and use that one' % (path, old.group(1), mine))
    t = strip_embedded(full)
    t = t.rstrip('\n') + '\n'
    if '\\end{document}' not in t:
        raise SystemExit('no \\end{document} in %s' % path)
    me = open(os.path.abspath(__file__), encoding='utf-8').read()
    # when run from a copy extracted with sed, drop its own BEGIN/END marker lines
    me = '\n'.join(l for l in me.splitlines() if not (l.startswith(EMBED_BEGIN) or l.startswith(EMBED_END))) + '\n'
    t = EMBED_HEADER + t + '\n' + EMBED_BEGIN + '  (CorrectVCG checker, ignored by LaTeX; extract with the sed line at the top)\n' + me.rstrip('\n') + '\n' + EMBED_END + '\n'
    open(out, 'w', encoding='utf-8').write(t)
    print('EMBEDDED checker (%d lines) in %s' % (me.count('\n') + 1, out))


def main():
    ap = argparse.ArgumentParser()
    ap.add_argument('file', nargs='?')
    ap.add_argument('--previous', help='previous version of the same chapter, to check that ids are stable')
    ap.add_argument('--verified', help='registry of verified sources (markdown); listed references are not re-checked')
    ap.add_argument('--bib-pdf', help='compiled PDF: render its bibliography pages as PNG and check them')
    ap.add_argument('--diff', help='previous version: every changed block must be a requested correction')
    ap.add_argument('--allowed', help='with --diff: labels / TP ids whose text may change (one per line)')
    ap.add_argument('--mark', help='write a copy with the review marks of --reviewer (Claude \\VCGreviewed{n}, Codex \\VCGreviewedCodex{n}, ChatGPT \\VCGreviewedChatGPT{n}, Copilot \\VCGreviewedCopilot{n}) to this path (then stop)')
    ap.add_argument('--mark-which', default='all', help='with --mark: "all" or a file of labels / TP ids reviewed')
    ap.add_argument('--require-review', action='store_true', help='ERROR for every statement without the review mark of --reviewer')
    ap.add_argument('--reviewer', default='claude', choices=sorted(REVIEWERS), help='who reviews: claude (default), codex, chatgpt or copilot; each has its own marks and never touches the others')
    ap.add_argument('--rules', action='store_true', help='print every correction Doktor has made (read before working)')
    ap.add_argument('--embed', help='write a copy of the .tex that carries this checker after \\end{document} (then stop)')
    ap.add_argument('--paper', action='store_true', help='check one of his manuscripts (format rules) instead of a TP chapter')
    ap.add_argument('--chapter', help='chapter number, e.g. 36 (inferred from tags if omitted)')
    a = ap.parse_args()
    if a.rules:
        print(RULES)
        return
    if a.embed:
        embed(a.file, a.embed)
        return
    if a.mark:
        mark_reviewed(a.file, a.mark, a.mark_which, a.reviewer)
        return
    marked = strip_embedded(open(a.file, encoding='utf-8', errors='replace').read())
    # review marks (of any reviewer) are invisible to every content check, so the result never depends on who reviewed
    text = re.sub(MARK_ANY, '', marked)
    errors, warns = [], []

    tags = re.findall(r'\\tag\{([^}]*)\}', text)
    ch = a.chapter or (Counter(t.split('.')[0] for t in tags).most_common(1)[0][0] if tags else None)
    targets = set(re.findall(r'\\UniqueTarget\{([^}]+)\}', text)) | set(re.findall(r'\\label\{([^}]+)\}', text)) \
        | set(re.findall(r'\\hypertarget\{([^}]+)\}', text)) | set(re.findall(r'\\NumberedFormula\{[^}]*\}\{([^}]+)\}', text))
    ids = [u for u, _, _ in entries(text)]
    idset = set(ids)

    # 1. Equation tags: no duplicates, no gaps inside this chapter.
    # Displays copied into Setting blocks repeat their tag; they must be identical to the original.
    settings = setting_spans(text)
    inside = lambda pos: any(a <= pos < b for a, b in settings)
    disp = {}
    for m in re.finditer(r'\\begin\{(equation\*?|align\*?|gather\*?|multline\*?)\}(.*?)\\end\{\1\}', text, re.S):
        for t in re.findall(r'\\tag\{([^}]*)\}', m.group(2)):
            body_eq = re.sub(r'\s+', '', m.group(2))
            disp.setdefault(t, []).append((inside(m.start()), body_eq))
    for t, lst in disp.items():
        originals = [b for ins, b in lst if not ins]
        if len(originals) > 1:
            errors.append('equation tag (%s) is displayed %d times outside Setting blocks' % (t, len(originals)))
        # A display that lives in source prose (not in an entry) appears in TP only inside Settings: allowed.
        ref = originals[0] if originals else lst[0][1]
        if any(b != ref for ins, b in lst if ins):
            errors.append('Setting copy of (%s) differs from the original display' % t)
    if ch and not a.paper:
        nums = sorted(int(t.split('.')[1]) for t in tags if re.fullmatch(r'%s\.\d+' % re.escape(ch), t))
        if nums:
            missing = sorted(set(range(nums[0], nums[-1] + 1)) - set(nums))
            if missing:
                errors.append('equation numbers missing in chapter %s: %s' % (ch, ', '.join('(%s.%d)' % (ch, k) for k in missing)))
        # every equation number quoted in the text must be displayed somewhere
        quoted = set(re.findall(r'\((%s\.\d+)\)' % re.escape(ch), re.sub(r'\\tag\{[^}]*\}', '', text)))
        for q in sorted(quoted - set(tags), key=lambda x: int(x.split('.')[1])):
            errors.append('(%s) is quoted but never displayed in the chapter' % q)

    # 2. Links must resolve.
    for lb in sorted(set(re.findall(r'\\ResolvedLink\{([^}]+)\}', text))):
        if lb not in targets:
            errors.append('ResolvedLink to missing target %s' % lb)
    for lb in sorted(set(re.findall(r'\\ref\{result:(U\d{5})\}', text))):
        if lb not in idset:
            warns.append('reference to entry %s outside this file (check it exists in the book)' % lb)
    for lb in sorted(set(re.findall(r'\\(?:CG)?(?:eq)?ref\{((?!result:)[^}]+)\}', text))):
        if lb not in targets:
            errors.append('raw \\ref/\\eqref/\\CGref to unknown label %s (generator did not resolve it)' % lb)

    # 3. Entry-level checks.
    for uid, title, body in entries(text):
        # an equation statement must not be replaced by the id of the theorem that contains it
        # The symptom: an entry id printed bare, without "Theorem~", "Lemma~", ... in front,
        # where the source referred to an equation inside that entry.
        for m in re.finditer(r'(\S{0,25}?)~?\\ref\{result:(U\d{5})\}', body):
            before = body[max(0, m.start() - 160):m.start() + len(m.group(1))]
            ok_single = re.search(KINDS + r's?(?:~|\s)*$', before)
            ok_list = re.search(KINDS + r's(?:~|\s)*(?:\\ref\{result:U\d{5}\}\s*(?:,|and|or|--)\s*)+(?:~|\s)*$', before)
            if not (ok_single or ok_list):
                errors.append('%s: bare entry id %s -- probably an equation reference resolved to the id of its theorem: %s'
                              % (uid, m.group(2), body[max(0, m.start() - 50):m.end() + 30].replace('\n', ' ')))
        sp = setting_spans(body)
        no_setting = ''.join(body[i] for i in range(len(body)) if not any(x <= i < y for x, y in sp)) if sp else body
        if re.search(r'\\ref\{result:%s\}' % uid, no_setting):
            errors.append('%s: entry refers to itself (circular reference)' % uid)
        for p in balance(re.sub(r'\\(?:begin|end)\{LitEntry\}', '', body)):
            errors.append('%s (%s): %s' % (uid, title[:50], p))
        sb = re.search(r'\\Setting(?:Block|AsIn)\{', body)
        if 'PROOF INCLUDED' in body and not sb and not title.lower().startswith(('remark', 'definition')):
            warns.append('%s (%s): no Setting block -- check that every symbol of the statement is defined' % (uid, title[:50]))

    # 4. Duplicate ids, id stability.
    for u, n in Counter(ids).items():
        if n > 1:
            errors.append('entry id %s used %d times' % (u, n))
    if a.previous:
        old = {t: u for u, t, _ in entries(strip_embedded(open(a.previous, encoding='utf-8', errors='replace').read()))}
        new = {t: u for u, t, _ in entries(text)}
        for t, u in new.items():
            if t in old and old[t] != u:
                errors.append('id changed for "%s": %s -> %s' % (t[:60], old[t], u))
        new_ids = set(new.values())
        for t in sorted(set(old) - set(new)):
            if old[t] not in new_ids:      # same id with a new heading (e.g. a number added) has not disappeared
                warns.append('entry disappeared since previous version: %s "%s"' % (old[t], t[:60]))

    # 5. Forbidden content.
    if a.paper:
        check_paper(text, errors, warns)
    else:
        # rule C1 holds for every .tex: a TP display is numbered (equation* with \tag) or it is an error
        for m in re.finditer(r'(?<!\\)\\\[|\$\$', text):
            where = unit_of(text, m.start())
            errors.append('unnumbered display (\\[ or $$) in %s: every display must be numbered: %s' % (where, text[m.start():m.start() + 70].replace('\n', ' ')))
    check_math(text, a.paper, warns, errors)
    check_unconditional_titles(text, a.paper, errors, warns)
    if not a.paper:
        byid = {u: b for u, _, b in entries(text)}
        units = []
        for u, _, b in entries(text):
            sp = setting_spans(b)
            parts = re.split(PROOF_MARK, b, maxsplit=1)
            units.append((u, setting_context(u, byid), tp_statement(b, sp), parts[1] if len(parts) > 1 else ''))
    else:
        body = text.split('\\begin{document}', 1)[-1]
        units = []
        for m in re.finditer(r'\\begin\{(' + REVIEW_ENVS + r')\}(.*?)\\end\{\1\}\s*(\\begin\{proof\}.*?\\end\{proof\})?', body, re.S):
            lab = re.search(r'\\label\{([^}]+)\}', m.group(2))
            units.append((lab.group(1) if lab else m.group(1), body[max(0, m.start() - 2500):m.start()], m.group(2), m.group(3) or ''))
    check_orphan_letters(units, errors if not a.paper else warns)
    check_deep(text, a.paper, warns, errors)
    check_names_exponents(text, a.paper, warns, errors)
    if a.diff:
        check_diff(a.diff, marked, a.allowed, errors, warns, a.reviewer)
    if a.require_review:
        check_review_marks(marked, a.paper, errors, a.reviewer)
    if a.verified:
        check_bib_registry(text, a.verified, warns)
    if a.bib_pdf:
        render_bib(a.bib_pdf, os.path.splitext(a.bib_pdf)[0] + '_bib_png', errors, warns)
    for pat in FORBIDDEN_WORDS:
        for m in re.finditer(pat, text, re.I):
            errors.append('forbidden word "%s" near: %s' % (m.group(0), text[max(0, m.start() - 40):m.end() + 40].replace('\n', ' ')))
    # His own work is never cited and never imported: his drafts may contain errors.  In a manuscript
    # his name is allowed only as author of the document itself, so only the bibliography and the citations are scanned.
    own_scope = text
    if a.paper:
        own_scope = '\n'.join(re.findall(r'\\begin\{thebibliography\}.*?\\end\{thebibliography\}', text, re.S)) \
            + '\n' + ' '.join(re.findall(r'\\cite[pt]?\*?(?:\[[^\]]*\])?\{([^}]*)\}', text))
    for pat in OWN_NAMES:
        for m in re.finditer(pat, own_scope):
            errors.append("his own work cited or imported (never cite Casana-Giner; his drafts may contain errors): %s"
                          % own_scope[max(0, m.start() - 60):m.end() + 40].replace('\n', ' '))
    body_txt = text.split('\\begin{document}', 1)[-1]
    for m in re.finditer(r'(?:our|my) (?:previous|earlier|forthcoming) (?:work|paper|note)|in preparation|forthcoming|\bdraft\b|compendium', body_txt, re.I):
        warns.append('material that may come from one of his drafts: re-prove it from scratch, a draft is no authority: %s'
                     % body_txt[max(0, m.start() - 60):m.end() + 40].replace('\n', ' '))
    for pat, why in LEFTOVERS:
        for m in re.finditer(pat, text):
            errors.append('%s: %s' % (why, text[max(0, m.start() - 40):m.end() + 40].replace('\n', ' ')))
    body_only = re.sub(r'\\SourceNote\{[^}]*\}', '', text)
    for pat in EDITION_MARKS:
        for m in re.finditer(pat, body_only):
            warns.append('possible edition/version mark: %s' % body_only[max(0, m.start() - 40):m.end() + 40].replace('\n', ' '))

    for e in errors:
        print('ERROR', e)
    for w in warns:
        print('WARN ', w)
    print('SUMMARY: %d entries, %d equation tags, %d errors, %d warnings' % (len(ids), len(tags), len(errors), len(warns)))
    sys.exit(1 if errors else 0)


if __name__ == '__main__':
    main()
# END CORRECTVCG
